\section{Invariants différentiels. Morphismes différentiellement lisses}
\label{section:IV.16}

\oldpage[IV-4]{5}
Dans ce paragraphe, nous présentons, sous forme globale, quelques notions de
calcul différentiel particulièrement utiles en Géométrie algébrique. Nous passons sous
silence de nombreux développements, classiques en Géométrie différentielle (connexions,
transformations infinitésimales associées à un champ de vecteurs, jets, etc.), bien que ces
notions s'écrivent de façon particulièrement naturelle dans le cadre des schémas. Nous
passons également sous silence ici les phénomènes spéciaux à la caractéristique $p>0$
(dont certains sont étudiés, dans le cadre affine, dans (0, 21)). Pour certains compléments
sur le formalisme différentiel dans les préschémas, le lecteur pourra consulter les
exposés II et VII de [42], ainsi que les chapitres ultérieurs de ce Traité.

\subsection{Invariants normaux d'une immersion}
\label{IV.16.1}

\begin{env}[16.1.1]
\label{IV.16.1.1}
Soient $(X,\mathcal{O}_X)$, $(Y,\mathcal{O}_Y)$ deux espaces annelés,
$f=(\psi,\theta):Y\to X$ un morphisme d'espaces annelés (0\textsubscript{I}, 4.1.1)
tel que l'homomorphisme
\[
  \theta^\#:\psi^*(\mathcal{O}_X)\longrightarrow\mathcal{O}_Y
\]
soit surjectif, de sorte que $\mathcal{O}_Y$ s'identifie à un faisceau d'anneaux quotient
$\psi^*(\mathcal{O}_X)/\mathscr{I}_f$. On peut alors munir
$\psi^*(\mathcal{O}_X)$ de la filtration $\mathscr{I}_f$-préadique.
\end{env}

\begin{definition}[16.1.2]
\label{IV.16.1.2}
Le faisceau d'anneaux $\mathcal{O}_Y$-augmenté
$\psi^*(\mathcal{O}_X)/\mathscr{I}_f^{n+1}$ est appelé le $n$-ème invariant normal
de $f$; l'espace annelé
$(Y,\psi^*(\mathcal{O}_X)/\mathscr{I}_f^{n+1})$ est appelé le $n$-ème voisinage
infinitésimal de $Y$ pour le morphisme $f$, et noté $Y_f^{(n)}$ ou simplement
$Y^{(n)}$. Le faisceau d'anneaux gradués associé au faisceau d'anneaux filtrés
$\psi^*(\mathcal{O}_X)$
\[
  \mathscr{G}r_\bullet(f)=\bigoplus_{n\geq 0}
  (\mathscr{I}_f^n/\mathscr{I}_f^{n+1})
  \tag{16.1.2.1}
  \label{IV.16.1.2.1}
\]
est appelé le faisceau d'anneaux gradués associé à $f$. Le faisceau
$\mathscr{G}r_1(f)=\mathscr{I}_f/\mathscr{I}_f^2$ est appelé le faisceau conormal
de $f$ (que l'on note aussi $\mathscr{N}_{Y/X}$ s'il n'en résulte pas de confusion).
\end{definition}

Il est clair que les
$\mathcal{O}_{Y^{(n)}}=\psi^*(\mathcal{O}_X)/\mathscr{I}_f^{n+1}$
(qu'on note aussi $\mathcal{O}_{Y_f^{(n)}}$) forment un
\oldpage[IV-4]{6}
système projectif de faisceaux d'anneaux sur $Y$, l'homomorphisme de transition
$\varphi_{nm}:\mathcal{O}_{Y^{(m)}}\to\mathcal{O}_{Y^{(n)}}$ pour $n\leq m$
identifiant $\mathcal{O}_{Y^{(n)}}$ au quotient de $\mathcal{O}_{Y^{(m)}}$ par la
puissance $(\mathscr{I}_f/\mathscr{I}_f^{n+1})^m$ de l'Idéal d'augmentation de
$\mathcal{O}_{Y^{(n)}}$, noyau de
$\varphi_{0n}:\mathcal{O}_{Y^{(n)}}\to\mathcal{O}_Y$. Les $Y^{(n)}$ forment par
suite un système inductif d'espaces annelés, ayant tous pour espace sous-jacent
l'espace $Y$, et on a des morphismes canoniques d'espaces annelés
$h_n:Y^{(n)}\to X$ égaux à $(\psi,\theta_n)$, où $\theta_n^\#$ est le morphisme
canonique
$\psi^*(\mathcal{O}_X)\to\psi^*(\mathcal{O}_X)/\mathscr{I}_f^{n+1}$.
Il est clair que le faisceau $\mathscr{G}r_\bullet(f)$ est un faisceau d'algèbres
graduées sur le faisceau d'anneaux
$\mathcal{O}_Y=\mathscr{G}r_0(f)$, et les $\mathscr{G}r_k(f)$ des
$\mathcal{O}_Y$-Modules.

Comme pour tout faisceau d'anneaux filtré, on a un homomorphisme canonique
surjectif de $\mathcal{O}_Y$-Algèbres graduées
\[
  \mathscr{S}_{\mathcal{O}_Y}^{\bullet}(\mathscr{G}r_1(f))
  \longrightarrow\mathscr{G}r_\bullet(f)
  \tag{16.1.2.2}
  \label{IV.16.1.2.2}
\]
coïncidant en degrés 0 et 1 avec les homomorphismes identiques.

\begin{examples}[16.1.3]
\label{IV.16.1.3}
\begin{enumerate}
\item[(i)] Supposons que $X$ soit un espace annelé en anneaux locaux, que $Y$ soit
réduit à un seul point $y$ (muni d'un anneau $\mathcal{O}_y$) et que, si
$x=\psi(y)$, $\theta^\#:\mathcal{O}_x\to\mathcal{O}_y$ soit un homomorphisme
surjectif d'anneaux ayant pour noyau l'idéal maximal $\mathfrak{m}_x$ de
$\mathcal{O}_x$. Alors les $\mathcal{O}_{Y^{(n)}}$ s'identifient aux anneaux
$\mathcal{O}_x/\mathfrak{m}_x^{n+1}$, et $\mathscr{G}r_\bullet(f)$ à l'anneau
gradué associé à l'anneau local $\mathcal{O}_x$ muni de sa filtration
$\mathfrak{m}_x$-préadique.
\item[(ii)] Supposons que $Y$ soit une partie fermée d'un sous-espace ouvert $U$
de $X$ et que $\mathcal{O}_Y$ soit induit sur $Y$ par un faisceau quotient
$\mathcal{O}_U/\mathscr{I}$, où $\mathscr{I}$ est un Idéal de $\mathcal{O}_U$
tel que $\mathscr{I}_x=\mathcal{O}_x$ pour tout $x\notin Y$; si $X$ est un
espace annelé en anneaux locaux, nous supposerons de plus que
$\mathscr{I}_x\neq\mathcal{O}_x$ pour $x\in Y$, de sorte que
$(Y,\mathcal{O}_Y)$ est encore un espace annelé en anneaux locaux.

Soit $\psi_0:Y\to U$ l'injection canonique, et notons
$\theta_0:\mathcal{O}_U\to(\psi_0)_*(\mathcal{O}_Y)$ l'homomorphisme tel que
$\theta_0^\#$ soit l'homomorphisme canonique
$\psi_0^*(\mathcal{O}_U)=\mathcal{O}_U|Y\to
(\mathcal{O}_U/\mathscr{I})|Y$, de sorte que
$j_0=(\psi_0,\theta_0):Y\to U$ est un morphisme d'espaces annelés (et d'espaces
annelés en anneaux locaux si $X$ est un espace annelé en anneaux locaux); si
$i:U\to X$ est l'injection canonique (morphisme d'espaces annelés),
$j=i\circ j_0$ est le morphisme $(\psi,\theta)$ de $Y$ dans $X$, où
$\psi:Y\to X$ est l'injection canonique, et
$\theta:\mathcal{O}_X\to\psi_*(\mathcal{O}_Y)$ est l'homomorphisme tel que
$\theta^\#=\theta_0^\#$. Comme $\theta^\#$ est surjectif, on peut appliquer les
définitions précédentes; $\mathcal{O}_{Y^{(n)}}$ est égal à
$\psi_0^*(\mathcal{O}_U/\mathscr{I}^{n+1})$, et l'on a
$(\psi_0)_*(\mathcal{O}_{Y^{(n)}})=\mathcal{O}_U/\mathscr{I}^{n+1}$, et
$\mathscr{G}r_n(j)=\mathscr{G}r_n(j_0)=
\psi_0^*(\mathscr{I}^n/\mathscr{I}^{n+1})=
j_0^*(\mathscr{I}^n/\mathscr{I}^{n+1})$.
\end{enumerate}
\end{examples}

\begin{env}[16.1.4]
\label{IV.16.1.4}
L'exemple (16.1.3, (ii)) montre qu'en général les
$\mathcal{O}_{Y^{(n)}}$ ne sont pas munis canoniquement d'une structure de
$\mathcal{O}_Y$-Module, ni a fortiori d'une structure de
$\mathcal{O}_Y$-Algèbre. La donnée d'une telle structure équivaut à celle d'un
homomorphisme de faisceaux d'anneaux
$\lambda_n:\mathcal{O}_Y\to\mathcal{O}_{Y^{(n)}}$, inverse à droite de
l'homomorphisme d'augmentation $\varphi_{0n}$; cela revient aussi à la donnée
d'un morphisme d'espaces annelés
$(1_Y,\lambda_n):Y^{(n)}\to Y$, inverse à gauche du morphisme canonique
$(1_Y,\varphi_{0n}):Y\to Y^{(n)}$.
\end{env}

\begin{proposition}[16.1.5]
\label{IV.16.1.5}
Soit $f=(\psi,\theta):Y\to X$ une immersion de préschémas. Alors:
\begin{enumerate}
\item[(i)] $\mathscr{G}r_\bullet(f)$ est une $\mathcal{O}_Y$-Algèbre graduée
quasi-cohérente.
\oldpage[IV-4]{7}
\item[(ii)] Les $Y^{(n)}$ sont des préschémas, canoniquement isomorphes à des
sous-préschémas de $X$.
\item[(iii)] Tout homomorphisme de faisceaux d'anneaux
$\lambda_n:\mathcal{O}_Y\to\mathcal{O}_{Y^{(n)}}$, inverse à droite de
l'homomorphisme d'augmentation $\varphi_{0n}$, fait de
$\mathcal{O}_{Y^{(n)}}$ et des $\mathcal{O}_{Y^{(k)}}$ pour $k\leq n$, des
$\mathcal{O}_Y$-Algèbres quasi-cohérentes; les structures de
$\mathcal{O}_Y$-Module déduites des structures précédentes sur les
$\mathscr{G}r_k(f)$ pour $k\leq n$ coïncident avec celles définies dans
(16.1.2).
\end{enumerate}
\end{proposition}

\noindent(i) La question étant locale sur $X$ et sur $Y$, on peut se borner au cas
où $Y$ est un sous-préschéma fermé de $X$ défini par un Idéal quasi-cohérent
$\mathscr{I}$ de $\mathcal{O}_X$; comme $\mathcal{O}_Y$ est la restriction à
$Y$ de $\mathcal{O}_X/\mathscr{I}$, l'assertion (i) est évidente, et $Y^{(n)}$
est le sous-préschéma fermé de $X$ défini par l'Idéal quasi-cohérent
$\mathscr{I}^{n+1}$ de $\mathcal{O}_X$. Enfin, pour prouver (iii), notons que la
donnée de $\lambda_n$ fait de l'Idéal
$\mathscr{I}/\mathscr{I}^{n+1}$ de l'augmentation $\varphi_{0n}$ et de ses
quotients $\mathscr{I}/\mathscr{I}^{k+1}$ ($1\leq k\leq n$) des
$\mathcal{O}_Y$-Modules, et il suffit de prouver par récurrence sur $k$ que les
$\mathscr{I}/\mathscr{I}^{k+1}$ sont des $\mathcal{O}_Y$-Modules
quasi-cohérents et que la structure de $\mathcal{O}_Y$-Module quotient qu'on en
déduit sur $\mathscr{I}^k/\mathscr{I}^{k+1}$ est la même que celle définie en
(16.1.2). La seconde assertion est immédiate,
$\mathscr{I}^k/\mathscr{I}^{k+1}$ étant annulé par
$\mathscr{I}/\mathscr{I}^{n+1}$; la première résulte, par récurrence sur $k$,
de ce qu'elle est triviale pour $k=1$ et de ce que
$\mathscr{I}/\mathscr{I}^{k+1}$ est une extension de
$\mathscr{I}/\mathscr{I}^k$ par
$\mathscr{I}^k/\mathscr{I}^{k+1}$ (III, 1.4.17).

\begin{corollary}[16.1.6]
\label{IV.16.1.6}
Sous les hypothèses générales de (16.1.5), si l'immersion $f$ est localement de
présentation finie, les $\mathscr{G}r_n(f)$ sont des $\mathcal{O}_Y$-Modules
quasi-cohérents de type fini.
\end{corollary}

En effet, avec les notations de la démonstration de (16.1.5), $\mathscr{I}$ est un
Idéal de type fini de $\mathcal{O}_X$ (1.4.7), donc les
$\mathscr{I}^n/\mathscr{I}^{n+1}$ sont des $\mathcal{O}_Y$-Modules de type fini,
d'où la conclusion.

\begin{corollary}[16.1.7]
\label{IV.16.1.7}
Sous les hypothèses générales de (16.1.5), soit $g:X\to Y$ un morphisme de
préschémas inverse à gauche de $f$. Alors, pour tout $n$, le morphisme composé
$(1,\lambda_n):Y^{(n)}\xrightarrow{h_n}X\xrightarrow{g}Y$ définit un
homomorphisme de faisceaux d'anneaux
$\lambda_n:\mathcal{O}_Y\to\mathcal{O}_{Y^{(n)}}$ inverse à droite de
l'augmentation $\varphi_{0n}$, faisant de $\mathcal{O}_{Y^{(n)}}$ une
$\mathcal{O}_Y$-Algèbre quasi-cohérente; pour ces homomorphismes, les
homomorphismes de transition
$\varphi_{nm}:\mathcal{O}_{Y^{(m)}}\to\mathcal{O}_{Y^{(n)}}$ ($n\leq m$) sont
des homomorphismes de $\mathcal{O}_Y$-Algèbres. En outre, si $g$ est localement
de type fini, les $\mathcal{O}_{Y^{(n)}}$ sont des $\mathcal{O}_Y$-Modules
quasi-cohérents de type fini.
\end{corollary}

La première assertion résulte aussitôt des définitions et de (16.1.5). D'autre
part, si $g$ est localement de type fini, alors $f$ est localement de
présentation finie (1.4.3, (v)); les $\mathscr{G}r_n(f)$ étant alors des
$\mathcal{O}_Y$-Modules quasi-cohérents de type fini par (16.1.6), il en est de
même des $\mathcal{O}_Y$-Modules
$\mathscr{I}/\mathscr{I}^{n+1}$, qui sont des extensions d'un nombre fini de
$\mathscr{G}r_k(f)$ (III, 1.4.17).

\begin{proposition}[16.1.8]
\label{IV.16.1.8}
Soient $X$ un préschéma localement noethérien, $j:Y\to X$ une immersion. Alors
les $Y^{(n)}$ sont des préschémas localement noethériens, les
$\mathscr{G}r_n(j)$ sont des $\mathcal{O}_Y$-Modules cohérents, et
$\mathscr{G}r_\bullet(j)$ est un faisceau d'anneaux cohérent sur l'espace $Y$.
\end{proposition}

Tout étant local sur $X$ et $Y$, on est ramené au cas où $X$ est affine et $j$
une immersion fermée, et alors toutes les assertions sont évidentes sauf la
dernière, qui résulte de ce que si $A$ est un anneau noethérien et
$\mathfrak{I}$ un idéal de $A$, $\operatorname{gr}_{\mathfrak{I}}^\bullet(A)$
est un anneau noethérien, compte tenu de l'exactitude du foncteur $\psi^*$ et de
(0\textsubscript{I}, 5.3.7).

\begin{proposition}[16.1.9]
\label{IV.16.1.9}
\oldpage[IV-4]{8}
Soient $X$ un préschéma, $j:Y\to X$ une immersion localement de présentation
finie, $y$ un point de $Y$. Les conditions suivantes sont équivalentes:
\begin{enumerate}
\item[(a)] Il existe un voisinage ouvert $U$ de $y$ dans $Y$ tel que $j|U$ soit
un homéomorphisme de $U$ sur un ouvert de $X$.
\item[(b)] Il existe un entier $n>0$ tel que l'homomorphisme canonique
\[
  (\varphi_{n,n-1})_y:
  \mathcal{O}_{Y^{(n)},y}\longrightarrow\mathcal{O}_{Y^{(n-1)},y}
\]
soit bijectif.
\item[(c)] Il existe un entier $n>0$ tel que
$(\mathscr{G}r_n(j))_y=0$.
\end{enumerate}
De plus, si l'entier $n$ vérifie b) ou c), il existe un voisinage $V$ de $y$
dans $Y$ tel que $\mathscr{G}r_m(j)|V=0$ pour $m\geq n$ et que
$\varphi_{nm}|V:\mathcal{O}_{Y^{(m)}}|V\to\mathcal{O}_{Y^{(n)}}|V$ soit
bijectif pour $m\geq n$.
\end{proposition}

La question étant locale sur $Y$, on peut se borner au cas où $j$ est une
immersion fermée, $Y$ étant défini par un Idéal quasi-cohérent de type fini
$\mathscr{I}$ de $\mathcal{O}_X$. L'équivalence de b) et c), pour un $n$ donné,
est alors immédiate; en outre, comme
$\mathscr{I}^n/\mathscr{I}^{n+1}$ est un $\mathcal{O}_X$-Module de type fini,
il existe un voisinage ouvert $U$ de $y$ dans $X$ tel que
$\mathscr{I}^n|U=\mathscr{I}^{n+1}|U$ (0\textsubscript{I}, 5.2.2), donc aussi
$\mathscr{I}^n|U=\mathscr{I}^m|U$ pour $m\geq n$, ce qui prouve les dernières
assertions. Pour prouver que a) entraîne b), on peut se borner au cas où
l'espace sous-jacent à $Y$ est égal à l'espace sous-jacent à $X$, et où
$\mathscr{I}$ est engendré par un nombre fini de ses sections au-dessus de $X$:
comme $\mathscr{I}$ est alors contenu dans le Nilradical $\mathscr{N}$ de
$\mathcal{O}_X$ (I, 5.1.2), il est nilpotent, ce qui prouve b). Enfin, pour
prouver que b) entraîne a), on peut aussi se borner au cas où
$\mathscr{I}^n=\mathscr{I}^{n+1}$; alors, pour tout $y\in Y$, comme
$\mathscr{I}_y\subset\mathfrak{m}_y$, idéal maximal de
$\mathcal{O}_{X,y}$, on a nécessairement $\mathscr{I}_y^n=0$ en vertu du
lemme de Nakayama, puisque $\mathscr{I}_y$ est un idéal de type fini.
L'ensemble des $x\in X$ tels que $\mathscr{I}_x^n=0$ est donc un ouvert $U$ de
$X$ contenant $Y$ (0\textsubscript{I}, 5.2.2); comme d'autre part
$\mathscr{I}_x\neq 0$ pour $x\notin Y$, on a nécessairement $U=Y$.

\begin{corollary}[16.1.10]
\label{IV.16.1.10}
Pour que la restriction de l'immersion $j$ à un voisinage de $y$ dans $Y$ soit
une immersion ouverte (autrement dit pour que $j$ soit un isomorphisme local au
point $y$), il faut et il suffit que
$(\mathscr{G}r_1(j))_y=(\mathscr{N}_{Y/X})_y=0$.
\end{corollary}

La condition est évidemment nécessaire, et le raisonnement précédent, appliqué
pour $n=1$, prouve qu'elle est suffisante.

\begin{remarks}[16.1.11]
\label{IV.16.1.11}
\begin{enumerate}
\item[(i)] Sous les conditions de la définition (16.1.1), la limite projective
du système projectif
$(\mathcal{O}_{Y^{(n)}},\varphi_{nm})$ de faisceaux d'anneaux sur $Y$ est
appelé l'invariant normal d'ordre infini de $f$, et parfois noté
$\mathcal{O}_{Y^{(\infty)}}$. Lorsque $X$ est un préschéma localement
noethérien, $j:Y\to X$ une immersion fermée, $Y$ étant donc un sous-préschéma
fermé de $X$ défini par un Idéal cohérent $\mathscr{I}$,
$\mathcal{O}_{Y^{(\infty)}}$ n'est autre que le complété formel de
$\mathcal{O}_X$ le long de $Y$ (I, 10.8.4), et
$Y^{(\infty)}=(Y,\mathcal{O}_{Y^{(\infty)}})$ le préschéma formel complété de
$X$ le long de $Y$ (I, 10.8.5). Dans tous les cas, on pourra dire que
$Y^{(\infty)}$ est le voisinage formel de $Y$ dans $X$ (pour le morphisme
$f$). Dans le cas particulier qu'on vient d'envisager, c'est donc le préschéma
formel limite inductive des voisinages infinitésimaux d'ordre $n$.
\item[(ii)] On notera que pour un morphisme de préschémas
$f=(\psi,\theta):Y\to X$, il peut se faire que l'homomorphisme
$\theta^\#:\psi^*(\mathcal{O}_X)\to\mathcal{O}_Y$ soit surjectif sans que $f$
soit une immersion
\oldpage[IV-4]{9}
locale, et sans que $f$ soit injectif. On en a un exemple en prenant pour $Y$
une somme de préschémas $Y_\lambda$ tous isomorphes à
$\operatorname{Spec}(\mathcal{O}_x)$, où $x\in X$, et pour $f$ le morphisme
égal au morphisme canonique dans chacun des $Y_\lambda$.
\end{enumerate}
\end{remarks}

\subsection{Propriétés fonctorielles des invariants normaux d'une immersion}
\label{IV.16.2}

\begin{env}[16.2.1]
\label{IV.16.2.1}
Soient $f=(\psi,\theta):Y\to X$,
$f'=(\psi',\theta'):Y'\to X'$ deux morphismes d'espaces annelés tels que les
homomorphismes $\theta^\#$ et $\theta'^\#$ soient surjectifs; considérons un
diagramme commutatif de morphismes d'espaces annelés
\[
  \xymatrix{
    Y \ar[r]^f & X \\
    Y' \ar[r]_{f'} \ar[u]^u & X' \ar[u]_v
  }
  \tag{16.2.1.1}
  \label{IV.16.2.1.1}
\]
Posons $u=(\rho,\lambda)$, $v=(\sigma,\mu)$. On a
$\rho^*(\psi^*(\mathcal{O}_X))=
\psi'^*(\sigma^*(\mathcal{O}_X))$ et on a par suite un diagramme commutatif
d'homomorphismes de faisceaux d'anneaux sur $Y'$
\[
  \xymatrix{
    \rho^*(\psi^*(\mathcal{O}_X))
      =\psi'^*(\sigma^*(\mathcal{O}_X))
      \ar[r]^-{\psi'^*(\mu^\#)}
      \ar[d]_{\rho^*(\theta^\#)}
      & \psi'^*(\mathcal{O}_{X'}) \ar[d]^{\theta'^\#} \\
    \rho^*(\mathcal{O}_Y) \ar[r]_-{\lambda^\#}
      & \mathcal{O}_{Y'}
  }
\]
d'où l'on conclut, si $\mathscr{I}$ et $\mathscr{I}'$ sont les noyaux de
$\theta^\#$ et $\theta'^\#$, que l'on a
$\psi'^*(\mu^\#)(\rho^*(\mathscr{I}))\subset\mathscr{I}'$, compte tenu de
l'exactitude du foncteur $\rho^*$. On en déduit aussitôt que pour tout entier
$n$,
$\psi'^*(\mu^\#)(\rho^*(\mathscr{I}^n))\subset\mathscr{I}'^n$, ce qui montre
que $\psi'^*(\mu^\#)$ définit, par passage aux quotients, un homomorphisme de
faisceaux d'anneaux
\[
  \nu_n:
  \rho^*(\psi^*(\mathcal{O}_X)/\mathscr{I}^{n+1})
  \longrightarrow
  \psi'^*(\mathcal{O}_{X'})/\mathscr{I}'^{n+1}
  \tag{16.2.1.2}
  \label{IV.16.2.1.2}
\]
et par suite un morphisme d'espaces annelés
$w_n=(\rho,\nu_n):Y'^{(n)}\to Y^{(n)}$ (qui, pour $n=0$, n'est autre que
$u$). Il résulte aussitôt de cette définition que les diagrammes
\[
  \xymatrix@R=1pc{
    Y^{(n)} \ar[r]^-{h_{mn}} & Y^{(m)} \ar[r]^-{h_m} & X \\
    & & & (n\leq m) \\
    Y'^{(n)} \ar[r]_-{h'_{mn}} \ar[uu]^-{w_n}
      & Y'^{(m)} \ar[r]_-{h'_m} \ar[uu]^-{w_m}
      & X' \ar[uu]_-v
  }
\]
(où les flèches horizontales sont les morphismes canoniques (16.1.2)) sont
commutatifs.

Par passage aux quotients à partir des homomorphismes (16.2.1.2), et en tenant
\oldpage[IV-4]{10}
compte de l'exactitude du foncteur $\rho^*$, on obtient un di-homomorphisme
d'Algèbres graduées (relatif à l'homomorphisme
$\lambda^\#:\rho^*(\mathcal{O}_Y)\to\mathcal{O}_{Y'}$)
\[
  \operatorname{gr}(u):
  \rho^*(\mathscr{G}r_\bullet(f))
  \longrightarrow\mathscr{G}r_\bullet(f')
  \tag{16.2.1.3}
  \label{IV.16.2.1.3}
\]
(ou, si on veut, un $\rho$-morphisme (0\textsubscript{I}, 3.5.1)
$\mathscr{G}r_\bullet(f)\to\mathscr{G}r_\bullet(f')$), et en particulier un
di-homomorphisme des faisceaux conormaux
\[
  \operatorname{gr}_1(u):
  \rho^*(\mathscr{G}r_1(f))\longrightarrow\mathscr{G}r_1(f').
\]
Il est immédiat en outre que ces homomorphismes donnent lieu à un diagramme
commutatif
\[
  \xymatrix{
    \rho^*(\mathscr{S}_{\mathcal{O}_Y}^{\bullet}
      (\mathscr{G}r_1(f)))
      \ar[r] \ar[d]_-{\mathscr{S}(\operatorname{gr}_1(u))}
      & \rho^*(\mathscr{G}r_\bullet(f))
        \ar[d]^-{\operatorname{gr}(u)} \\
    \mathscr{S}_{\mathcal{O}_{Y'}}^{\bullet}
      (\mathscr{G}r_1(f'))
      \ar[r]
      & \mathscr{G}r_\bullet(f')
  }
  \tag{16.2.1.4}
  \label{IV.16.2.1.4}
\]
où les flèches horizontales sont les homomorphismes canoniques (16.1.2.2).

Enfin, si l'on a un diagramme commutatif de morphismes d'espaces annelés
\[
  \xymatrix{
    Y \ar[r]^f & X \\
    Y' \ar[r]_{f'} \ar[u]^u & X' \ar[u]_v \\
    Y'' \ar[r]_{f''} \ar[u]^{u'} & X'' \ar[u]_{v'}
  }
\]
où $f''=(\psi'',\theta'')$ est tel que $\theta''^\#$ soit surjectif, et si
$w'_n$ et $w''_n$ sont définis à partir de $u'$, $v'$ d'une part, et de
$u''=u\circ u'$, $v''=v\circ v'$ de l'autre, alors on a
$w''_n=w_n\circ w'_n$, comme il résulte aussitôt des définitions et de
(0\textsubscript{I}, 3.5.5); on a de même
$\operatorname{gr}(u'')=
\operatorname{gr}(u')\circ\rho'^*(\operatorname{gr}(u))$ si
$u'=(\rho',\lambda')$. On peut donc dire que les $Y^{(n)}$ et
$\mathscr{G}r_\bullet(f)$ dépendent fonctoriellement de $f$.
\end{env}

\begin{proposition}[16.2.2]
\label{IV.16.2.2}
Avec les notations et hypothèses de (16.2.1), supposons de plus que $f$, $f'$,
$u$ et $v$ soient des morphismes de préschémas. Alors:
\begin{enumerate}
\item[(i)] Les morphismes $w_n:Y'^{(n)}\to Y^{(n)}$ sont des morphismes de
préschémas.
\item[(ii)] Si $Y'=Y\times_X X'$, $u$ et $f'$ étant les projections
canoniques, et si $f$ est une immersion ou si $v$ est plat, on a
$Y'^{(n)}=Y^{(n)}\times_X X'$.
\item[(iii)] Si $Y'=Y\times_X X'$ et si $v$ est plat (resp. si $f$ est une
immersion), l'homomorphisme
\[
  \operatorname{Gr}(u)=\operatorname{gr}(u)\otimes 1:
  \mathscr{G}r_\bullet(f)\otimes_{\mathcal{O}_Y}\mathcal{O}_{Y'}
  \longrightarrow\mathscr{G}r_\bullet(f')
\]
est bijectif (resp. surjectif).
\end{enumerate}
\end{proposition}

\noindent(i) Les hypothèses entraînent aussitôt que pour tout $y'\in Y'$,
$\rho_{y'}^*(\theta_{\psi'(y')}^\#)$ est un homomorphisme local (I, 1.6.2),
donc $w_n$ est un morphisme de préschémas (I, 2.2.1).

\oldpage[IV-4]{11}
\noindent(ii) et (iii) Si $f$ est une immersion, on peut se borner au cas où
$f$ est une immersion fermée, $Y$ étant donc défini par l'Idéal
quasi-cohérent $\mathscr{I}$ de $\mathcal{O}_X$, et $Y^{(n)}$ par l'Idéal
$\mathscr{I}^{n+1}$; les assertions résultent alors de (I, 4.4.5).

Supposons en second lieu que $v$ soit plat; on peut se borner au cas où
$X=\operatorname{Spec}(A)$, $Y=\operatorname{Spec}(B)$,
$X'=\operatorname{Spec}(A')$ sont affines, $A'$ étant un $A$-module plat;
alors $Y'=\operatorname{Spec}(B')$ avec $B'=B\otimes_A A'$; en outre, si
$\mathfrak{I}$ est le noyau de l'homomorphisme $A\to B$, le noyau
$\mathfrak{I}'$ de $A'\to B'$ s'identifie à $\mathfrak{I}\otimes_A A'$ par
platitude, et
\[
  \mathfrak{I}'^n/\mathfrak{I}'^{n+1}
  =(\mathfrak{I}^n/\mathfrak{I}^{n+1})\otimes_A A'.
\]
On en déduit aussitôt, compte tenu de (0\textsubscript{I}, 4.3.3), que le
$\mathcal{O}_{Y'}$-Module
$\mathscr{I}'^n/\mathscr{I}'^{n+1}$ est égal à
\begin{align*}
&\psi'^*\!\left(\sigma^*
  \left((\mathfrak{I}^n/\mathfrak{I}^{n+1})^{\sim}\right)\right)
  \otimes_{\sigma^*(\mathcal{O}_X)}\mathcal{O}_{X'} \\
={}&
\psi'^*\!\left(\sigma^*
  \left((\mathfrak{I}^n/\mathfrak{I}^{n+1})^{\sim}\right)\right)
  \otimes_{\psi'^*(\sigma^*(\mathcal{O}_X))}
  \psi'^*(\mathcal{O}_{X'}) \\
={}&
\rho^*(\mathscr{I}^n/\mathscr{I}^{n+1})
  \otimes_{\rho^*(\psi^*(\mathcal{O}_X))}
  \psi'^*(\mathcal{O}_{X'}),
\end{align*}
et comme en particulier pour $n=0$, on a
\[
  \mathcal{O}_{Y'}=
  \rho^*(\mathcal{O}_Y)
  \otimes_{\rho^*(\psi^*(\mathcal{O}_X))}
  \psi'^*(\mathcal{O}_{X'}),
\]
on obtient un isomorphisme canonique de
$\mathscr{I}'^n/\mathscr{I}'^{n+1}$ sur
\[
  \rho^*(\mathscr{I}^n/\mathscr{I}^{n+1})
  \otimes_{\rho^*(\mathcal{O}_Y)}\mathcal{O}_{Y'}
  =
  (\mathscr{I}^n/\mathscr{I}^{n+1})
  \otimes_{\mathcal{O}_Y}\mathcal{O}_{Y'},
\]
ce qui prouve (iii). Posons maintenant
$C_n=\Gamma(Y,\mathcal{O}_{Y^{(n)}})$,
$C'_n=\Gamma(Y',\mathcal{O}_{Y'^{(n)}})$. Puisque $Y^{(n)}$ et
$Y'^{(n)}$ sont des schémas affines (16.1.5), le noyau
$\mathfrak{K}_n$ (resp. $\mathfrak{K}'_n$) de l'homomorphisme
$C_n\to C_{n-1}$ (resp. $C'_n\to C'_{n-1}$) est
$\Gamma(Y,\mathscr{I}^n/\mathscr{I}^{n+1})$ (resp.
$\Gamma(Y',\mathscr{I}'^n/\mathscr{I}'^{n+1})$); on déduit donc de ce qui
précède que $\mathfrak{K}'_n=\mathfrak{K}_n\otimes_A A'$. Or, on a un
diagramme commutatif
\[
  \xymatrix{
    0 \ar[r] &
    \mathfrak{K}_n\otimes_A A' \ar[r]\ar[d]^r &
    C_n\otimes_A A' \ar[r]\ar[d]^{s_n} &
    C_{n-1}\otimes_A A' \ar[r]\ar[d]^{s_{n-1}} & 0 \\
    0 \ar[r] &
    \mathfrak{K}'_n \ar[r] &
    C'_n \ar[r] &
    C'_{n-1} \ar[r] & 0
  }
\]
où la flèche verticale de gauche est bijective et les deux lignes sont exactes
($A'$ étant un $A$-module plat). On en déduit par récurrence que $s_n$ est
bijectif pour tout $n$, car c'est vrai par hypothèse pour $n=0$, et se déduit
par application du lemme des cinq pour $n$ quelconque. Cela prouve la seconde
assertion de (ii).

\begin{corollary}[16.2.3]
\label{IV.16.2.3}
Soient $g:X\to Y$, $u:Y'\to Y$ deux morphismes de préschémas,
$X'=X\times_Y Y'$, $g':X'\to Y'$ et $v:X'\to X$ les projections canoniques.
Soient $f:Y\to X$ une $Y$-section de $X$ (donc une immersion),
$f'=f_{(Y')}:Y'\to X'$ la $Y'$-section de $X'$ déduite de $f$ par le
changement de base $u$. Alors:
\begin{enumerate}
\item[(i)] Le morphisme
$w_n:Y_{f'}'^{(n)}\to Y_f^{(n)}$ correspondant à $f$, $f'$, $u$, $v$
(16.2.1) et le morphisme canonique
$h'_n:Y_{f'}'^{(n)}\to X'$ identifient $Y_{f'}'^{(n)}$ au produit
$Y_f^{(n)}\times_X X'$.
\item[(ii)] Si l'on munit $\mathcal{O}_{Y_f^{(n)}}$ (resp.
$\mathcal{O}_{Y_{f'}'^{(n)}}$) de la structure de $\mathcal{O}_Y$-Algèbre
définie par $g$ (resp. de la structure de $\mathcal{O}_{Y'}$-Algèbre définie
par $g'$) (16.1.6), l'homomorphisme de $\mathcal{O}_{Y'}$-Algèbres
\[
  \rho^*(\mathcal{O}_{Y_f^{(n)}})
  \otimes_{\mathcal{O}_Y}\mathcal{O}_{Y'}
  \longrightarrow
  \mathcal{O}_{Y_{f'}'^{(n)}}
  \tag{16.2.3.1}
  \label{IV.16.2.3.1}
\]
\oldpage[IV-4]{12}
déduit de l'homomorphisme $\nu_n$ (16.2.1.2) est bijectif. En outre,
l'homomorphisme de $\mathcal{O}_{Y'}$-Modules
\[
  \operatorname{Gr}_1(u):
  \mathscr{G}r_1(f)\otimes_{\mathcal{O}_Y}\mathcal{O}_{Y'}
  \longrightarrow\mathscr{G}r_1(f')
  \tag{16.2.3.2}
  \label{IV.16.2.3.2}
\]
est bijectif.
\end{enumerate}
\end{corollary}

\noindent(i) Notons d'abord que les morphismes $f':Y'\to X'$ et $u:Y'\to Y$
identifient $Y'$ au produit $Y\times_X X'$ (pour les morphismes structuraux
$f:Y\to X$ et $v:X'\to X$) (14.5.12.1). La conclusion de (i) résulte alors de
(16.2.2, (ii)), le morphisme $g$ étant une immersion.

\noindent(ii) Le diagramme commutatif
\[
  \xymatrix{
    Y_f^{(n)} \ar[d]^{h_n}
      & Y_{f'}'^{(n)} \ar[d]^-{h'_n} \ar[l]^{w_n} \\
    X \ar[d]^g
      & X' \ar[d]^{g'} \ar[l]^v \\
    Y
      & Y' \ar[l]^u
  }
\]
identifie $Y_{f'}'^{(n)}$ au produit $Y_f^{(n)}\times_X X'$ et $X'$ au produit
$X\times_Y Y'$, donc (I, 3.3.9) il identifie (pour les morphismes
$g'\circ h'_n$ et $w_n$) $Y_{f'}'^{(n)}$ au produit
$Y_f^{(n)}\times_Y Y'$. Comme $Y_f^{(n)}$ (resp. $Y_{f'}'^{(n)}$) est le
préschéma affine au-dessus de $Y$ (resp. $Y'$) associé à la
$\mathcal{O}_Y$-Algèbre $\mathcal{O}_{Y_f^{(n)}}$ (resp. à la
$\mathcal{O}_{Y'}$-Algèbre $\mathcal{O}_{Y_{f'}'^{(n)}}$), le fait que
l'homomorphisme canonique (16.2.3.1) soit bijectif résulte de (II, 1.5.2).
Enfin, l'homomorphisme canonique (16.2.3.1) est compatible avec les
augmentations
$\mathcal{O}_{Y_f^{(n)}}\to\mathcal{O}_Y$ et
$\mathcal{O}_{Y_{f'}'^{(n)}}\to\mathcal{O}_{Y'}$; comme
$\mathcal{O}_{Y_f^{(n)}}$ est somme directe (en tant que
$\mathcal{O}_Y$-Module) de $\mathcal{O}_Y$ et de l'Idéal d'augmentation
$\mathscr{I}/\mathscr{I}^{n+1}$, on voit donc que l'homomorphisme canonique
(16.2.3.1), restreint à
$(\mathscr{I}/\mathscr{I}^{n+1})\otimes_{\mathcal{O}_Y}\mathcal{O}_{Y'}$,
est une bijection de ce dernier sur
$\mathscr{I}'/\mathscr{I}'^{n+1}$. Pour $n=1$, cela montre que
$\operatorname{Gr}_1(u)$ est bijectif.

On notera que, sous les hypothèses de (16.2.3), les homomorphismes
$\operatorname{Gr}_n(u)$ sont surjectifs en vertu de ce qui précède, mais ne
sont pas bijectifs en général pour $n\geq 2$. Toutefois:

\begin{corollary}[16.2.4]
\label{IV.16.2.4}
Sous les hypothèses de (16.2.3), supposons que $u:Y'\to Y$ soit un morphisme
plat (resp. que les $\mathscr{G}r_n(f)$ soient des $\mathcal{O}_Y$-Modules
plats pour $n\leq m$). Alors l'homomorphisme
\[
  \operatorname{Gr}_n(u):
  \mathscr{G}r_n(f)\otimes_{\mathcal{O}_Y}\mathcal{O}_{Y'}
  \longrightarrow\mathscr{G}r_n(f')
\]
est bijectif pour tout $n$ (resp. pour $n\leq m$).
\end{corollary}

Si $u$ est plat, il en est de même de $v:X'\to X$ qui s'en déduit par
changement de base, et on sait déjà dans ce cas que $\operatorname{Gr}(u)$ est
bijectif (16.2.2, (iii)). Si les $\mathscr{G}r_n(f)$ sont plats pour
$n\leq m$, on voit d'abord par récurrence sur $n$ qu'il en est de même des
$\mathscr{I}/\mathscr{I}^{n+1}$ pour $n\leq m$, en vertu des suites exactes
\[
  0\longrightarrow\mathscr{I}^n/\mathscr{I}^{n+1}
  \longrightarrow\mathscr{I}/\mathscr{I}^{n+1}
  \longrightarrow\mathscr{I}/\mathscr{I}^n
  \longrightarrow 0
\]
\oldpage[IV-4]{13}
(0\textsubscript{I}, 6.1.2); en outre, on a alors des diagrammes commutatifs
\[
  \xymatrix{
    0 \ar[r] &
    (\mathscr{I}^n/\mathscr{I}^{n+1})
      \otimes_{\mathcal{O}_Y}\mathcal{O}_{Y'}
      \ar[r]\ar[d] &
    (\mathscr{I}/\mathscr{I}^{n+1})
      \otimes_{\mathcal{O}_Y}\mathcal{O}_{Y'}
      \ar[r]\ar[d] &
    (\mathscr{I}/\mathscr{I}^n)
      \otimes_{\mathcal{O}_Y}\mathcal{O}_{Y'}
      \ar[r]\ar[d] & 0 \\
    0 \ar[r] &
    \mathscr{I}'^n/\mathscr{I}'^{n+1}
      \ar[r] &
    \mathscr{I}'/\mathscr{I}'^{n+1}
      \ar[r] &
    \mathscr{I}'/\mathscr{I}'^n
      \ar[r] & 0
  }
\]
dans lesquels les lignes sont exactes (la première par platitude
(0\textsubscript{I}, 6.1.2)) et les deux dernières flèches verticales sont
bijectives en vertu de (16.2.3, (ii)); d'où la conclusion.

\begin{remarks}[16.2.5]
\label{IV.16.2.5}
\begin{enumerate}
\item[(i)] Le raisonnement de (16.2.2, (i)) s'applique encore lorsque dans
(16.2.1.1) il s'agit de morphismes d'espaces annelés en anneaux locaux
($\mathbf{Err}_{\mathrm{II}}$, (1.8.2)).
\item[(ii)] Dans (16.2.2, (ii)), la conclusion n'est plus nécessairement valable
lorsqu'on suppose seulement que $v$ et $f$ sont des morphismes de préschémas
($f$ vérifiant la condition de (16.1.1)). Par exemple (avec les notations de la
démonstration de (16.2.2, (ii)), il peut se faire que $\mathfrak{I}=0$ mais que
le noyau $\mathfrak{I}'$ de $A'\to B'=B\otimes_A A'$ ne soit pas nul et que
$B'\neq 0$, auquel cas on a $Y^{(n)}=Y$ pour tout $n$, mais
$Y'^{(n)}\neq Y'$. On a un exemple de ce fait en prenant
$A=\mathbb{Z}$, $B=\mathbb{Q}$,
$A'=\prod_{h=1}^{\infty}(\mathbb{Z}/m^h\mathbb{Z})$ où $m>1$.
\end{enumerate}
\end{remarks}

\begin{env}[16.2.6]
\label{IV.16.2.6}
Considérons le cas particulier du diagramme (16.2.1.1) où $X'=X$, $v$ étant
l'identité, $X$ est un préschéma, $Y$ un sous-préschéma de $X$, $Y'$ un
sous-préschéma de $Y$, $f$, $u$ et $f'=f\circ u$ les injections canoniques; le
di-homomorphisme (16.2.1.3) donne, par tensorisation avec
$\mathcal{O}_{Y'}$ au-dessus de $\rho^*(\mathcal{O}_Y)$, un homomorphisme de
$\mathcal{O}_{Y'}$-Algèbres graduées
\[
  u^*(\mathscr{G}r_\bullet(f))
  \longrightarrow\mathscr{G}r_\bullet(f')
  \tag{16.2.6.1}
  \label{IV.16.2.6.1}
\]
D'autre part, on identifie $\mathcal{O}_Y$ à
$\psi^*(\mathcal{O}_X)/\mathscr{I}_f$, $\mathcal{O}_{Y'}$ à
$\rho^*(\mathcal{O}_Y)/\mathscr{I}_u$; comme $\rho^*$ est un foncteur exact,
on a
$\rho^*(\mathcal{O}_Y)=
\rho^*(\psi^*(\mathcal{O}_X))/\rho^*(\mathscr{I}_f)=
\psi'^*(\mathcal{O}_X)/\rho^*(\mathscr{I}_f)$, et comme
$\mathcal{O}_{Y'}$ s'identifie par ailleurs à
$\psi'^*(\mathcal{O}_X)/\mathscr{I}_{f'}$, on voit que l'on a
$\mathscr{I}_u=\mathscr{I}_{f'}/\rho^*(\mathscr{I}_f)$. On en déduit pour
tout entier $n$ un homomorphisme canonique
$\mathscr{I}_{f'}^n/\mathscr{I}_{f'}^{n+1}\to
\mathscr{I}_u^n/\mathscr{I}_u^{n+1}$, d'où un homomorphisme canonique de
$\mathcal{O}_{Y'}$-Algèbres graduées
\[
  \mathscr{G}r_\bullet(f')
  \longrightarrow\mathscr{G}r_\bullet(u).
  \tag{16.2.6.2}
  \label{IV.16.2.6.2}
\]
\end{env}

\begin{proposition}[16.2.7]
\label{IV.16.2.7}
Soient $X$ un préschéma, $Y$ un sous-préschéma de $X$, $Y'$ un
sous-préschéma de $Y$, $j:Y'\to Y$ l'injection canonique. On a alors une suite
exacte de faisceaux conormaux ($\mathcal{O}_{Y'}$-Modules)
\[
  \xymatrix{
    j^*(\mathscr{N}_{Y/X}) \ar[r] &
    \mathscr{N}_{Y'/X} \ar[r] &
    \mathscr{N}_{Y'/Y} \ar[r] & 0
  }
  \tag{16.2.7.1}
  \label{IV.16.2.7.1}
\]
où les flèches sont les composantes de degré 1 des homomorphismes canoniques
(16.2.6.1) et (16.2.6.2).
\end{proposition}

La question étant locale, on peut se borner au cas où
$X=\operatorname{Spec}(A)$, $Y=\operatorname{Spec}(A/\mathfrak{I})$,
$Y'=\operatorname{Spec}(A/\mathfrak{K})$, $\mathfrak{I}$ et $\mathfrak{K}$
étant deux idéaux de $A$ tels que $\mathfrak{I}\subset\mathfrak{K}$; tout
revient à voir alors
\oldpage[IV-4]{14}
que la suite d'homomorphismes canoniques
$\mathfrak{I}/\mathfrak{K}\mathfrak{I}\to
\mathfrak{K}/\mathfrak{K}^2\to
(\mathfrak{K}/\mathfrak{I})/(\mathfrak{K}/\mathfrak{I})^2\to 0$
est exacte, ce qui est immédiat puisque l'image de
$\mathfrak{I}/\mathfrak{K}\mathfrak{I}$ dans
$\mathfrak{K}/\mathfrak{K}^2$ est
$(\mathfrak{I}+\mathfrak{K}^2)/\mathfrak{K}^2$ et que
$(\mathfrak{K}/\mathfrak{I})/(\mathfrak{K}/\mathfrak{I})^2$ s'identifie à
$\mathfrak{K}/(\mathfrak{I}+\mathfrak{K}^2)$.

Il est facile de donner des exemples où la suite (16.2.7.1) prolongée à gauche
par un 0 n'est plus exacte; avec les notations précédentes, il suffit de
prendre $A=k[T]$, $\mathfrak{I}=AT^2$, $\mathfrak{K}=AT$, car on a alors
$(\mathfrak{I}+\mathfrak{K}^2)/\mathfrak{K}^2=0$ et
$\mathfrak{I}/\mathfrak{K}\mathfrak{I}\neq 0$. Voir cependant (16.9.13) et
(19.1.5) pour des cas utiles où la suite prolongée reste encore exacte.

\subsection{Invariants différentiels fondamentaux d'un morphisme de préschémas}
\label{IV.16.3}

\begin{definition}[16.3.1]
\label{IV.16.3.1}
Soient $f:X\to S$ un morphisme de préschémas, $\Delta_f:X\to X\times_S X$
le morphisme diagonal correspondant, qui est une immersion (I, 5.3.9 et
$\mathbf{Err}_{\mathrm{III}}$, 10). On désigne par $\mathscr{P}_f^n$ ou
$\mathscr{P}_{X/S}^n$, et l'on appelle faisceau des parties principales d'ordre $n$
du $S$-préschéma $X$, le faisceau d'anneaux $\mathcal{O}_X$-augmentés, $n$-ème
invariant normal de $\Delta_f$ (16.1.2). On pose
$\mathscr{P}_f^\infty=\mathscr{P}_{X/S}^\infty=
\varprojlim_n\mathscr{P}_{X/S}^n$,
$\mathscr{G}r_n(\mathscr{P}_f)=\mathscr{G}r_n(\mathscr{P}_{X/S})=
\mathscr{G}r_n(\Delta_f)$ (16.1.2); le $\mathcal{O}_X$-Module
$\mathscr{G}r_1(\Delta_f)$, Idéal d'augmentation de $\mathscr{P}_{X/S}^1$, est
aussi noté $\Omega_f^1$ ou $\Omega_{X/S}^1$, et appelé le $\mathcal{O}_X$-Module
des 1-différentielles de $f$, ou de $X$ par rapport à $S$, ou du $S$-préschéma $X$.
\end{definition}

Il résulte de cette définition que $\mathscr{P}_{X/S}^0$ s'identifie
canoniquement à $\mathcal{O}_X$ (16.1.2).

On a (16.1.2.2) un homomorphisme canonique surjectif de
$\mathcal{O}_X$-Algèbres graduées
\[
  \mathscr{S}_{\mathcal{O}_X}^{\bullet}(\Omega_{X/S}^1)
  \longrightarrow\mathscr{G}r_\bullet(\mathscr{P}_{X/S}).
  \tag{16.3.1.1}
  \label{IV.16.3.1.1}
\]
Il résulte aussi de la définition (16.3.1) que pour tout ouvert $U$ de $X$, on a
$\mathscr{P}_{f|U}^n=\mathscr{P}_f^n|U$,
$\mathscr{P}_{f|U}^\infty=\mathscr{P}_f^\infty|U$,
$\mathscr{G}r_n(\mathscr{P}_{f|U})=\mathscr{G}r_n(\mathscr{P}_f)|U$,
$\Omega_{f|U}^1=\Omega_f^1|U$ (autrement dit, les notions introduites sont
locales sur $X$).

\begin{env}[16.3.2]
\label{IV.16.3.2}
Notons $p_1$, $p_2$ les deux projections canoniques du produit $X\times_S X$;
comme $\Delta_f$ est une $X$-section de $X\times_S X$ à la fois pour les
morphismes $p_1$ et $p_2$, chacun de ces morphismes définit, pour tout $n$, un
homomorphisme de faisceaux d'anneaux $\mathcal{O}_X\to\mathscr{P}_{X/S}^n$,
inverse à droite de l'augmentation $\mathscr{P}_{X/S}^n\to\mathcal{O}_X$
(16.1.7); on peut encore dire que l'on définit ainsi sur
$\mathscr{P}_{X/S}^n$ deux structures de $\mathcal{O}_X$-Algèbres augmentées
quasi-cohérentes; les structures correspondantes de $\mathcal{O}_X$-Module sur
$\mathscr{G}r_n(\mathscr{P}_{X/S})$ sont les mêmes. On a de même, par passage à
la limite, deux structures de $\mathcal{O}_X$-Algèbre sur
$\mathscr{P}_{X/S}^\infty$.
\end{env}

\begin{env}[16.3.3]
\label{IV.16.3.3}
Le morphisme $s=(p_2,p_1)_S:X\times_S X\to X\times_S X$ est un automorphisme
involutif de $X\times_S X$, appelé la symétrie canonique, tel que
\[
  p_1\circ s=p_2,\qquad p_2\circ s=p_1,\qquad s\circ\Delta_f=\Delta_f.
  \tag{16.3.3.1}
  \label{IV.16.3.3.1}
\]
Si l'on pose $s=(\rho,\lambda)$, $p_i=(\pi_i,\mu_i)$ ($i=1,2$),
$\Delta_f=(\delta,\nu)$, $\lambda^\#$ est donc un isomorphisme de
$\rho^*(\pi_1^*(\mathcal{O}_X))$ sur $\pi_2^*(\mathcal{O}_X)$, et
$\delta^*(\lambda^\#)$ laisse invariant $\delta^*(\mathcal{O}_{X\times_S X})$
et le noyau $\mathscr{I}$ de l'homomorphisme
$\nu^\#:\delta^*(\mathcal{O}_{X\times_S X})\to\mathcal{O}_X$. Par suite:
\end{env}

\begin{proposition}[16.3.4]
\label{IV.16.3.4}
L'homomorphisme $\sigma=\delta^*(\lambda^\#)$ déduit de $s$ (et appelé encore
symétrie canonique) est un automorphisme involutif du système projectif
$(\mathscr{P}_{X/S}^n)$ de faisceaux
\oldpage[IV-4]{15}
d'anneaux $\mathcal{O}_X$-augmentés, et par suite aussi de leur limite projective
$\mathscr{P}_{X/S}^\infty$. Cet automorphisme permute les deux structures de
$\mathcal{O}_X$-Algèbre sur les $\mathscr{P}_{X/S}^n$ et sur
$\mathscr{P}_{X/S}^\infty$.
\end{proposition}

\begin{env}[16.3.5]
\label{IV.16.3.5}
Dans ce qui suit, les deux structures de $\mathcal{O}_X$-Algèbre définies sur les
$\mathscr{P}_{X/S}^n$ et sur $\mathscr{P}_{X/S}^\infty$ joueront des rôles très
différents: nous conviendrons désormais, sauf mention expresse du contraire,
que lorsque $\mathscr{P}_{X/S}^n$ ou $\mathscr{P}_{X/S}^\infty$ sera considéré
comme une $\mathcal{O}_X$-Algèbre, c'est de la structure d'Algèbre définie par
$p_1$ qu'il s'agira.
\end{env}

Pour tout ouvert $U$ de $X$ et toute section $t\in\Gamma(U,\mathcal{O}_X)$, on
notera simplement $t.1$ ou même $t$ l'image de $t$ par l'homomorphisme
structural $\Gamma(U,\mathcal{O}_X)\to\Gamma(U,\mathscr{P}_{X/S}^n)$ (resp.
$\Gamma(U,\mathcal{O}_X)\to\Gamma(U,\mathscr{P}_{X/S}^\infty)$)
(c'est-à-dire l'homomorphisme correspondant à $p_1$).

\begin{definition}[16.3.6]
\label{IV.16.3.6}
On désigne par $d_f^n$, ou $d_{X/S}^n$ (resp. $d_f^\infty$, ou
$d_{X/S}^\infty$), ou simplement $d^n$ (resp. $d^\infty$), l'homomorphisme de
faisceaux d'anneaux $\mathcal{O}_X\to\mathscr{P}_f^n=\mathscr{P}_{X/S}^n$
(resp. $\mathcal{O}_X\to\mathscr{P}_f^\infty=\mathscr{P}_{X/S}^\infty$)
déduit de $p_2$. Pour tout ouvert $U$ de $X$, et tout
$t\in\Gamma(U,\mathcal{O}_X)$, $d^nt$ (resp. $d^\infty t$) est appelé la partie
principale d'ordre $n$ (resp. partie principale d'ordre infini) de $t$. On pose
$dt=d^1t-t$, et on dit que $dt$ est la différentielle de $t$ (élément de
$\Gamma(U,\Omega_{X/S}^1)$, aussi noté $d_{X/S}(t)$).
\end{definition}

Il résulte aussitôt de cette définition que l'on a
\[
  d(t_1t_2)=t_1dt_2+t_2dt_1
  \tag{16.3.6.1}
  \label{IV.16.3.6.1}
\]
quels que soient $t_1,t_2$ dans $\Gamma(U,\mathcal{O}_X)$, autrement dit, $d$
est une dérivation de l'anneau $\Gamma(U,\mathcal{O}_X)$ dans le
$\Gamma(U,\mathcal{O}_X)$-module $\Gamma(U,\Omega_{X/S}^1)$.

Dans toutes les notations introduites dans (16.3.1) et (16.3.6), on remplacera
parfois $S$ par $A$ lorsque $S=\operatorname{Spec}(A)$.

\begin{env}[16.3.7]
\label{IV.16.3.7}
Supposons en particulier que $S=\operatorname{Spec}(A)$ et
$X=\operatorname{Spec}(B)$ soient des schémas affines, $B$ étant donc une
$A$-algèbre. Alors $\Delta_f$ correspond à l'homomorphisme canonique surjectif
$\pi:B\otimes_A B\to B$ tel que $\pi(b\otimes b')=bb'$, de noyau
$\mathfrak{I}=\mathfrak{I}_{B/A}$ (0, 20.4.1); $\mathscr{P}_f^n$ est le faisceau
structural du préschéma $\operatorname{Spec}(P_{B/A}^n)$, où
\[
  P_{B/A}^n=(B\otimes_A B)/\mathfrak{I}^{n+1};
\]
$\mathscr{G}r_\bullet(\mathscr{P}_f)$ est le $\mathcal{O}_X$-Module
quasi-cohérent correspondant au $B$-module gradué
\[
  \operatorname{gr}_{\mathfrak{I}}^\bullet(B\otimes_A B)
  =\bigoplus_{n\geq 0}(\mathfrak{I}^n/\mathfrak{I}^{n+1});
\]
en particulier $\Omega_f^1=\Omega_{X/S}^1$ est le $\mathcal{O}_X$-Module
quasi-cohérent correspondant au $B$-module des 1-différentielles de $B$ par
rapport à $A$, $\Omega_{B/A}^1$ (0, 20.4.3). Les morphismes projections
$p_1:X\times_S X\to X$, $p_2:X\times_S X\to X$ correspondent aux deux
homomorphismes d'anneaux $j_1:B\to B\otimes_A B$, $j_2:B\to B\otimes_A B$
tels que $j_1(b)=b\otimes 1$, $j_2(b)=1\otimes b$, de sorte que (par la
convention de (16.3.5)), $P_{B/A}^n$ est toujours considérée comme une
$B$-algèbre pour l'homomorphisme composé
$B\xrightarrow{j_1}B\otimes_A B\to P_{B/A}^n$; l'homomorphisme d'anneaux
$B\xrightarrow{j_2}B\otimes_A B\to P_{B/A}^n$ se note $d_{B/A}^n$ et correspond
à $d_{X/S}^n$ opérant sur $\Gamma(X,\mathcal{O}_X)$; pour tout $t\in B$, $dt$
est égal à $d_{B/A}t$, défini dans (0, 20.4.6) (cf. $\mathbf{Err}_{\mathrm{IV}}$,
11).

Si $\pi_n:B\otimes_A B\to P_{B/A}^n$ est l'homomorphisme canonique, on a donc,
en vertu des définitions précédentes
\[
  \pi_n(b\otimes b')=b.\pi_n(1\otimes b')=b.d_{B/A}^n(b')
  \quad\text{pour }b\in B,\ b'\in B.
  \tag{16.3.7.1}
  \label{IV.16.3.7.1}
\]
\end{env}

\oldpage[IV-4]{16}
\begin{proposition}[16.3.8]
\label{IV.16.3.8}
L'image de l'homomorphisme $d_{X/S}^n:\mathcal{O}_X\to\mathscr{P}_{X/S}^n$
engendre le $\mathcal{O}_X$-Module $\mathscr{P}_{X/S}^n$.
\end{proposition}

On se ramène aussitôt au cas où $X=\operatorname{Spec}(B)$ et
$S=\operatorname{Spec}(A)$ sont affines et la proposition résulte de (16.3.7.1)
puisque $\pi_n$ est surjectif. On notera qu'en général $d_{X/S}^n$ n'est pas
surjectif (même déjà pour $n=1$).

\begin{proposition}[16.3.9]
\label{IV.16.3.9}
Supposons que $f:X\to S$ soit un morphisme localement de type fini. Alors les
$\mathscr{P}_f^n$ et les $\mathscr{G}r_n(\mathscr{P}_f)$ sont des
$\mathcal{O}_X$-Modules quasi-cohérents de type fini.
\end{proposition}

Cela résulte de (16.1.6) et du fait que $\Delta_f$ est localement de présentation
finie (I, 4.3.1).

\subsection{Propriétés fonctorielles des invariants différentiels}
\label{IV.16.4}

\begin{env}[16.4.1]
\label{IV.16.4.1}
Considérons un diagramme commutatif de morphismes de préschémas
\[
  \xymatrix{
    X \ar[d]_-{f} & X' \ar[l]_-u \ar[d]^-{f'} \\
    S & S' \ar[l]^-w
  }
  \tag{16.4.1.1}
  \label{IV.16.4.1.1}
\]
On en déduit un diagramme commutatif
\[
  \xymatrix{
    X \ar[d]_-{\Delta_f} & X' \ar[l]_-u \ar[d]^-{\Delta_{f'}} \\
    X\times_S X & X'\times_{S'}X' \ar[l]^-v
  }
\]
où $v$ est l'homomorphisme composé (I, 5.3.5 et 5.3.15)
\[
  X'\times_{S'}X'
  \xrightarrow{(p'_1,p'_2)_S}X'\times_S X'
  \xrightarrow{u\times_S u}X\times_S X.
  \tag{16.4.1.2}
  \label{IV.16.4.1.2}
\]
On déduit donc de $u$ et $v$, comme il a été expliqué dans (16.2.1), des
homomorphismes de faisceaux d'anneaux augmentés
\[
  \nu_n:\rho^*(\mathscr{P}_{X/S}^n)\longrightarrow
  \mathscr{P}_{X'/S'}^n
  \tag{16.4.1.3}
  \label{IV.16.4.1.3}
\]
(où l'on a posé $u=(\rho,\lambda)$); ces homomorphismes forment un système
projectif, et donnent donc à la limite un homomorphisme de faisceaux d'anneaux
augmentés
\[
  \nu_\infty:\rho^*(\mathscr{P}_{X/S}^\infty)\longrightarrow
  \mathscr{P}_{X'/S'}^\infty;
  \tag{16.4.1.4}
  \label{IV.16.4.1.4}
\]
d'autre part, par passage aux quotients, les homomorphismes $\nu_n$ donnent un
di-homomorphisme d'Algèbres graduées (relatif à $\lambda^\#$):
\[
  \operatorname{gr}(u):\rho^*(\mathscr{G}r_\bullet(\mathscr{P}_{X/S}))
  \longrightarrow\mathscr{G}r_\bullet(\mathscr{P}_{X'/S'}).
  \tag{16.4.1.5}
  \label{IV.16.4.1.5}
\]
\end{env}

\begin{env}[16.4.2]
\label{IV.16.4.2}
\oldpage[IV-4]{17}
Si l'on a un diagramme commutatif
\[
  \xymatrix{
    X \ar[d]_-{f} & X' \ar[l]_-u \ar[d]^-{f'} &
      X'' \ar[l]_-{u'} \ar[d]^-{f''} \\
    S & S' \ar[l]^-w & S'' \ar[l]^-{w'}
  }
\]
on en déduit un diagramme commutatif
\[
  \xymatrix{
    X \ar[d]_-{\Delta_f} & X' \ar[l]_-u \ar[d]^-{\Delta_{f'}} &
      X'' \ar[l]_-{u'} \ar[d]^-{\Delta_{f''}} \\
    X\times_S X & X'\times_{S'}X' \ar[l]^-v &
      X''\times_{S''}X'' \ar[l]^-{v'}
  }
\]
où $v'$ est défini à partir de $u'$, $w'$, $f'$, $f''$ comme $v$ à partir de
$u$, $w$, $f$, $f'$. On vérifie aussitôt que si $u''=u\circ u'$,
$w''=w\circ w'$, alors le morphisme composé $v\circ v'$ est égal au morphisme
$v''$ défini à partir de $u''$, $w''$, $f$, $f''$ comme $v$ à partir de $u$,
$w$, $f$, $f'$. Si l'on pose $u'=(\rho',\lambda')$,
$u''=(\rho'',\lambda'')$ il résulte alors de (16.2.1) que l'homomorphisme
$\nu_n'':\rho''^*(\mathscr{P}_{X/S}^n)\to\mathscr{P}_{X''/S''}^n$ est égal au
composé
\[
  \rho'^*(\rho^*(\mathscr{P}_{X/S}^n))
  \xrightarrow{\rho'^*(\nu_n^\#)}
  \rho'^*(\mathscr{P}_{X'/S'}^n)
  \xrightarrow{\nu_n'}\mathscr{P}_{X''/S''}^n
\]
et l'on a des propriétés de transitivité analogues pour les homomorphismes
(16.4.1.4) et (16.4.1.5), ce qui permet de dire que les
$\mathscr{P}_{X/S}^n$, $\mathscr{P}_{X/S}^\infty$ et
$\mathscr{G}r_\bullet(\mathscr{P}_{X/S})$ dépendent fonctoriellement de $f$.
\end{env}

\begin{env}[16.4.3]
\label{IV.16.4.3}
On vérifie aussitôt (par exemple en se ramenant au cas affine à l'aide de
(16.3.7)) qu'avec les notations de (16.4.1), le diagramme
\[
  \xymatrix{
    \rho^*(\mathcal{O}_X) \ar[r]^{\lambda^\#} \ar[d] &
      \mathcal{O}_{X'} \ar[d] \\
    \rho^*(\mathscr{P}_{X/S}^n) \ar[r]_{\nu_n} &
      \mathscr{P}_{X'/S'}^n
  }
  \tag{16.4.3.1}
  \label{IV.16.4.3.1}
\]
où les flèches verticales sont celles définissant les structures d'Algèbre
choisies dans (16.3.5) (c'est-à-dire celles provenant des premières projections)
est commutatif; il en est de même du diagramme
\[
  \xymatrix{
    \rho^*(\mathcal{O}_X) \ar[r]^{\lambda^\#}
      \ar[d]_{\rho^*(d_{X/S}^n)} &
      \mathcal{O}_{X'} \ar[d]^{d_{X'/S'}^n} \\
    \rho^*(\mathscr{P}_{X/S}^n) \ar[r]_{\nu_n} &
      \mathscr{P}_{X'/S'}^n
  }
  \tag{16.4.3.2}
  \label{IV.16.4.3.2}
\]
\oldpage[IV-4]{18}
les flèches verticales définissant ici les structures d'Algèbre provenant des
secondes projections; d'ailleurs, si $\sigma$ et $\sigma'$ sont les symétries
canoniques correspondant à $f$ et $f'$ (16.3.4), on a
\[
  \nu_n\circ\rho^*(\sigma)=\sigma'\circ\nu_n
\]
qui fait passer d'un des diagrammes précédents à l'autre. On déduit donc de
(16.4.3.1) un homomorphisme canonique de $\mathcal{O}_{X'}$-Algèbres augmentées
\[
  P^n(u):u^*(\mathscr{P}_{X/S}^n)
  =\mathscr{P}_{X/S}^n\otimes_{\mathcal{O}_X}\mathcal{O}_{X'}
  \longrightarrow\mathscr{P}_{X'/S'}^n
  \tag{16.4.3.3}
  \label{IV.16.4.3.3}
\]
et il résulte de (16.4.3.2) que le diagramme
\[
  \xymatrix{
    \mathcal{O}_{X'} \ar[r]^{\operatorname{id}}
      \ar[d]_{u^*(d_{X/S}^n)} &
      \mathcal{O}_{X'} \ar[d]^{d_{X'/S'}^n} \\
    u^*(\mathscr{P}_{X/S}^n) \ar[r]_{P^n(u)} &
      \mathscr{P}_{X'/S'}^n
  }
  \tag{16.4.3.4}
  \label{IV.16.4.3.4}
\]
est commutatif. On en déduit un homomorphisme de $\mathcal{O}_{X'}$-Algèbres
graduées
\[
  \operatorname{Gr}_\bullet(u):u^*(\mathscr{G}r_\bullet(\mathscr{P}_{X/S}))
  \longrightarrow\mathscr{G}r_\bullet(\mathscr{P}_{X'/S'})
  \tag{16.4.3.5}
  \label{IV.16.4.3.5}
\]
et en particulier un homomorphisme de $\mathcal{O}_{X'}$-Modules
\[
  \operatorname{Gr}_1(u):
  \Omega_{X/S}^1\otimes_{\mathcal{O}_X}\mathcal{O}_{X'}
  \longrightarrow\Omega_{X'/S'}^1
  \tag{16.4.3.6}
  \label{IV.16.4.3.6}
\]
donnant lieu à un diagramme commutatif
\[
  \xymatrix{
    \mathcal{O}_{X'} \ar[r]^{\operatorname{id}}
      \ar[d]_{d_{X/S}\otimes 1} &
      \mathcal{O}_{X'} \ar[d]^{d_{X'/S'}} \\
    \Omega_{X/S}^1\otimes_{\mathcal{O}_X}\mathcal{O}_{X'} \ar[r] &
      \Omega_{X'/S'}^1
  }
  \tag{16.4.3.7}
  \label{IV.16.4.3.7}
\]
\end{env}

\begin{env}[16.4.4]
\label{IV.16.4.4}
Lorsque $S=\operatorname{Spec}(A)$, $S'=\operatorname{Spec}(A')$,
$X=\operatorname{Spec}(B)$, $X'=\operatorname{Spec}(B')$ sont affines, de sorte
que l'on a un diagramme commutatif d'homomorphismes d'anneaux
\[
  \xymatrix{
    B \ar[r] & B' \\
    A \ar[u] \ar[r] & A' \ar[u]
  }
\]
l'image de $\mathfrak{I}_{B/A}$ dans $B'\otimes_{A'}B'$ est contenue dans
$\mathfrak{I}_{B'/A'}$, et l'homomorphisme $\nu_n$ correspond à
l'homomorphisme d'anneaux $P_{B/A}^n\to P_{B'/A'}^n$ déduit de l'homomorphisme
$B\otimes_A B\to B'\otimes_{A'}B'$ par passage aux quotients. L'homomorphisme
(16.4.3.6) correspond à l'homomorphisme défini dans (0, 20.5.4.1), et le
diagramme commutatif (16.4.3.7) au diagramme (0, 20.5.4.2).
\end{env}

\oldpage[IV-4]{19}
\begin{proposition}[16.4.5]
\label{IV.16.4.5}
Supposons que $X'=X\times_S S'$, $f'$ et $u$ étant les projections canoniques.
Alors les homomorphismes canoniques $P^n(u)$ (16.4.3.3) et
$\operatorname{Gr}_1(u)$ (16.4.3.6) sont bijectifs.
\end{proposition}

On a en effet alors $X'\times_{S'}X'=(X\times_S X)\times_S S'$, et on peut donc
appliquer (16.2.3, (ii)) en remplaçant $g$ par la première projection
$p_1:X\times_S X\to X$ et $f$ par la diagonale $\Delta_f$.

On notera aussi que sous les hypothèses de (16.4.5) l'homomorphisme
$\operatorname{Gr}_\bullet(u)$ (16.4.3.5) est surjectif, mais non bijectif en
général. Toutefois (16.2.4):

\begin{corollary}[16.4.6]
\label{IV.16.4.6}
Les hypothèses étant celles de (16.4.5), supposons de plus que $w:S'\to S$ soit
plat (resp. que les $\mathscr{G}r_n(\mathscr{P}_{X/S})$ soient des
$\mathcal{O}_X$-Modules plats pour $n\leq m$); alors l'homomorphisme
\[
  \operatorname{Gr}_n(u):u^*(\mathscr{G}r_n(\mathscr{P}_{X/S}))
  \longrightarrow\mathscr{G}r_n(\mathscr{P}_{X'/S'})
\]
est bijectif pour tout $n$ (resp. pour $n\leq m$).
\end{corollary}

En effet, si $w$ est plat, il en est de même de
$v:X'\times_{S'}X'\to X\times_S X$, donc la conclusion résulte de (16.2.4).

\begin{env}[16.4.7]
\label{IV.16.4.7}
Soient $S$ un préschéma, $\mathscr{E}$ un $\mathcal{O}_S$-Module
quasi-cohérent, et posons $X=\mathbf{V}(\mathscr{E})$ (II, 1.7.8), fibré
vectoriel associé à $\mathscr{E}$, égal à
$\operatorname{Spec}(\mathscr{S}_{\mathcal{O}_S}(\mathscr{E}))$. Soit
$f:X\to S$ le morphisme structural. Pour tout ouvert $U$ de $S$, et toute
section $t\in\Gamma(U,\mathscr{E})$, $t$ s'identifie à une section de
$\mathscr{S}_{\mathcal{O}_S}(\mathscr{E})$ au-dessus de $U$; soit $t'$ son
image dans
$\Gamma(f^{-1}(U),\mathcal{O}_X)=\Gamma(U,f_*(\mathcal{O}_X))=
\Gamma(U,\mathscr{S}_{\mathcal{O}_S}(\mathscr{E}))$, et posons
\[
  \delta(t)=d_{X/S}^n(t')-t'\in
  \Gamma(f^{-1}(U),\mathscr{P}_{X/S}^n);
  \tag{16.4.7.1}
  \label{IV.16.4.7.1}
\]
il est clair que $\delta$ est un di-homomorphisme de modules (correspondant à
l'homomorphisme d'anneaux
$\Gamma(U,\mathcal{O}_S)\to\Gamma(f^{-1}(U),\mathcal{O}_X)$) de
$\Gamma(U,\mathscr{E})$ dans $\Gamma(f^{-1}(U),\mathscr{P}_{X/S}^n)$, dont
l'image appartient d'ailleurs à l'idéal d'augmentation de
$\Gamma(f^{-1}(U),\mathscr{P}_{X/S}^n)$. On en déduit (en faisant varier $U$)
un homomorphisme canonique de $\mathcal{O}_X$-Algèbres
\[
  f^*(\mathscr{S}_{\mathcal{O}_S}(\mathscr{E}))
  \longrightarrow\mathscr{P}_{X/S}^n
  \tag{16.4.7.2}
  \label{IV.16.4.7.2}
\]
et d'après la remarque précédente, si $\mathscr{K}$ est l'Idéal noyau de
l'augmentation $\mathscr{S}_{\mathcal{O}_S}(\mathscr{E})\to\mathcal{O}_S$,
l'image de $\mathscr{K}^{n+1}$ par (16.4.7.2) est nulle, si bien qu'en
factorisant par $\mathscr{K}^{n+1}$, on a finalement un homomorphisme canonique
\[
  \delta_n:f^*(\mathscr{S}_{\mathcal{O}_S}(\mathscr{E})/
  \mathscr{K}^{n+1})\longrightarrow\mathscr{P}_{X/S}^n.
  \tag{16.4.7.3}
  \label{IV.16.4.7.3}
\]
\end{env}

\begin{proposition}[16.4.8]
\label{IV.16.4.8}
Sous les conditions de (16.4.7), les homomorphismes $\delta_n$ sont bijectifs et
forment un système projectif d'isomorphismes; on en déduit un isomorphisme de
$\mathcal{O}_X$-Algèbres graduées
\[
  f^*(\mathscr{S}_{\mathcal{O}_S}^{\bullet}(\mathscr{E}))
  \longrightarrow\mathscr{G}r_\bullet(\mathscr{P}_{X/S}).
  \tag{16.4.8.1}
  \label{IV.16.4.8.1}
\]
\end{proposition}

Le fait que les homomorphismes (16.4.7.3) forment un système projectif résulte
aussitôt de leur définition. Pour prouver que ce sont des isomorphismes, il
suffira de
\oldpage[IV-4]{20}
démontrer que (16.4.8.1) est un isomorphisme, les filtrations des deux membres de
(16.4.7.3) étant finies (Bourbaki, \emph{Alg. comm.}, chap. III, \S~2,
n\textsuperscript{o} 8, cor. 3 du th. 1). Pour cela, considérons la suite exacte
scindée de $\mathcal{O}_S$-Modules
\[
  \xymatrix{
    0 \ar[r] & \mathscr{E} \ar[r]^{u} &
    \mathscr{E}\oplus\mathscr{E} \ar[r]^{v} &
    \mathscr{E} \ar[r] & 0
  }
  \tag{16.4.8.2}
  \label{IV.16.4.8.2}
\]
où, pour tout couple de sections $s$, $t$ de $\mathscr{E}$ au-dessus d'un ouvert
$U$ de $S$, on prend $u(s)=(-s,s)$ et $v(s,t)=s+t$. On a
\[
  X\times_S X=\operatorname{Spec}
  (\mathscr{S}_{\mathcal{O}_S}(\mathscr{E})\otimes_{\mathcal{O}_S}
  \mathscr{S}_{\mathcal{O}_S}(\mathscr{E}))
  =\operatorname{Spec}(\mathscr{S}_{\mathcal{O}_S}
  (\mathscr{E}\oplus\mathscr{E}))
\]
(II, 1.4.6 et 1.7.11), et le morphisme diagonal $X\to X\times_S X$ correspond
(II, 1.2.7) à l'homomorphisme de $\mathcal{O}_S$-Algèbres
$\mathscr{S}(v):\mathscr{S}_{\mathcal{O}_S}(\mathscr{E}\oplus\mathscr{E})
\to\mathscr{S}_{\mathcal{O}_S}(\mathscr{E})$ (II, 1.7.4), de sorte que si
$\mathscr{I}$ est le noyau de cet homomorphisme, on a
\[
  \mathscr{P}_{X/S}^n=f^*(\mathscr{S}_{\mathcal{O}_S}
  (\mathscr{E}\oplus\mathscr{E})/\mathscr{I}^{n+1}).
\]

La proposition sera conséquence du lemme suivant:

\begin{lemma}[16.4.8.3]
\label{IV.16.4.8.3}
Soient $Y$ un espace annelé,
$0\to\mathscr{F}'\xrightarrow{u}\mathscr{F}\xrightarrow{v}\mathscr{F}''\to0$
une suite exacte de $\mathcal{O}_Y$-Modules telle que tout point $y\in Y$
possède un voisinage ouvert $V$ tel que la suite
$0\to\mathscr{F}'|V\to\mathscr{F}|V\to\mathscr{F}''|V\to0$ soit scindée.
Soit $\mathscr{I}$ l'Idéal noyau de
$\mathscr{S}(v):\mathscr{S}_{\mathcal{O}_Y}(\mathscr{F})\to
\mathscr{S}_{\mathcal{O}_Y}(\mathscr{F}'')$, et soit
$\operatorname{gr}_{\mathscr{I}}^\bullet
(\mathscr{S}_{\mathcal{O}_Y}(\mathscr{F}))$ la $\mathcal{O}_Y$-Algèbre graduée
associée à la $\mathcal{O}_Y$-Algèbre
$\mathscr{S}_{\mathcal{O}_Y}(\mathscr{F})$ munie de la filtration
$\mathscr{I}$-préadique. Alors l'homomorphisme de $\mathcal{O}_Y$-Algèbres
graduées
\[
  \mathscr{S}_{\mathcal{O}_Y}^{\bullet}(\mathscr{F}')
  \otimes_{\mathcal{O}_Y}
  \mathscr{S}_{\mathcal{O}_Y}^{\bullet}(\mathscr{F}'')
  \longrightarrow
  \operatorname{gr}_{\mathscr{I}}^\bullet
  (\mathscr{S}_{\mathcal{O}_Y}(\mathscr{F}))
  \tag{16.4.8.4}
  \label{IV.16.4.8.4}
\]
(où le premier membre est le produit tensoriel gradué des
$\mathcal{O}_Y$-Algèbres symétriques munies de leur graduation canonique
(II, 1.7.4 et 2.1.2)), provenant de l'injection canonique
\[
  \mathscr{F}'\longrightarrow\mathscr{I}
  =\operatorname{gr}_{\mathscr{I}}^1
  (\mathscr{S}_{\mathcal{O}_Y}(\mathscr{F})),
\]
est bijectif.
\end{lemma}

L'injection $\mathscr{F}'\to\mathscr{I}$ donne en effet canoniquement un
homomorphisme de $\mathcal{O}_Y$-Algèbres graduées\linebreak
$\mathscr{S}_{\mathcal{O}_Y}^{\bullet}(\mathscr{F}')\to\allowbreak
\operatorname{gr}_{\mathscr{I}}^\bullet
(\mathscr{S}_{\mathcal{O}_Y}(\mathscr{F}))$, et comme le second membre est par
définition une $\mathscr{S}_{\mathcal{O}_Y}^{\bullet}(\mathscr{F}'')$-Algèbre
graduée, on en déduit l'homomorphisme canonique (16.4.8.4) par tensorisation du
précédent avec $\mathscr{S}_{\mathcal{O}_Y}^{\bullet}(\mathscr{F}'')$. Pour
prouver le lemme, on peut, la question étant locale, se borner au cas où
$\mathscr{F}=\mathscr{F}'\oplus\mathscr{F}''$, $u$ et $v$ étant les
homomorphismes canoniques. Alors l'Algèbre graduée
$\mathscr{S}_{\mathcal{O}_Y}^{\bullet}(\mathscr{F})$ s'identifie canoniquement au
produit tensoriel gradué
$\mathscr{S}_{\mathcal{O}_Y}^{\bullet}(\mathscr{F}')
\otimes_{\mathcal{O}_Y}
\mathscr{S}_{\mathcal{O}_Y}^{\bullet}(\mathscr{F}'')$ (II, 1.7.4), et il est
immédiat que $\mathscr{I}$ est alors l'Idéal
$\mathscr{I}'\otimes_{\mathcal{O}_Y}
\mathscr{S}_{\mathcal{O}_Y}^{\bullet}(\mathscr{F}'')$, où $\mathscr{I}'$ est
l'Idéal d'augmentation de
$\mathscr{S}_{\mathcal{O}_Y}^{\bullet}(\mathscr{F}')$, c'est-à-dire la somme
(directe) des $\mathscr{S}_{\mathcal{O}_Y}^m(\mathscr{F}')$ pour $m\geq1$. On en
conclut que $\mathscr{I}^n=\mathscr{I}'^n\otimes_{\mathcal{O}_Y}
\mathscr{S}_{\mathcal{O}_Y}^{\bullet}(\mathscr{F}'')$, où cette fois
$\mathscr{I}'^n$ est la somme (directe) des
$\mathscr{S}_{\mathcal{O}_Y}^m(\mathscr{F}')$ pour $m\geq n$; on a par suite
$\mathscr{I}^n/\mathscr{I}^{n+1}=
\mathscr{S}_{\mathcal{O}_Y}^n(\mathscr{F}')\otimes_{\mathcal{O}_Y}
\mathscr{S}_{\mathcal{O}_Y}^{\bullet}(\mathscr{F}'')$, ce qui prouve que
(16.4.8.4) est bijectif.

\oldpage[IV-4]{21}
Ce lemme étant démontré, il reste à voir que l'homomorphisme (16.4.8.1) est
bien l'image par $f^*$ de l'homomorphisme (16.4.8.4) correspondant à la suite
exacte (16.4.8.2); on constate aisément que cela résulte de la définition de
$u$ (16.4.8.2) et de celle de $\delta$ (16.4.7.1), compte tenu de la définition
de la structure de $\mathcal{O}_X$-Algèbre sur $\mathscr{P}_{X/S}^n$ et de
celle de $d_{X/S}^n$ (16.3.5 et 16.3.6).

En particulier:

\begin{corollary}[16.4.9]
\label{IV.16.4.9}
Sous les conditions de (16.4.7), on a un isomorphisme canonique
\[
  \operatorname{gr}_1(\delta):f^*(\mathscr{E})
  \xrightarrow{\sim}\Omega_{X/S}^1.
  \tag{16.4.9.1}
  \label{IV.16.4.9.1}
\]
\end{corollary}

\begin{corollary}[16.4.10]
\label{IV.16.4.10}
Si $S=\operatorname{Spec}(A)$, $\mathscr{E}=\mathcal{O}_S^m$, de sorte que
\[
  X=\operatorname{Spec}(A[T_1,\ldots,T_m]),
\]
$\mathscr{P}_{X/S}^n$ s'identifie canoniquement à la
$\mathcal{O}_X$-Algèbre correspondant à la $A[T_1,\ldots,T_m]$-algèbre quotient
\[
  A[T_1,\ldots,T_m,U_1,\ldots,U_m]/\mathfrak{K}^{n+1},
\]
où les $U_i$ ($1\leq i\leq m$) sont $m$ nouvelles indéterminées et
$\mathfrak{K}$ est l'idéal engendré par $U_1,\ldots,U_m$.
\end{corollary}

On retrouve ainsi en particulier la structure de $\Omega_{X/S}^1$ dans ce cas
(0, 20.5.13).

Notons en outre que $d_{X/S}^n$ fait alors correspondre à un polynôme
$F(T_1,\ldots,T_m)$ la classe mod.\ $\mathfrak{K}^{n+1}$ de
$F(T_1+U_1,\ldots,T_m+U_m)$, comme il résulte de la définition (16.4.7.1).

\begin{proposition}[16.4.11]
\label{IV.16.4.11}
Soient $f:X\to S$ un morphisme, $g:S\to X$ une $S$-section de $X$,
$S_g^{(n)}$ le $n$-ème voisinage infinitésimal de $S$ pour l'immersion $g$
(16.1.2). Il existe alors un isomorphisme et un seul de
$\mathcal{O}_S$-Algèbres
\[
  \varpi_n:g^*(\mathscr{P}_{X/S}^n)\longrightarrow
  \mathcal{O}_{S_g^{(n)}}
  \tag{16.4.11.1}
  \label{IV.16.4.11.1}
\]
(pour la structure de $\mathcal{O}_S$-Algèbre sur
$\mathcal{O}_{S_g^{(n)}}$ définie par $f$ (16.1.7)), rendant commutatif le
diagramme
\[
  \xymatrix{
    \mathcal{O}_S=g^*(\mathcal{O}_X)
      \ar[rr]^-{\lambda_n}\ar[dr]_-{g^*(d_{X/S}^n)}
      &&\mathcal{O}_{S_g^{(n)}}\\
    &g^*(\mathscr{P}_{X/S}^n)\ar[ur]_-{\varpi_n}
  }
  \tag{16.4.11.2}
  \label{IV.16.4.11.2}
\]
(où $\lambda_n$ est l'homomorphisme structural).
\end{proposition}

En vertu de (I, 5.3.7), où l'on remplace $X$, $Y$, $S$ par $S$, $X$, $S$
respectivement et $f$ par $g$, les diagrammes
\[
  \xymatrix@=3pc{
    S\ar[r]^-g\ar[d]_-g&X\ar[d]^-{\Delta_f}
      &S\ar[r]^-g\ar[d]_-g&X\ar[d]^-{\Delta_f}\\
    X\ar[r]_{(g\circ f,1_X)_S}&X\times_S X
      &X\ar[r]_{(1_X,g\circ f)_S}&X\times_S X
  }
  \tag{16.4.11.3}
  \label{IV.16.4.11.3}
\]
\oldpage[IV-4]{22}
identifient $S$ au produit des $(X\times_S X)$-préschémas $X$ et $X$ pour les
morphismes $\Delta_f$ et $(g\circ f,1_X)_S$ (resp.\ $(1_X,g\circ f)_S$).
D'autre part, les diagrammes
\[
  \xymatrix@=3pc{
    X\ar[r]^{(g\circ f,1_X)_S}\ar[d]_-f&X\times_S X\ar[d]^-{p_1}
      &X\ar[r]^{(1_X,g\circ f)_S}\ar[d]_-f&X\times_S X\ar[d]^-{p_2}\\
    S\ar[r]_{g}&X&S\ar[r]_{g}&X
  }
  \tag{16.4.11.4}
  \label{IV.16.4.11.4}
\]
identifient $X$ au produit des $X$-préschémas $S$ et $X\times_S X$ pour les
morphismes $g$ et $p_1$ (resp.\ $p_2$) (cas particulier de la formule
d'associativité (I, 3.3.9.1)). On peut dire que $\Delta_f$, considéré comme
$X$-section de $X\times_S X$ (relativement à $p_1$ ou $p_2$) joue le rôle
d'une \emph{section universelle} pour les $S$-sections de $X$: chacune de ces
sections $g$ s'en déduit en effet par le \emph{changement de base}
$(g\circ f,1_X)_S:X\to X\times_S X$. La définition de l'homomorphisme
$\varpi_n$ et le fait qu'il est bijectif résultent donc de ces remarques et de
(16.2.3, (ii)) appliqué au premier diagramme (16.4.11.4). La commutativité du
diagramme (16.4.11.2) résulte de même de (16.2.3, (ii)) appliqué cette fois au
second diagramme (16.4.11.4). Pour expliciter $\varpi_n$, on peut se borner au
cas où $g$ est une immersion \emph{fermée}: en effet, pour tout $s\in S$, il y
a un voisinage ouvert $W$ de $s$ dans $S$ tel que $g(W)$ soit fermé dans un
ouvert $U$ de $X$, et il est clair que $g|W$ est une $W$-section du morphisme
$U\cap f^{-1}(W)\to W$, restriction de $f$, et $g(W)$ est \emph{a fortiori}
fermé dans $U\cap f^{-1}(W)$. On peut donc supposer que $S$ est un
sous-préschéma fermé de $X$ défini par un Idéal quasi-cohérent $\mathscr{K}$.
Alors les définitions précédentes montrent que si $W$ est un ouvert de $S$,
$t$ une section de $\mathcal{O}_X$ au-dessus de $f^{-1}(W)$,
$\varpi_n(d^n t|W)$ est égal à l'image canonique de $t$ dans
$\Gamma(W,(\mathcal{O}_X/\mathscr{K}^{n+1})|W)$. L'unicité de $\varpi_n$
résulte donc de ce que l'image de $\mathcal{O}_X$ par $d_{X/S}^n$ engendre le
$\mathcal{O}_X$-Module $\mathscr{P}_{X/S}^n$ (16.3.8).

\begin{corollary}[16.4.12]
\label{IV.16.4.12}
Soient $k$ un corps, $X$ un $k$-préschéma, $x$ un point de $X$ rationnel sur
$k$. Alors $(\mathscr{P}_{X/k}^n)_x\otimes_{\mathcal{O}_x}k(x)$ est
canoniquement isomorphe (en tant que $k(x)$-algèbre augmentée) à
$\mathcal{O}_x/\mathfrak{m}_x^{n+1}$.
\end{corollary}

Il suffit de considérer l'unique $k$-section $g$ de $X$ telle que
$g(\operatorname{Spec}(k))=\{x\}$.

\begin{corollary}[16.4.13]
\label{IV.16.4.13}
Soient $f:X\to S$ un morphisme, $s$ un point de $S$,
$X_s=X\times_S\operatorname{Spec}(k(s))$ la fibre de $f$ en $s$. Si
$x\in X_s$ est rationnel sur $k(s)$,
$(\mathscr{P}_{X/S}^n)_x\otimes_{\mathcal{O}_s}k(s)$ est canoniquement
isomorphe à $\mathcal{O}_{X_s,x}/\mathfrak{m}'_x{}^{\,n+1}$, où
$\mathfrak{m}'_x$ est l'idéal maximal de $\mathcal{O}_{X_s,x}$; de façon
précise, cet isomorphisme fait correspondre à $(d^n t)_x\otimes1$ (où $t$ est
une section de $\mathcal{O}_X$ au-dessus d'un voisinage ouvert de $x$ dans
$X$) la classe de $t_x\otimes1$ modulo $\mathfrak{m}'_x{}^{\,n+1}$.
\end{corollary}

Cela résulte de (16.4.5) et (16.4.12).

Les corollaires précédents justifient la terminologie de «faisceau des parties
principales d'ordre $n$».

\begin{proposition}[16.4.14]
\label{IV.16.4.14}
Soit $\rho:A\to B$ un homomorphisme d'anneaux, $S$ une partie
multiplicativement stable de $B$. Alors les homomorphismes canoniques
\[
  S^{-1}P_{B/A}^n\longrightarrow P_{S^{-1}B/A}^n
  \tag{16.4.14.1}
  \label{IV.16.4.14.1}
\]
\oldpage[IV-4]{23}
déduits des homomorphismes canoniques
$P_{B/A}^n\to P_{S^{-1}B/A}^n$ (16.4.4), forment un système projectif et sont
bijectifs.
\end{proposition}

Il suffit de remarquer que
\[
  S^{-1}((B\otimes_A B)/\mathfrak{I}^{n+1})
  =S^{-1}(B\otimes_A B)/(S^{-1}\mathfrak{I})^{n+1}
\]
par platitude, et que
$S^{-1}(B\otimes_A B)=(S^{-1}B)\otimes_A(S^{-1}B)$ (I, 1.3.4).

\begin{corollary}[16.4.15]
\label{IV.16.4.15}
Les notations étant celles de (16.4.14), soit $R$ une partie multiplicative de
$A$ telle que $\rho(R)\subset S$. Alors on a des isomorphismes canoniques
\[
  S^{-1}P_{B/A}^n\xrightarrow{\sim}P_{S^{-1}B/R^{-1}A}^n
  \tag{16.4.15.1}
  \label{IV.16.4.15.1}
\]
formant un système projectif.
\end{corollary}

Il suffit évidemment de définir des isomorphismes canoniques
\[
  P_{S^{-1}B/A}^n\xrightarrow{\sim}P_{S^{-1}B/R^{-1}A}^n,
  \tag{16.4.15.2}
  \label{IV.16.4.15.2}
\]
c'est-à-dire que l'on est ramené au cas où $\rho(R)$ est formé d'éléments
\emph{inversibles} de $B$. Mais alors l'isomorphisme (16.4.15.2) est simplement
déduit de l'isomorphisme canonique
$B\otimes_A B\to B\otimes_{R^{-1}A}B$ par passage aux quotients
(0$_{\rm I}$, 1.5.3).

\begin{corollary}[16.4.16]
\label{IV.16.4.16}
Soient $f:X\to S$ un morphisme de préschémas, $x$ un point de $X$,
$s=f(x)$. Alors on a des isomorphismes canoniques
\[
  (\mathscr{P}_{X/S}^n)_x\xrightarrow{\sim}
  P_{\mathcal{O}_x/\mathcal{O}_s}^n
  \tag{16.4.16.1}
  \label{IV.16.4.16.1}
\]
formant un système projectif.
\end{corollary}

On déduit de là des isomorphismes pour les gradués associés, et en particulier
un isomorphisme canonique
\[
  (\Omega_{X/S}^1)_x\xrightarrow{\sim}
  \Omega_{\mathcal{O}_x/\mathcal{O}_s}^1.
  \tag{16.4.16.2}
  \label{IV.16.4.16.2}
\]

\begin{corollary}[16.4.17]
\label{IV.16.4.17}
Soient $k$ un corps, $K$ le corps des fractions rationnelles
$k(T_1,\ldots,T_r)$. Alors, pour tout entier $n$, l'homomorphisme de
$K[U_1,\ldots,U_r]$ ($U_i$ indéterminées) dans $P_{K/k}^n$ qui, à tout $U_i$,
fait correspondre $d^nT_i-T_i\cdot1$, est surjectif et définit un isomorphisme
du quotient $K[U_1,\ldots,U_r]/\mathfrak{m}^{n+1}$ (où $\mathfrak{m}$ est
l'idéal engendré par les $U_i$) sur $P_{K/k}^n$.
\end{corollary}

Cela résulte de (16.4.8), (16.4.10) et (16.4.14), où l'on fait $A=k$,
$B=k[T_1,\ldots,T_r]$ et $S=B-\{0\}$.

On retrouve ainsi le fait que les $dT_i$ forment une base du $K$-espace
vectoriel $\Omega_{K/k}^1$ (0, 20.5.10).

\begin{proposition}[16.4.18]
\label{IV.16.4.18}
Soient $f:X\to Y$, $g:Y\to Z$ deux morphismes de préschémas; et considérons
les homomorphismes canoniques de $\mathcal{O}_X$-Algèbres augmentées
(16.4.3.3)
\[
  g_{X/Y/Z}:\mathscr{P}_{X/Z}^n\longrightarrow\mathscr{P}_{X/Y}^n
  \tag{16.4.18.1}
  \label{IV.16.4.18.1}
\]
\[
  f_{X/Y/Z}:f^*(\mathscr{P}_{Y/Z}^n)
  \longrightarrow\mathscr{P}_{X/Z}^n.
  \tag{16.4.18.2}
  \label{IV.16.4.18.2}
\]
Alors $g_{X/Y/Z}$ est surjectif, et son noyau est l'Idéal engendré par l'image
par $f_{X/Y/Z}$ de l'Idéal d'augmentation de
$f^*(\mathscr{P}_{Y/Z}^n)$.
\end{proposition}

\oldpage[IV-4]{24}
Notons d'abord que $g_{X/Y/Z}$ correspond au cas où dans (16.4.3.3) on fait
$X'=X$, $S'=Y$, $S=Z$ et $u=1_X$, et $f_{X/Y/Z}$ au cas où l'on remplace
$X'$, $X$, $S$, $S'$ par $X$, $Y$, $Z$, $Z$ respectivement et $u$, $f$ par
$f$, $g$ respectivement.

On a un diagramme commutatif (I, 5.3.5)
\[
  \xymatrix{
    X\ar[r]^-{\Delta_f}\ar[dr]_-f
      &X\times_Y X\ar[d]^-p\ar[r]^-j
      &X\times_Z X\ar[d]^-{f\times_Z f}\\
    &Y\ar[r]_-{\Delta_g}&Y\times_Z Y
  }
  \tag{16.4.18.3}
  \label{IV.16.4.18.3}
\]
où $j=(1_X,1_X)_Z$ est une immersion, $j\circ\Delta_f=\Delta_{g\circ f}$ et
$p$ est le morphisme structural. Comme on peut se borner au cas où $X$, $Y$,
$Z$ sont affines, on peut supposer les immersions $\Delta_f$, $\Delta_g$ et
$j$ fermées, de sorte que $\mathcal{O}_X$ et
$\mathcal{O}_{X\times_Y X}$ s'identifient respectivement à
$\mathcal{O}_{X\times_Z X}/\mathscr{J}$ et
$\mathcal{O}_{X\times_Z X}/\mathscr{L}$, où
$\mathscr{J}\supset\mathscr{L}$ sont les deux Idéaux quasi-cohérents
correspondant respectivement aux immersions $\Delta_{g\circ f}$ et $j$. La
$\mathcal{O}_X$-Algèbre $\mathscr{P}_{X/Z}^n$ s'identifie donc à
$\mathcal{O}_{X\times_Z X}/\mathscr{J}^{n+1}$, et
$\mathscr{P}_{X/Y}^n$ s'identifie à
$\mathcal{O}_{X\times_Y X}/(\mathscr{J}/\mathscr{L})^{n+1}$, c'est-à-dire à
$\mathcal{O}_{X\times_Z X}/(\mathscr{J}^{n+1}+\mathscr{L})$, et par suite au
quotient de $\mathscr{P}_{X/Z}^n$ par
$(\mathscr{J}^{n+1}+\mathscr{L})/\mathscr{J}^{n+1}$. Mais on sait
(\emph{loc. cit.}) que $p$ et $j$ font de $X\times_Y X$ le produit des
$(Y\times_Z Y)$-préschémas $Y$ et $X\times_Z X$, donc si $\mathcal{O}_Y$ est
identifié à $\mathcal{O}_{Y\times_Z Y}/\mathscr{K}$, où $\mathscr{K}$ est
l'Idéal correspondant à $\Delta_g$, $\mathscr{L}$ est égal à
$(f\times_Z f)^*(\mathscr{K})\cdot\mathcal{O}_{X\times_Z X}$ (I, 4.4.5).
Comme $(\mathscr{J}^{n+1}+\mathscr{L})/\mathscr{J}^{n+1}$ est l'Idéal de
$\mathscr{P}_{X/Z}^n$ engendré par l'image de $\mathscr{L}$, on en déduit la
proposition.

\begin{corollary}[16.4.19]
\label{IV.16.4.19}
Avec les notations de (16.4.18), on a une suite exacte de
$\mathcal{O}_X$-Modules quasi-cohérents
\[
  \xymatrix{
    f^*(\Omega_{Y/Z}^1)\ar[r]^-{f_{X/Y/Z}}&
    \Omega_{X/Z}^1\ar[r]^-{g_{X/Y/Z}}&
    \Omega_{X/Y}^1\ar[r]&0.
  }
  \tag{16.4.19.1}
  \label{IV.16.4.19.1}
\]
\end{corollary}

Lorsque $X$, $Y$, $Z$ sont affines, on retrouve ainsi la suite exacte
(0, 20.5.7.1).

\begin{proposition}[16.4.20]
\label{IV.16.4.20}
Soient $f:Y\to Z$ un morphisme, $j:X\to Y$ une immersion fermée,
$\mathscr{K}$ l'Idéal quasi-cohérent de $\mathcal{O}_Y$ correspondant à $j$.
Alors on a
$\mathscr{P}_{X/Y}^n=\mathcal{O}_X=\mathcal{O}_Y/\mathscr{K}$,
l'homomorphisme canonique
$j_{X/Y/Z}:j^*(\mathscr{P}_{Y/Z}^n)\to\mathscr{P}_{X/Z}^n$ est surjectif, et
son noyau est l'Idéal de $j^*(\mathscr{P}_{Y/Z}^n)$ engendré par
$j^*(\mathcal{O}_Y\cdot d_{Y/Z}^n(\mathscr{K}))$ (il faut noter que
$d_{Y/Z}^n(\mathscr{K})$ est un sous-faisceau de groupes commutatifs de
$\mathscr{P}_{Y/Z}^n$, mais non un $\mathcal{O}_Y$-Module en général).
\end{proposition}

On sait (I, 5.3.8) que la diagonale
$\Delta_j:X\to X\times_Y X$ est un isomorphisme, d'où la première assertion.
Si $\varpi_1$ et $\varpi_2$ sont les deux homomorphismes d'Algèbres
$\mathcal{O}_Y\to\mathscr{P}_{Y/Z}^n$ correspondant respectivement aux deux
projections canoniques $p_1$, $p_2$ de $Y\times_Z Y$ sur $Y$, rappelons que
par définition (16.3.5 et 16.3.6) $\varpi_1$ est l'homomorphisme structural de
la $\mathcal{O}_Y$-Algèbre $\mathscr{P}_{Y/Z}^n$ et
$\varpi_2=d_{Y/Z}^n$. La $\mathcal{O}_X$-Algèbre
$j^*(\mathscr{P}_{Y/Z}^n)$ s'identifie donc à
$\mathscr{P}_{Y/Z}^n/\varpi_1(\mathscr{K})\mathscr{P}_{Y/Z}^n$ et son quotient
par l'Idéal engendré par $j^*(d_{Y/Z}^n(\mathscr{K}))$ à
\[
  \mathscr{P}_{Y/Z}^n/
  (\varpi_1(\mathscr{K})+\varpi_2(\mathscr{K}))\mathscr{P}_{Y/Z}^n.
\]
Notons maintenant qu'on a un diagramme commutatif
\oldpage[IV-4]{25}
\[
  \xymatrix{
    Y\ar[d]_-{\Delta_f}&X\ar[l]_-j\ar[d]^-{\Delta_{f\circ j}}\\
    Y\times_Z Y&X\times_Z X\ar[l]^-{j\times_Z j}
  }
\]
identifiant $X$ au produit des $(Y\times_Z Y)$-préschémas $Y$ et
$X\times_Z X$ (I, 5.3.7). Comme $j\times_Z j$ est une immersion, on déduit
donc de cette remarque et de (16.2.2) que si $\Delta_{Y/Z}^n$ et
$\Delta_{X/Z}^n$ désignent les voisinages infinitésimaux d'ordre $n$ de $Y$
et $X$ pour les immersions canoniques $\Delta_f$ et $\Delta_{f\circ j}$
respectivement, on a un diagramme
\[
  \xymatrix{
    \Delta_{Y/Z}^n\ar[d]&\Delta_{X/Z}^n\ar[l]\ar[d]\\
    Y\times_Z Y&X\times_Z X\ar[l]^-{j\times_Z j}
  }
\]
faisant de $\Delta_{X/Z}^n$ le produit des $(Y\times_Z Y)$-préschémas
$\Delta_{Y/Z}^n$ et $X\times_Z X$. On peut encore dire que
$\mathscr{P}_{X/Z}^n$ s'identifie au faisceau d'anneaux
$\mathscr{P}_{Y/Z}^n\otimes_{\mathcal{O}_{Y\times_Z Y}}
\mathcal{O}_{X\times_Z X}$. Mais on voit aussitôt (par exemple en se ramenant
au cas affine) que
\[
  \mathcal{O}_{X\times_Z X}=
  \mathcal{O}_{Y\times_Z Y}/
  (p_1^*(\mathscr{K})+p_2^*(\mathscr{K}))
  \mathcal{O}_{Y\times_Z Y}.
\]
Donc $\mathscr{P}_{X/Z}^n$ s'identifie au quotient de
$\mathscr{P}_{Y/Z}^n$ par l'Idéal engendré par l'image dans
$\mathscr{P}_{Y/Z}^n$ de $p_1^*(\mathscr{K})+p_2^*(\mathscr{K})$. Mais par
définition cet Idéal est aussi engendré par
$\varpi_1(\mathscr{K})+\varpi_2(\mathscr{K})$. C.Q.F.D.

\begin{corollary}[16.4.21]
\label{IV.16.4.21}
Soient $f:Y\to Z$ un morphisme, $j:X\to Y$ une immersion. On a une suite
exacte de $\mathcal{O}_X$-Modules quasi-cohérents
\[
  \xymatrix{
    \mathscr{N}_{X/Y}\ar[r]&j^*(\Omega_{Y/Z}^1)\ar[r]&
    \Omega_{X/Z}^1\ar[r]&0.
  }
  \tag{16.4.21.1}
  \label{IV.16.4.21.1}
\]
\end{corollary}

Lorsque $X$, $Y$ et $Z$ sont affines, on retrouve ainsi la suite exacte
(0, 20.5.12.1).

\begin{corollary}[16.4.22]
\label{IV.16.4.22}
Si $f:X\to S$ est un morphisme localement de présentation finie,
$\mathscr{P}_{X/S}^n$ et $\Omega_{X/S}^1$ sont des
$\mathcal{O}_X$-Modules quasi-cohérents de présentation finie.
\end{corollary}

On est aussitôt ramené au cas où $S=\operatorname{Spec}(A)$ est affine,
$X=\operatorname{Spec}(B)$, où
$B=A[T_1,\ldots,T_r]/\mathfrak{K}$, $\mathfrak{K}$ étant un idéal de type
fini de $C=A[T_1,\ldots,T_r]$. On applique alors (16.4.20) avec $Z=S$,
$Y=\operatorname{Spec}(C)$ et $\mathscr{K}=\widetilde{\mathfrak{K}}$. Alors
$j^*(\mathscr{P}_{Y/Z}^n)$ est un $\mathcal{O}_X$-Module libre de rang fini
(16.4.10) et l'hypothèse sur $\mathfrak{K}$ entraîne que
$j^*(\mathcal{O}_Y\cdot d_{Y/Z}^n(\mathscr{K}))$ engendre un
$\mathcal{O}_X$-Module quasi-cohérent de type fini; d'où la conclusion.

\begin{proposition}[16.4.23]
\label{IV.16.4.23}
Soient $X$, $Y$ deux $S$-préschémas, $Z=X\times_S Y$ leur produit,
$p:X\times_S Y\to X$ et $q:X\times_S Y\to Y$ les projections canoniques.
Alors l'homomorphisme canonique
\[
  p_{Z/X/S}\oplus q_{Z/Y/S}:
  p^*(\Omega_{X/S}^1)\oplus q^*(\Omega_{Y/S}^1)
  \longrightarrow\Omega_{(X\times_S Y)/S}^1
  \tag{16.4.23.1}
  \label{IV.16.4.23.1}
\]
est bijectif.
\end{proposition}

\oldpage[IV-4]{26}
Le diagramme commutatif
\[
  \xymatrix{
    Y\ar[d]_-g&X\times_S Y\ar[l]_-q\ar[d]_-h&
      X\times_S Y\ar[l]_-{\operatorname{id}}\ar[d]^-p\\
    S&S\ar[l]^-{\operatorname{id}}&X\ar[l]^-f
  }
\]
donne une factorisation de l'\emph{isomorphisme} canonique $P^n(p)$ (16.4.5)
\[
  p^*(\mathscr{P}_{X/S}^n)\longrightarrow
  \mathscr{P}_{Z/S}^n\longrightarrow\mathscr{P}_{Z/Y}^n
\]
et de même, en intervertissant les rôles de $X$ et $Y$, on a une factorisation
de l'\emph{isomorphisme} $P^n(q)$
\[
  q^*(\mathscr{P}_{Y/S}^n)\longrightarrow
  \mathscr{P}_{Z/S}^n\longrightarrow\mathscr{P}_{Z/X}^n.
\]
Ceci prouve que l'homomorphisme canonique (16.4.18.1)
\[
  p_{Z/X/S}:p^*(\mathscr{P}_{X/S}^n)\longrightarrow\mathscr{P}_{Z/S}^n
  \quad\text{(resp.\
  $q_{Z/Y/S}:q^*(\mathscr{P}_{Y/S}^n)\to\mathscr{P}_{Z/S}^n$)}
\]
est \emph{injectif}, et que le noyau de l'homomorphisme canonique surjectif
(16.4.18.2)
\[
  \mathscr{P}_{Z/S}^n\longrightarrow\mathscr{P}_{Z/Y}^n
  \quad\text{(resp.\
  $\mathscr{P}_{Z/S}^n\to\mathscr{P}_{Z/X}^n$)}
\]
est supplémentaire de l'image de $p_{Z/X/S}$ (resp.\ $q_{Z/Y/S}$). Mais
d'autre part, ce noyau est, en vertu de (16.4.18), engendré par l'image par
$q_{Z/Y/S}$ (resp.\ $p_{Z/X/S}$) de l'Idéal d'augmentation de
$q^*(\mathscr{P}_{Y/S}^n)$ (resp.\ $p^*(\mathscr{P}_{X/S}^n)$). On en conclut
la proposition en considérant le cas $n=1$.

On généralise aussitôt (16.4.23) au cas d'un produit d'un nombre fini
quelconque de $S$-préschémas.

\begin{remarks}[16.4.24]
\label{IV.16.4.24}
\begin{enumerate}
\item[(i)] Nous verrons (17.2.3) que lorsque le morphisme $f:X\to Y$ dans
(16.4.18) est \emph{lisse}, l'homomorphisme $f_{X/Y/Z}$ dans (16.4.19.1) est
localement \emph{inversible à gauche} et en particulier injectif. De même,
lorsque le morphisme $f\circ j:X\to Z$ de (16.4.20) est \emph{lisse},
l'homomorphisme de gauche dans (16.4.21.2) est localement
\emph{inversible à gauche} et \emph{a fortiori} injectif (17.2.5). Au chap. V,
nous donnerons aussi une variante, dans le cas des Modules sur les
préschémas, des «modules d'imperfection» étudiés dans (0, 20.6), et des suites
exactes où ils figurent.

\item[(ii)] Soient $X$ un espace topologique, $\mathscr{A}$ un faisceau
d'anneaux sur $X$ et $\mathscr{B}$ une $\mathscr{A}$-Algèbre sur $X$. Alors il
est clair que
\[
  U\longmapsto
  P_{\Gamma(U,\mathscr{B})/\Gamma(U,\mathscr{A})}^n
  \qquad (U\text{ ouvert dans }X)
\]
est un préfaisceau de $\Gamma(U,\mathscr{B})$-algèbres augmentées, donc le
faisceau associé $\mathscr{P}_{\mathscr{B}/\mathscr{A}}^n$ est une
$\mathscr{B}$-Algèbre augmentée. Dans le cas particulier où $X$ est un
préschéma, $f=(\psi,\theta):X\to S$ un morphisme de préschémas, il résulte
aisément de (16.4.16) et de l'exactitude du foncteur $\varinjlim$ que
$\mathscr{P}_{X/S}^n$ est canoniquement isomorphe à
$\mathscr{P}_{\mathcal{O}_X/\psi^*(\mathcal{O}_S)}^n$. Il s'ensuit que le
formalisme développé dans le présent paragraphe pourrait être considéré comme
un
\oldpage[IV-4]{27}
cas particulier d'un formalisme différentiel pour des espaces annelés munis
d'un faisceau d'algèbres sur le faisceau structural. Cependant, nous n'avons
pas voulu partir de ce point de vue, moins intuitif et moins commode pour les
applications. Il semble d'ailleurs que, pour les diverses espèces de
«variétés», la construction «globale» des $\mathscr{P}^n$ analogue à celle que
nous utilisons ici soit également mieux adaptée aux applications.
\end{enumerate}
\end{remarks}

\subsection{Faisceaux et fibrés tangents relatifs; dérivations.}
\label{IV.16.5}

\begin{env}[16.5.1]
\label{IV.16.5.1}
Soit $f=(\psi,\theta):X\to S$ un morphisme d'espaces annelés. Pour tout
$\mathcal{O}_X$-Module $\mathscr{F}$, on appelle \emph{$S$-dérivation}
(ou \emph{$(X/S)$-dérivation}, ou \emph{$f$-dérivation}) \emph{de
$\mathcal{O}_X$ dans $\mathscr{F}$} tout homomorphisme de \emph{faisceaux de
groupes additifs} $D:\mathcal{O}_X\to\mathscr{F}$, vérifiant les conditions
suivantes:
\begin{enumerate}
\item[(a)] pour tout ouvert $V$ de $X$, et tout couple de sections
$(t_1,t_2)$ de $\mathcal{O}_X$ au-dessus de $V$, on a
\[
  D(t_1t_2)=t_1D(t_2)+D(t_1)t_2;
  \tag{16.5.1.1}
  \label{IV.16.5.1.1}
\]
\item[(b)] pour tout ouvert $V$ de $X$, toute section $t$ de $\mathcal{O}_X$
au-dessus de $V$ et toute section $s$ de $\mathcal{O}_S$ au-dessus d'un ouvert
$U$ de $S$ tel que $V\subset f^{-1}(U)$, on a
\[
  D((s|V)t)=(s|V)D(t).
  \tag{16.5.1.2}
  \label{IV.16.5.1.2}
\]
\end{enumerate}
Il est clair qu'il revient au même de dire que pour tout $x\in X$,
l'homomorphisme de groupes additifs
$D_x:\mathcal{O}_x\to\mathscr{F}_x$ est une
$\mathcal{O}_{f(x)}$-dérivation.

Une autre interprétation consiste à considérer la $\mathcal{O}_X$-Algèbre
$\mathscr{D}_{\mathcal{O}_X}(\mathscr{F})$ égale à
$\mathcal{O}_X\oplus\mathscr{F}$, la structure d'Algèbre étant définie par la
condition que pour tout ouvert $V$ de $X$, le produit de deux sections de
$\mathcal{O}_X$ (resp. d'une section de $\mathcal{O}_X$ et d'une section de
$\mathscr{F}$) au-dessus de $V$ est défini par la structure d'anneau de
$\Gamma(V,\mathcal{O}_X)$ (resp. la structure de
$\Gamma(V,\mathcal{O}_X)$-module sur $\Gamma(V,\mathscr{F})$), et le produit de
deux sections de $\mathscr{F}$ au-dessus de $V$ est pris égal à 0; alors,
$\mathscr{F}$ est un Idéal de
$\mathscr{D}_{\mathcal{O}_X}(\mathscr{F})$, noyau de l'augmentation canonique
$\mathscr{D}_{\mathcal{O}_X}(\mathscr{F})\to\mathcal{O}_X$, et dire que $D$
est une $S$-dérivation de $\mathcal{O}_X$ dans $\mathscr{F}$ signifie que
$1_{\mathcal{O}_X}+D$ est un $\mathcal{O}_S$-homomorphisme d'Algèbres de
$\mathcal{O}_X$ dans $\mathscr{D}_{\mathcal{O}_X}(\mathscr{F})$, qui, composé
avec l'augmentation, donne $1_{\mathcal{O}_X}$.

Les $S$-dérivations de $\mathcal{O}_X$ dans $\mathscr{F}$ forment évidemment
un $\Gamma(X,\mathcal{O}_X)$-module
$\operatorname{Dér}_S(\mathcal{O}_X,\mathscr{F})$.

Lorsque $\mathscr{F}=\mathcal{O}_X$, une $S$-dérivation de $\mathcal{O}_X$ dans
lui-même est simplement appelée une \emph{$S$-dérivation de $\mathcal{O}_X$}.
\end{env}

\begin{proposition}[16.5.2]
\label{IV.16.5.2}
Soient $A$ un anneau, $B$ une $A$-algèbre, $L$ un $B$-module; on pose
$S=\operatorname{Spec}(A)$, $X=\operatorname{Spec}(B)$,
$\mathscr{F}=\widetilde{L}$. Alors l'application $D\mapsto\Gamma(D)$ qui, à
toute $S$-dérivation $D$ de $\mathcal{O}_X$ dans $\mathscr{F}$ fait
correspondre l'application $\Gamma(D):t\mapsto D(t)$ de $B$ dans $L$, est un
isomorphisme de $B$-modules de
$\operatorname{Dér}_S(\mathcal{O}_X,\mathscr{F})$ sur
$\operatorname{Dér}_A(B,L)$ (cf. (0, 20.1.2)).
\end{proposition}

Cela résulte aussitôt de l'interprétation donnée ci-dessus des $S$-dérivations
en termes
\oldpage[IV-4]{28}
d'homomorphismes d'Algèbres, de l'interprétation analogue donnée dans
(0, 20.1.6), et de la correspondance canonique entre homomorphismes de
$\mathcal{O}_X$-Algèbres et homomorphismes de $B$-algèbres (I, 1.3.13 et
1.3.8).

\begin{proposition}[16.5.3]
\label{IV.16.5.3}
Soit $f=(\psi,\theta):X\to S$ un morphisme de préschémas.
\begin{enumerate}
\item[(i)] La différentielle
$d_{X/S}:\mathcal{O}_X\to\Omega_{X/S}^1$ (16.3.6) est une $S$-dérivation.
\item[(ii)] Pour tout $\mathcal{O}_X$-Module $\mathscr{F}$, l'application
$u\mapsto u\circ d_{X/S}$ est un isomorphisme de
$\Gamma(X,\mathcal{O}_X)$-modules
\[
  \operatorname{Hom}_{\mathcal{O}_X}(\Omega_{X/S}^1,\mathscr{F})
  \xrightarrow{\sim}
  \operatorname{Dér}_S(\mathcal{O}_X,\mathscr{F}).
  \tag{16.5.3.1}
  \label{IV.16.5.3.1}
\]
\end{enumerate}
\end{proposition}

L'assertion (i) a déjà été notée (16.3.6). D'autre part, il est immédiat (en
vertu de (0, 20.4.8)) que $u\mapsto u\circ d_{X/S}$ est injectif, en
considérant les restrictions à une fibre $\mathcal{O}_x$ des deux membres et
utilisant (16.4.16.2). Pour voir que l'homomorphisme (16.5.3.1) est surjectif,
considérons une $S$-dérivation $D:\mathcal{O}_X\to\mathscr{F}$; pour tout
ouvert affine $V=\operatorname{Spec}(B)$ de $X$, tel que $f(V)$ soit contenu
dans un ouvert affine $U=\operatorname{Spec}(A)$ de $S$,
$D_V:B\to\Gamma(V,\mathscr{F})$ est une $A$-dérivation, donc il existe un
\emph{unique} $B$-homomorphisme
$u_V:\Omega_{B/A}^1\to\Gamma(V,\mathscr{F})$ tel que
$D_V=u_V\circ d_{B/A}$ (0, 20.4.8); en outre l'unicité de $u_V$ montre
aussitôt que pour tout ouvert affine $W\subset V$, on a $u_W=u_V|W$, donc les
$u_V$ définissent un homomorphisme de $\mathcal{O}_X$-Modules
$u:\mathcal{O}_X\to\mathscr{F}$ répondant à la question.

\begin{env}[16.5.4]
\label{IV.16.5.4}
Avec les notations de (16.5.1), pour tout ouvert $U$ de $X$,
$\operatorname{Dér}_S(\mathcal{O}_U,\mathscr{F}|U)$ est un
$\Gamma(U,\mathcal{O}_X)$-module et il est clair que l'application
$U\mapsto\operatorname{Dér}_S(\mathcal{O}_U,\mathscr{F}|U)$ est un
préfaisceau; en fait, c'est même un \emph{faisceau} (donc un
$\mathcal{O}_X$-Module), en vertu de la caractérisation ponctuelle des
$S$-dérivations, vue en (16.5.1). Cet $\mathcal{O}_X$-Module se note
$\mathscr{D}\!{\rm ér}_S(\mathcal{O}_X,\mathscr{F})$ et est appelé
\emph{faisceau des $S$-dérivations de $\mathcal{O}_X$ dans $\mathscr{F}$}, et
ce qu'on vient de voir s'exprime encore par le corollaire suivant:
\end{env}

\begin{corollary}[16.5.5]
\label{IV.16.5.5}
Pour tout $\mathcal{O}_X$-Module $\mathscr{F}$, l'homomorphisme de
$\mathcal{O}_X$-Modules déduit de $u\mapsto u\circ d_{X/S}$
\[
  \mathscr{H}\!om_{\mathcal{O}_X}(\Omega_{X/S}^1,\mathscr{F})
  \longrightarrow
  \mathscr{D}\!{\rm ér}_S(\mathcal{O}_X,\mathscr{F})
  \tag{16.5.5.1}
  \label{IV.16.5.5.1}
\]
est bijectif.
\end{corollary}

\begin{corollary}[16.5.6]
\label{IV.16.5.6}
\begin{enumerate}
\item[(i)] Si le morphisme $f:X\to S$ est localement de présentation finie et
si $\mathscr{F}$ est un $\mathcal{O}_X$-Module quasi-cohérent, alors
$\mathscr{D}\!{\rm ér}_S(\mathcal{O}_X,\mathscr{F})$ est un
$\mathcal{O}_X$-Module quasi-cohérent.
\item[(ii)] Si de plus $S$ est localement noethérien et si $\mathscr{F}$ est
cohérent, alors $\mathscr{D}\!{\rm ér}_S(\mathcal{O}_X,\mathscr{F})$ est un
$\mathcal{O}_X$-Module cohérent.
\end{enumerate}
\end{corollary}

L'assertion (i) résulte de l'isomorphisme (16.5.5.1), de (16.4.22) et
(I, 1.3.12); l'assertion (ii) résulte de (0$_{\rm I}$, 5.3.5).

\begin{env}[16.5.7]
\label{IV.16.5.7}
On pose
\[
  \mathfrak{G}_{X/S}=
  \mathscr{H}\!om_{\mathcal{O}_X}(\Omega_{X/S}^1,\mathcal{O}_X)
  =\mathscr{D}\!{\rm ér}_S(\mathcal{O}_X,\mathcal{O}_X)
  \tag{16.5.7.1}
  \label{IV.16.5.7.1}
\]
et on dit que c'est le \emph{faisceau des $S$-dérivations de
$\mathcal{O}_X$}, ou encore le \emph{faisceau tangent de $X$ relativement à
$S$}: c'est donc le \emph{dual} du $\mathcal{O}_X$-Module
$\Omega_{X/S}^1$. Si $f$ est localement de présentation
\oldpage[IV-4]{29}
finie, $\mathfrak{G}_{X/S}$ est un $\mathcal{O}_X$-Module quasi-cohérent; si
de plus $S$ est localement noethérien, $\mathfrak{G}_{X/S}$ est cohérent
(16.5.6).
\end{env}

\begin{env}[16.5.8]
\label{IV.16.5.8}
Supposons plus particulièrement que $\Omega_{X/S}^1$ soit un
$\mathcal{O}_X$-Module \emph{localement libre} (de rang fini) (ce qui sera le
cas lorsque $f$ est \emph{lisse} (17.2.3)); alors $\mathfrak{G}_{X/S}$ est un
$\mathcal{O}_X$-Module localement libre et a même rang que $\Omega_{X/S}^1$
en chaque point. De façon plus précise, supposons que $\Omega_{X/S}^1$ soit de
rang $n$ en un point $x$; il y a alors $n$ sections $s_i$ ($1\leq i\leq n$)
de $\mathcal{O}_X$ au-dessus d'un voisinage affine $U$ de $x$ telles que les
images canoniques des $ds_i$ dans
$\Omega_{X/S}^1\otimes_{\mathcal{O}_X}k(x)$ forment une base de ce
$k(x)$-espace vectoriel; en vertu du lemme de Nakayama, les germes
$(ds_i)_x=d(s_i)_x$ des $ds_i$ au point $x$ forment une base du
$\mathcal{O}_x$-module $(\Omega_{X/S}^1)_x$, donc, en restreignant $U$, on
peut supposer que les $ds_i$ forment une \emph{base} du
$\Gamma(U,\mathcal{O}_X)$-module $\Gamma(U,\Omega_{X/S}^1)$. Alors le
$\Gamma(U,\mathcal{O}_X)$-module $\Gamma(U,\mathfrak{G}_{X/S})$ est dual du
précédent; on note $(D_i)_{1\leq i\leq n}$ ou
$\left(\frac{\partial}{\partial s_i}\right)_{1\leq i\leq n}$ la base duale de
$(ds_i)_{1\leq i\leq n}$, de sorte que, par (16.5.3), on a
\[
  D_is_j=\langle D_i,ds_j\rangle
  =\left\langle\frac{\partial}{\partial s_i},ds_j\right\rangle
  =\delta_{ij}\qquad\text{(indice de Kronecker).}
  \tag{16.5.8.1}
  \label{IV.16.5.8.1}
\]
Toute $\Gamma(S,\mathcal{O}_S)$-dérivation de la
$\Gamma(S,\mathcal{O}_S)$-algèbre $\Gamma(U,\mathcal{O}_X)$ s'écrit donc
d'une seule manière
\[
  D=\sum_{i=1}^n a_iD_i=\sum_{i=1}^n a_i\frac{\partial}{\partial s_i}
\]
où les $a_i$ ($1\leq i\leq n$) sont des sections de $\mathcal{O}_X$ au-dessus
de $U$. Pour toute section $g\in\Gamma(U,\mathcal{O}_X)$, si l'on pose
$dg=\sum_{i=1}^n c_i ds_i$, on a
$c_i=\langle D_i,dg\rangle=D_ig$ en vertu de (16.5.8.1), autrement dit
\[
  dg=\sum_{i=1}^n(D_ig)ds_i
  =\sum_{i=1}^n\frac{\partial g}{\partial s_i}ds_i.
  \tag{16.5.8.2}
  \label{IV.16.5.8.2}
\]
\end{env}

\begin{env}[16.5.9]
\label{IV.16.5.9}
Soient $D_1$, $D_2$ deux $S$-dérivations de $\mathcal{O}_X$. Pour tout ouvert
$U$ de $X$, si $D_1^U$, $D_2^U$ sont les dérivations correspondantes de
l'anneau $\Gamma(U,\mathcal{O}_X)$, le \emph{crochet}
\[
  [D_1^U,D_2^U]=D_1^U\circ D_2^U-D_2^U\circ D_1^U
\]
est aussi une dérivation de cet anneau, donc le
$\psi^*(\mathcal{O}_S)$-endomorphisme de $\mathcal{O}_X$
\[
  [D_1,D_2]=D_1\circ D_2-D_2\circ D_1
  \tag{16.5.9.1}
  \label{IV.16.5.9.1}
\]
est encore une $S$-dérivation; comme on vérifie aussitôt que ce crochet
vérifie l'identité de Jacobi, on voit qu'on a ainsi défini sur
$\operatorname{Dér}_S(\mathcal{O}_X,\mathcal{O}_X)$ une structure de
$\Gamma(S,\mathcal{O}_S)$-algèbre de Lie. Comme la définition de cette
structure commute à la restriction à un ouvert de $X$, on voit que
$\mathfrak{G}_{X/S}$ est canoniquement muni d'une structure de
$\psi^*(\mathcal{O}_S)$-Algèbre de Lie. On notera que l'application
$(D_1,D_2)\mapsto[D_1,D_2]$ n'est pas
$\Gamma(X,\mathcal{O}_X)$-bilinéaire.
\end{env}

\begin{env}[16.5.10]
\label{IV.16.5.10}
Pour tout changement de base $g:S'\to S$, si l'on pose $X'=X\times_S S'$, on
a vu (16.4.5) que l'on a un isomorphisme canonique
\[
  \Omega_{X/S}^1\otimes_S S'\xrightarrow{\sim}\Omega_{X'/S'}^1
  \tag{16.5.10.1}
  \label{IV.16.5.10.1}
\]
\oldpage[IV-4]{30}
d'où l'on déduit, en vertu de (16.5.10.1), un homomorphisme canonique
(Bourbaki, \emph{Alg.}, chap. II, 3\textsuperscript{e} éd., \S~5,
n\textsuperscript{o} 3)
\[
  \mathfrak{G}_{X/S}\otimes_{\mathcal{O}_S}\mathcal{O}_{S'}
  \longrightarrow\mathfrak{G}_{X'/S'}
  \tag{16.5.10.2}
  \label{IV.16.5.10.2}
\]
qui en général n'est ni injectif ni surjectif. Toutefois:
\end{env}

\begin{proposition}[16.5.11]
\label{IV.16.5.11}
\begin{enumerate}
\item[(i)] Si $g:S'\to S$ est un morphisme plat et si $f$ est localement de
type fini (resp.\ localement de présentation finie), l'homomorphisme
(16.5.10.2) est injectif (resp.\ bijectif).
\item[(ii)] Si $\Omega_{X/S}^1$ est un $\mathcal{O}_X$-Module localement libre
de type fini, l'homomorphisme (16.5.10.2) est bijectif.
\end{enumerate}
\end{proposition}

En effet, l'assertion (ii) résulte de Bourbaki, \emph{Alg.}, chap. II,
3\textsuperscript{e} éd., \S~5, n\textsuperscript{o} 3, prop. 7. L'assertion
(i) résulte de même de Bourbaki, \emph{Alg. comm.}, chap. I, \S~2,
n\textsuperscript{o} 10, prop. 11 et du fait que si $f$ est localement de type
fini (resp.\ localement de présentation finie), $\Omega_{X/S}^1$ est un
$\mathcal{O}_X$-Module de type fini (resp.\ de présentation finie) ((16.3.9)
et (16.4.22)).

\begin{env}[16.5.12]
\label{IV.16.5.12}
Puisque $\Omega_{X/S}^1$ est un $\mathcal{O}_X$-Module quasi-cohérent, on peut
considérer le fibré vectoriel sur $X$ défini par $\Omega_{X/S}^1$
(II, 1.7.8)
\[
  T_{X/S}=\mathbf{V}(\Omega_{X/S}^1)
  \tag{16.5.12.1}
  \label{IV.16.5.12.1}
\]
qu'on appelle \emph{fibré tangent de $X$ relativement à $S$}. On a donc une
bijection canonique (II, 1.7.9)
\[
  \Gamma(T_{X/S}/S)\xrightarrow{\sim}
  \operatorname{Hom}_{\mathcal{O}_X}(\Omega_{X/S}^1,\mathcal{O}_X)
  =\Gamma(X,\mathfrak{G}_{X/S})
\]
par définition de $\mathfrak{G}_{X/S}$, et on peut dans cet isomorphisme
remplacer $X$ par un ouvert quelconque $U$ de $X$; on peut donc dire que le
\emph{faisceau tangent} de $X$ relativement à $S$ est isomorphe au
\emph{faisceau des germes de $S$-sections} du fibré tangent de $X$
relativement à $S$. Si $f:X\to Y$ est un $S$-morphisme, on a vu (16.4.19)
qu'on a un homomorphisme canonique
$f_{X/Y/S}:f^*(\Omega_{Y/S}^1)\to\Omega_{X/S}^1$, ce qui donne, compte tenu de
ce que
\[
  \mathbf{V}(f^*(\Omega_{Y/S}^1))
  =\mathbf{V}(\Omega_{Y/S}^1)\times_Y X\qquad\text{(II, 1.7.11),}
\]
un $X$-morphisme
$T_{X/S}(f):T_{X/S}\to T_{Y/S}\times_Y X$. Si $g:Y\to Z$ est un second
$S$-morphisme, on a
$T_{X/S}(g\circ f)=(T_{Y/S}(g)\times1_X)\circ T_{X/S}(f)$ (0, 20.5.4.1).

Il résulte de (16.5.10.1) et de (II, 1.7.11) que pour tout changement de base
$g:S'\to S$, on a un isomorphisme canonique
\[
  T_{X'/S'}\xrightarrow{\sim}T_{X/S}\times_S S'
  =T_{X/S}\times_X X'.
  \tag{16.5.12.2}
  \label{IV.16.5.12.2}
\]
\end{env}

\begin{env}[16.5.13]
\label{IV.16.5.13}
Pour tout point $x\in X$, on appelle \emph{espace tangent à $X$ au point $x$}
(relativement à $S$) l'ensemble des points de la fibre
$T_{X/S}\times_X\operatorname{Spec}(k(x))$ \emph{rationnels sur $k(x)$}, donc
l'ensemble
\[
  T_{X/S}(x)=
  \operatorname{Hom}_{k(x)}
  (\Omega_{X/S}^1\otimes_{\mathcal{O}_x}k(x),k(x))
  \tag{16.5.13.1}
  \label{IV.16.5.13.1}
\]
qui est le \emph{dual} du $k(x)$-espace vectoriel
$\Omega_{\mathcal{O}_x/\mathcal{O}_s}^1/
\mathfrak{m}_x\cdot\Omega_{\mathcal{O}_x/\mathcal{O}_s}^1$. Lorsque
$\Omega_{X/S}^1$ est un $\mathcal{O}_X$-Module de \emph{type fini},
$T_{X/S}(x)$ est donc un espace vectoriel de rang fini sur $k(x)$, et pour
tout
\oldpage[IV-4]{31}
changement de base $g:S'\to S$, et tout point
$x'\in X'=X\times_S S'$ au-dessus de $x$, on a un isomorphisme canonique
\[
  T_{X'/S'}(x')\xrightarrow{\sim}
  T_{X/S}(x)\otimes_{k(x)}k(x').
  \tag{16.5.13.2}
  \label{IV.16.5.13.2}
\]

Si $x$ est \emph{rationnel sur $k(s)$}, où $s=f(x)$ (de sorte que
$k(s)\to k(x)$ est un isomorphisme), il résulte de (16.4.13) que l'on a un
isomorphisme canonique
\[
  T_{X/S}(x)=T_{X_s/k(s)}(x)
  =\operatorname{Hom}_{k(s)}
  (\mathfrak{m}'_x/\mathfrak{m}'_x{}^2,k(s))
  \tag{16.5.13.3}
  \label{IV.16.5.13.3}
\]
où $\mathfrak{m}'_x$ est l'idéal maximal de
$\mathcal{O}_{X_s,x}=\mathcal{O}_{X,x}/
\mathfrak{m}_s\mathcal{O}_{X,x}$. Dans le cas où $S$ est le spectre d'un corps
$k$, on retrouve ainsi la définition de Zariski de l'espace tangent en un
point $x\in X$ \emph{rationnel sur $k$}, comme dual de
$\mathfrak{m}_x/\mathfrak{m}_x^2$.

Soit $Y$ un second $S$-préschéma et soit $g:Y\to X$ un $S$-morphisme; on a
alors un homomorphisme canonique de $\mathcal{O}_Y$-Modules (16.4.19)
\[
  g_{Y/X/S}:g^*(\Omega_{X/S}^1)\longrightarrow\Omega_{Y/S}^1.
  \tag{16.5.13.4}
  \label{IV.16.5.13.4}
\]
Notons maintenant que si $y\in Y$ et $x=g(y)$, on a
\[
  g^*(\Omega_{X/S}^1)\otimes_{\mathcal{O}_Y}k(y)
  =(\Omega_{X/S}^1\otimes_{\mathcal{O}_X}k(x))
  \otimes_{k(x)}k(y)
\]
et par suite, si $\Omega_{X/S}^1$ est un $\mathcal{O}_X$-Module
\emph{de type fini}, on peut identifier
\[
  \operatorname{Hom}_{k(y)}
  (g^*(\Omega_{X/S}^1)\otimes_{\mathcal{O}_Y}k(y),k(y))
\]
à $T_{X/S}(x)\otimes_{k(x)}k(y)$. On déduit donc alors de l'homomorphisme
(16.5.13.4) un homomorphisme de \emph{$k(y)$-espaces vectoriels}
\[
  T_y(g):T_{Y/S}(y)\longrightarrow
  T_{X/S}(x)\otimes_{k(x)}k(y)
  \tag{16.5.13.5}
  \label{IV.16.5.13.5}
\]
dit \emph{application linéaire tangente à $g$ au point $y$}. Lorsque $y$ est
\emph{rationnel sur $k(s)$}, on peut identifier $k(s)$, $k(y)$ et $k(x)$ et
$T_y(g)$ est alors un homomorphisme de $k(s)$-espaces vectoriels
$T_{Y/S}(y)\to T_{X/S}(x)$; on notera d'ailleurs que dans ce cas,
$g^*(\Omega_{X/S}^1)\otimes_{\mathcal{O}_Y}k(y)$ s'identifie à
$\Omega_{X/S}^1\otimes_{\mathcal{O}_X}k(x)$, et l'homomorphisme précédent est
donc défini sans aucune condition de finitude sur $\Omega_{X/S}^1$ et n'est
autre que l'homomorphisme $T_{Y/S}(g)$ (16.5.12) restreint à la fibre au point
$y$ de $T_{Y/S}$.
\end{env}

\begin{env}[16.5.14]
\label{IV.16.5.14}
L'interprétation des dérivations d'une $A$-algèbre $B$ dans un $B$-module
$L$, donnée dans (0, 20.1.1), se traduit dans le langage des préschémas de la
façon suivante.

Considérons deux morphismes de préschémas $f:X\to S$, $g:Y\to S$, et un
sous-préschéma fermé $Y_0$ de $Y$ défini par un Idéal de carré nul
$\mathscr{J}$ de $\mathcal{O}_Y$ (de sorte que $Y$ et $Y_0$ ont
\emph{même espace topologique sous-jacent}). Supposons donné un $S$-morphisme
$u_0:Y_0\to X$, de sorte qu'on a un diagramme commutatif
\[
  \xymatrix{
    X\ar[d]_-f & Y_0\ar[l]_-{u_0}\ar[d]^j\\
    S & Y\ar@{-->}[ul]_-u\ar[l]^-g
  }
  \tag{16.5.14.1}
  \label{IV.16.5.14.1}
\]
\oldpage[IV-4]{32}
et proposons-nous de chercher s'il existe des \emph{$S$-morphismes}
$u:Y\to X$ tels que $u_0=u\circ j$ (autrement dit, s'il est possible de
compléter le diagramme précédent par la flèche pointillée $u$, de façon à le
laisser \emph{commutatif}).

Pour cela, considérons un ouvert affine $U=\operatorname{Spec}(C)$ de $Y$;
son image réciproque $j^{-1}(U)$ est l'ouvert affine
$U_0=\operatorname{Spec}(C/\mathfrak{L})$, où
$\mathfrak{L}=\Gamma(U,\mathscr{J})$, idéal \emph{de carré nul} dans $C$;
nous supposerons $U$ assez petit pour que $u_0(U_0)$ soit contenu dans un
ouvert affine $V=\operatorname{Spec}(B)$ de $X$, et
$g(U)=f(u_0(U_0))$ contenu dans un ouvert affine
$W=\operatorname{Spec}(A)$ de $S$, de sorte que $B$ et $C$ sont des
$A$-algèbres et que $u_0|U_0$ correspond à un $A$-homomorphisme $\psi$ de
$B$ dans $C/\mathfrak{L}$; soit $P(U_0)$ l'ensemble des restrictions $u|U$
des homomorphismes cherchés, qui correspondent canoniquement aux
$A$-homomorphismes d'algèbres $\varphi:B\to C$ tels que le \emph{composé}
\[
  B\xrightarrow{\varphi}C\longrightarrow C/\mathfrak{L}
\]
\emph{soit égal à $\psi$}. On sait donc (0, 20.1.1) que l'ensemble de ces
homomorphismes est vide ou de la forme
$\varphi_1+\operatorname{D\acute{e}r}_A(B,\mathfrak{L})$; lorsque $P(U_0)$ n'est pas
vide, le groupe additif $\operatorname{D\acute{e}r}_A(B,\mathfrak{L})$ opère par
addition dans $P(U_0)$, qui est donc un \emph{espace affine} pour le groupe
additif $\operatorname{D\acute{e}r}_A(B,\mathfrak{L})$ (ou encore un
\emph{espace homogène principal} (ou \emph{torseur}) \emph{sous}
$\operatorname{D\acute{e}r}_A(B,\mathfrak{L})$).

Remarquons maintenant que, puisque $\mathfrak{L}$ est muni d'une structure de
$B$-module au moyen de $\psi$, on a un \emph{isomorphisme}
$v\mapsto v\circ d_{B/A}$ de
$\operatorname{Hom}_B(\Omega_{B/A}^1,\mathfrak{L})$ sur
$\operatorname{D\acute{e}r}_A(B,\mathfrak{L})$ (0, 20.4.8). D'ailleurs, comme
$\mathfrak{L}$ est de carré nul, donc un $(C/\mathfrak{L})$-module, tout
$B$-homomorphisme $v:\Omega_{B/A}^1\to\mathfrak{L}$ peut être considéré
comme un $(C/\mathfrak{L})$-homomorphisme
$\Omega_{B/A}^1\otimes_B(C/\mathfrak{L})\to\mathfrak{L}$. Comme
$\mathscr{J}$ est de carré nul, il peut être considéré comme un
$\mathcal{O}_{Y_0}$-Module quasi-cohérent; introduisons le
$\mathcal{O}_{Y_0}$-Module
\[
  \mathscr{G}=\mathscr{H}om(u_0^*(\Omega_{X/S}^1),\mathscr{J});
  \tag{16.5.14.2}
  \label{IV.16.5.14.2}
\]
il résulte alors du fait que
$\Omega_{B/A}^1=\Gamma(V,\Omega_{X/S}^1)$ (16.3.7) que l'on peut écrire
$\operatorname{D\acute{e}r}_A(B,\mathfrak{L})=\Gamma(U_0,\mathscr{G})$.

Comme $P(U_0)$ est défini comme ensemble de $S$-morphismes $U\to X$, il est
clair que $U_0\mapsto P(U_0)$ est un \emph{faisceau d'ensembles}
$\mathscr{P}$ sur $Y_0$. Utilisons ce fait pour prouver que l'application
\[
  h:\Gamma(U_0,\mathscr{G})\times P(U_0)\longrightarrow P(U_0)
\]
définissant la structure de torseur sur $P(U_0)$ est aussi indépendante du
choix de $V$ et $W$, et en outre que, si $U'\subset U$ est un second ouvert
affine de $Y$, $U'_0$ son image réciproque dans $Y_0$, le diagramme
\[
  \xymatrix{
    \Gamma(U_0,\mathscr{G})\times P(U_0)\ar[d]\ar[r]^h
      & P(U_0)\ar[d]\\
    \Gamma(U'_0,\mathscr{G})\times P(U'_0)\ar[r]_{h'}
      & P(U'_0)
  }
  \tag{16.5.14.3}
  \label{IV.16.5.14.3}
\]
est commutatif (les flèches verticales étant les opérateurs de restriction).
En vertu de la remarque précédente, on est ramené à prouver la commutativité
du diagramme précédent lorsque $h$ est défini comme ci-dessus à partir des
ouverts affines $V$, $W$ et $h'$ à partir d'ouverts
\oldpage[IV-4]{33}
affines $V'\subset V$ et $W'\subset W$. Mais en vertu de la description
précédente de $h$, cela résulte de la commutativité du diagramme (0, 20.5.4.2).

Les applications
$\Gamma(U_0,\mathscr{G})\times P(U_0)\to P(U_0)$ définissent donc un
homomorphisme de \emph{faisceaux d'ensembles}
$m:\mathscr{G}\times\mathscr{P}\to\mathscr{P}$ tel que, pour tout ouvert
$U_0$ pour lequel $\Gamma(U_0,\mathscr{P})\neq\varnothing$,
\[
  m_{U_0}:\Gamma(U_0,\mathscr{G})\times\Gamma(U_0,\mathscr{P})
  \longrightarrow\Gamma(U_0,\mathscr{P})
\]
soit une loi externe définissant sur $\Gamma(U_0,\mathscr{P})$ une structure
de torseur sous le groupe $\Gamma(U_0,\mathscr{G})$.
\end{env}

\begin{env}[16.5.15]
\label{IV.16.5.15}
D'une façon générale, lorsqu'on se donne sur un espace topologique $Z$ un
faisceau d'ensembles $\mathscr{P}$, un faisceau de groupes $\mathscr{G}$ (non
nécessairement commutatif) et un homomorphisme de faisceaux d'ensembles
$m:\mathscr{G}\times\mathscr{P}\to\mathscr{P}$ tel que, pour tout ouvert
$U\subset Z$ tel que $\Gamma(U,\mathscr{P})\neq\varnothing$,
\[
  m_U:\Gamma(U,\mathscr{G})\times\Gamma(U,\mathscr{P})
  \longrightarrow\Gamma(U,\mathscr{P})
\]
fasse de $\Gamma(U,\mathscr{P})$ un \emph{torseur} sous le groupe
$\Gamma(U,\mathscr{G})$, on dit que $\mathscr{P}$ est un
\emph{pseudo-torseur} (ou \emph{faisceau formellement principal homogène})
sous le faisceau de groupes $\mathscr{G}$. On dit que $\mathscr{P}$ est un
\emph{torseur} (ou \emph{faisceau principal homogène}) sous $\mathscr{G}$ si
de plus $\Gamma(U,\mathscr{P})\neq\varnothing$ pour tout ouvert
$U\neq\varnothing$ d'une base convenable de la topologie de $Z$.

Pour la théorie générale des torseurs, nous renvoyons à [42]; nous nous
bornerons ici à rappeler la correspondance canonique entre les classes
d'isomorphie de ces torseurs (pour un $\mathscr{G}$ donné) et les éléments de
l'ensemble de cohomologie $H^1(Z,\mathscr{G})$. Considérons en effet un
torseur $\mathscr{P}$ sous $\mathscr{G}$ et un recouvrement ouvert
$(U_\lambda)$ de $Z$ tel que
$\Gamma(U_\lambda,\mathscr{P})\neq\varnothing$ pour tout $\lambda$;
désignons par $p_\lambda$ un élément de
$\Gamma(U_\lambda,\mathscr{P})$. Pour tout couple d'indices $\lambda$, $\mu$
tel que $U_\lambda\cap U_\mu\neq\varnothing$, il existe alors un élément et
un seul $\gamma_{\lambda\mu}$ de
$\Gamma(U_\lambda\cap U_\mu,\mathscr{G})$ tel que
\[
  \gamma_{\lambda\mu}\cdot
  (p_\mu|U_\lambda\cap U_\mu)=p_\lambda|U_\lambda\cap U_\mu;
\]
en outre, si $\lambda$, $\mu$, $\nu$ sont trois indices tels que
$U_\lambda\cap U_\mu\cap U_\nu\neq\varnothing$, les restrictions
$\gamma'_{\lambda\mu}$, $\gamma'_{\mu\nu}$, $\gamma'_{\lambda\nu}$ de
$\gamma_{\lambda\mu}$, $\gamma_{\mu\nu}$, $\gamma_{\lambda\nu}$ à
$U_\lambda\cap U_\mu\cap U_\nu$ vérifient la condition
$\gamma'_{\lambda\nu}=\gamma'_{\lambda\mu}\gamma'_{\mu\nu}$; autrement dit
$(\lambda,\mu)\mapsto\gamma_{\lambda\mu}$ est un \emph{1-cocycle} du
recouvrement $(U_\lambda)$ à valeurs dans $\mathscr{G}$. Si, pour tout
$\lambda$, $p'_\lambda$ est un second élément de
$\Gamma(U_\lambda,\mathscr{P})$, il existe un élément et un seul
$\beta_\lambda\in\Gamma(U_\lambda,\mathscr{G})$ tel que
$p'_\lambda=\beta_\lambda\cdot p_\lambda$ et le 1-cocycle
$(\gamma'_{\lambda\mu})$ correspondant à la famille $(p'_\lambda)$ est donné
par
\[
  \gamma'_{\lambda\mu}=\beta_\lambda\gamma_{\lambda\mu}\beta_\mu^{-1},
\]
autrement dit, est \emph{cohomologue} à $(\gamma_{\lambda\mu})$.
Inversement, la donnée d'un 1-cocycle $(\gamma_{\lambda\mu})$ définit, pour
tout couple $(\lambda,\mu)$, un automorphisme $\theta_{\lambda\mu}$ du
faisceau \emph{d'ensembles}
$\mathscr{G}|U_\lambda\cap U_\mu$, savoir la translation à droite par
$\gamma_{\lambda\mu}$, et le fait qu'il s'agisse d'un cocycle montre que l'on
peut \emph{recoller} les faisceaux d'ensembles $\mathscr{G}|U_\lambda$ à
l'aide des automorphismes $\theta_{\lambda\mu}$ (0, 3.3.1); on obtient ainsi
de façon évidente un torseur sous $\mathscr{G}$, soit $\mathscr{P}$, et si
l'on prend pour $p_\lambda$ la section unité au-dessus de $U_\lambda$, le
1-cocycle correspondant n'est autre que le 1-cocycle donné
$(\gamma_{\lambda\mu})$; en outre, si l'on remplace
$(\gamma_{\lambda\mu})$ par un 1-cocycle
$\gamma'_{\lambda\mu}=\beta_\lambda\gamma_{\lambda\mu}\beta_\mu^{-1}$ qui
lui est cohomologue, on vérifie aussitôt que le torseur obtenu est isomorphe
à $\mathscr{P}$.

En particulier, si $(\gamma_{\lambda\mu})$ est un \emph{1-cobord}, autrement
dit de la forme
$\gamma_{\lambda\mu}=\beta_\lambda\beta_\mu^{-1}$, le torseur obtenu
$\mathscr{P}$ est \emph{isomorphe à $\mathscr{G}$} (considéré comme torseur
sous lui-même pour les translations à gauche); on dit dans ce cas que
$\mathscr{P}$ est \emph{trivial}, et la réciproque est évidente.

Plus particulièrement, il résulte de (III, 1.3.1) que l'on a:
\end{env}

\begin{proposition}[16.5.16]
\label{IV.16.5.16}
Soient $Z$ un schéma affine, $\mathscr{G}$ un $\mathcal{O}_Z$-Module
quasi-cohérent; alors tout torseur sous $\mathscr{G}$ est trivial.
\end{proposition}

\oldpage[IV-4]{34}
Revenant au problème considéré en (16.5.13), on obtient donc:

\begin{proposition}[16.5.17]
\label{IV.16.5.17}
Soient $X$, $Y$ deux $S$-préschémas, $Y_0$ un sous-préschéma fermé de $Y$
défini par un Idéal quasi-cohérent $\mathscr{J}$ de $\mathcal{O}_Y$ tel que
$\mathscr{J}^2=0$, $j:Y_0\to Y$ l'injection canonique. Soit
$u_0:Y_0\to X$ un $S$-morphisme, et $\mathscr{P}$ le faisceau d'ensembles sur
$Y$ tel que, pour tout ouvert $U$ de $Y$, $\Gamma(U,\mathscr{P})$ soit
l'ensemble des $S$-morphismes $u:U\to X$ tels que
$u_0|U_0=u\circ(j|U_0)$, où $U_0=j^{-1}(U)$. Alors il existe sur
$\mathscr{P}$ une structure de pseudo-torseur sous le
$\mathcal{O}_{Y_0}$-Module
\[
  \mathscr{G}=\mathscr{H}om_{\mathcal{O}_{Y_0}}
  (u_0^*(\Omega_{X/S}^1),\mathscr{J}).
\]
\end{proposition}

En particulier:

\begin{corollary}[16.5.18]
\label{IV.16.5.18}
Avec les notations de (16.5.16), supposons $Y$ affine et $\Omega_{X/S}^1$ de
présentation finie; s'il existe un recouvrement ouvert $(U_\alpha)$ de $Y$,
et, pour chaque indice $\alpha$, un $S$-morphisme $v_\alpha:U_\alpha\to X$
tel que, si $U_\alpha^0=j^{-1}(U_\alpha)$, on ait
$v_\alpha\circ(j|U_\alpha^0)=u_0|U_\alpha^0$, alors il existe un
$S$-morphisme $u:Y\to X$ tel que $u\circ j=u_0$.
\end{corollary}

En effet, $\mathscr{G}$ est alors un $\mathcal{O}_{Y_0}$-Module
quasi-cohérent (I, 1.3.12); en vertu de (16.5.16) et du fait que $Y_0$ est
alors affine, le faisceau $\mathscr{P}$, qui est par hypothèse un torseur sous
$\mathscr{G}$, et non seulement un pseudo-torseur, est \emph{trivial}; mais si
$w$ est un isomorphisme de $\mathscr{G}$ sur $\mathscr{P}$ (en tant que
torseurs sous $\mathscr{G}$), l'image par $w$ de la section 0 de
$\mathscr{G}$ est le $S$-morphisme $u$ cherché.

\subsection{Faisceaux de $p$-différentielles et différentielle extérieure.}
\label{IV.16.6}

\begin{env}[16.6.1]
\label{IV.16.6.1}
Soit $f:X\to S$ un morphisme de préschémas. On appelle \emph{faisceau des
$p$-différentielles de $X$ relativement à $S$} ($p$ entier) la
\emph{puissance extérieure $p$-ème} (0, 4.1.5) du $\mathcal{O}_X$-Module
$\Omega_{X/S}^1$, notée
\[
  \Omega_{X/S}^p=\bigwedge^p(\Omega_{X/S}^1).
  \tag{16.6.1.1}
  \label{IV.16.6.1.1}
\]
On a donc $\Omega_{X/S}^0=\mathcal{O}_X$, et $\Omega_{X/S}^p=0$ pour $p<0$;
les $\Omega_{X/S}^p$ sont les composants homogènes de l'Algèbre extérieure
sur $\Omega_{X/S}^1$
\[
  \Omega_{X/S}^{\bullet}=\bigwedge(\Omega_{X/S}^1)
  =\bigoplus_{p\in\mathbf{Z}}\bigwedge^p(\Omega_{X/S}^1),
  \tag{16.6.1.2}
  \label{IV.16.6.1.2}
\]
qui est donc une $\mathcal{O}_X$-Algèbre graduée quasi-cohérente,
anticommutative et dont les éléments de degré 1 sont de carré nul. Pour tout
ouvert affine $U$ de $X$, on a
\[
  \Gamma(U,\Omega_{X/S}^{\bullet})
  =\bigwedge(\Gamma(U,\Omega_{X/S}^1)),
\]
où $\Gamma(U,\Omega_{X/S}^1)$ est considéré comme
$\Gamma(U,\mathcal{O}_X)$-module.

Lorsque $S=\operatorname{Spec}(A)$ et $X=\operatorname{Spec}(B)$ sont
affines, $B$ étant donc une $A$-algèbre, on a (0, 4.1.5)
$\Omega_{X/S}^p=(\Omega_{B/A}^p)^{\sim}$, en posant
$\Omega_{B/A}^p=\bigwedge^p\Omega_{B/A}^1$.
\end{env}

\begin{theorem}[16.6.2]
\label{IV.16.6.2}
Il existe un endomorphisme et un seul $d$ du faisceau de groupes additifs
$\Omega_{X/S}^{\bullet}$, ayant les propriétés suivantes:
\begin{enumerate}
\item[(i)] $d\circ d=0$.
\item[(ii)] Pour tout ouvert $U$ de $X$ et toute section
$f\in\Gamma(U,\mathcal{O}_X)$, on a $df=d_{X/S}f$.
\oldpage[IV-4]{35}
\item[(iii)] Pour tout ouvert $U$ de $X$, tout couple d'entiers $p$, $q$ et
tout couple de sections
$\omega'_p\in\Gamma(U,\Omega_{X/S}^p)$,
$\omega''_q\in\Gamma(U,\Omega_{X/S}^q)$, on a
\[
  d(\omega'_p\wedge\omega''_q)
  =(d\omega'_p)\wedge\omega''_q
  +(-1)^p\omega'_p\wedge d\omega''_q.
  \tag{16.6.2.1}
  \label{IV.16.6.2.1}
\]
\end{enumerate}
En outre, $d$ est un endomorphisme de $\psi^*(\mathcal{O}_X)$-Modules
gradués, de degré $+1$.
\end{theorem}

Supposons prouvée l'existence de l'endomorphisme $d$. Pour tout ouvert affine
$U$ de $X$, toute section de $\Omega_{X/S}^p$ au-dessus de $U$ est (en vertu
de (i)) combinaison linéaire d'un nombre fini d'éléments de la forme
$g(df_1\wedge df_2\wedge\cdots\wedge df_p)$, où $g$ et les $f_i$ sont des
sections de $\mathcal{O}_X$ au-dessus de $U$ (0, 20.4.7). Les conditions (i)
et (iii) montrent alors, par récurrence sur $p$, que l'on a nécessairement
\[
  d(g(df_1\wedge df_2\wedge\cdots\wedge df_p))
  =dg\wedge df_1\wedge df_2\wedge\cdots\wedge df_p.
  \tag{16.6.2.2}
  \label{IV.16.6.2.2}
\]
Cela prouve donc l'\emph{unicité} de $d$ et la dernière assertion du théorème.
En vertu de cette propriété d'unicité, pour démontrer l'existence de $d$, on
peut se borner au cas où $S=\operatorname{Spec}(A)$ et
$X=\operatorname{Spec}(B)$ sont affines. Or (Bourbaki, \emph{Alg.}, chap.
III, 3\textsuperscript{e} éd., \S~10) pour définir une $A$-antidérivation
$D$ de degré $+1$ d'une algèbre extérieure $\bigwedge(M)$ (où $M$ est un
$B$-module et $B$ une $A$-algèbre), cette antidérivation prenant ses valeurs
dans une $A$-algèbre \emph{graduée anticommutative}
$C=\bigoplus_{n=0}^{\infty}C_n$, dont les éléments de degré 1 sont de carré
nul, il suffit de se donner arbitrairement une $A$-dérivation $D_0$ de $B$
dans $C_1$ et un $A$-homomorphisme $D_1$ de $M$ dans $C_2$; il existe alors
une $A$-antidérivation $D$ et une seule de $\bigwedge(M)$ dans $C$
coïncidant avec $D_0$ dans $B$ et avec $D_1$ dans $M$.

Dans le cas actuel, $D_0$ est nécessairement égal à $d_{B/A}$ en vertu de
(ii); tout revient à voir, compte tenu de (16.6.2.2), qu'il y a un
$A$-homomorphisme $u$ de $\Omega_{B/A}^1$ dans $\Omega_{B/A}^2$ tel que
l'on ait
\[
  u(g.df)=dg\wedge df
  \tag{16.6.2.3}
  \label{IV.16.6.2.3}
\]
quels que soient $f$, $g$ dans $A$; il suffira pour cela de montrer qu'il
existe un $A$-homomorphisme
$v:B\otimes_A\Omega_{B/A}^1\to\Omega_{B/A}^2$ tel que
\[
  v(g.\omega)=dg\wedge\omega
  \tag{16.6.2.4}
  \label{IV.16.6.2.4}
\]
pour $g\in B$ et $\omega\in\Omega_{B/A}^1$. Finalement, comme
$\Omega_{B/A}^1=\mathfrak{I}/\mathfrak{I}^2$ (où
$\mathfrak{I}=\mathfrak{I}_{B/A}$ est le noyau de l'homomorphisme canonique
$B\otimes_A B\to B$) et que $\Omega_{B/A}^1$ est engendré par les éléments
de la forme $g.df$, il suffit de définir un $A$-homomorphisme
$w:B\otimes_A(B\otimes_A B)\to\Omega_{B/A}^2$ tel que
\[
  w(g'\otimes g\otimes f)=dg'\wedge(g.df)
  \tag{16.6.2.5}
  \label{IV.16.6.2.5}
\]
et tel que $w$ soit nul dans l'image de $B\otimes_A\mathfrak{I}^2$. Or,
comme le second membre de (16.6.2.5) est $A$-trilinéaire en $g'$, $g$, $f$,
l'existence de $w$ vérifiant (16.6.2.5) est immédiate. Comme, d'autre part,
$\mathfrak{I}$ est engendré par les éléments $1\otimes x-x\otimes1$
($x\in B$), on est ramené à vérifier que lorsque
$z=(1\otimes x-x\otimes1)(1\otimes y-y\otimes1)$, on a
$w(g'\otimes z)=0$. Or, comme
$z=1\otimes(xy)+(xy)\otimes1-x\otimes y-y\otimes x$, la formule (16.6.2.4)
montre qu'il
\oldpage[IV-4]{36}
suffit de voir que l'on a $d(xy)-x.dy-y.dx=0$, ce qui exprime que $d$ est une
dérivation.

Il reste à prouver que $d$ vérifie la condition (i). Or, le carré d'une
antidérivation est une dérivation (Bourbaki, \emph{loc. cit.}), et comme
$\Omega_{B/A}^{\bullet}$ est engendré par $\Omega_{B/A}^1$ comme
$B$-algèbre, il suffit de vérifier que $d(dz)=0$ pour $z\in B$ et
$z\in\Omega_{B/A}^1$; dans le premier cas, cela résulte de la formule
(16.6.2.3) avec $g=1$; dans le second, on peut se borner au cas où $z=g.df$
avec $f$, $g$ dans $B$, et alors, en vertu de (16.6.2.1) et (16.6.2.3), on a
\[
  d(d(g.df))=d(dg\wedge df)
  =(d(dg))\wedge(df)-(dg)\wedge(d(df))=0.
\]
\hfill C.Q.F.D.

\begin{definition}[16.6.3]
\label{IV.16.6.3}
L'antidérivation $d$ définie dans (16.6.2) (aussi notée $d_{X/S}$) est
appelée \emph{différentielle extérieure sur $X$} (relativement à $S$).
\end{definition}

\begin{proposition}[16.6.4]
\label{IV.16.6.4}
Pour tout changement de base $g:S'\to S$, si l'on pose
$X'=X\times_S S'$, l'homomorphisme canonique
\[
  \Omega_{X/S}^{\bullet}\otimes_S S'
  \longrightarrow\Omega_{X'/S'}^{\bullet}
  \tag{16.6.4.1}
  \label{IV.16.6.4.1}
\]
déduit de l'isomorphisme (16.5.9.1) est bijectif. En outre, si $s$ est une
section de $\Omega_{X/S}^{\bullet}$ au-dessus d'un ouvert $U$ de $X$,
$s\otimes1$ son image réciproque, section de $\Omega_{X'/S'}^{\bullet}$
au-dessus de l'image réciproque $U'$ de $U$ dans $X'$, on a
$d_{X'/S'}(s\otimes1)=d_{X/S}(s)\otimes1$.
\end{proposition}

La première assertion est immédiate, la formation de l'algèbre extérieure
d'un module permutant avec toute extension de l'anneau des scalaires. Pour
prouver la seconde, on peut, en vertu de (16.6.2.2), se borner au cas où
$s\in\Gamma(U,\mathcal{O}_X)$, et dans ce cas l'assertion a déjà été prouvée
(16.4.3.7).

\begin{env}[16.6.5]
\label{IV.16.6.5}
Supposons que $\Omega_{X/S}^1$ soit un $\mathcal{O}_X$-Module localement libre
de rang $n$ en un point $x$, de sorte qu'il y a $n$ sections
$s_i\in\Gamma(U,\mathcal{O}_X)$ telles que les $ds_i$ forment une base du
$\Gamma(U,\mathcal{O}_X)$-module $\Gamma(U,\Omega_{X/S}^1)$ (16.5.8).
Alors, pour tout entier $p\geq1$, les $p$-différentielles
$ds_{i_1}\wedge ds_{i_2}\wedge\cdots\wedge ds_{i_p}$ (pour
$i_1<i_2<\cdots<i_p$, éléments de $[1,n]$) forment une base de
$\binom{n}{p}$ éléments de $\Gamma(U,\Omega_{X/S}^p)$ sur
$\Gamma(U,\mathcal{O}_X)$. En outre la formule (16.6.2.2) montre que pour
toute section $g\in\Gamma(U,\mathcal{O}_X)$, on a
\[
  d(g.ds_{i_1}\wedge ds_{i_2}\wedge\cdots\wedge ds_{i_p})
  =\sum_k(-1)^r\frac{\partial g}{\partial s_k}
  ds_{i_1}\wedge\cdots\wedge ds_{i_r}\wedge ds_k\wedge
  ds_{i_{r+1}}\wedge\cdots\wedge ds_{i_p}
  \tag{16.6.5.1}
  \label{IV.16.6.5.1}
\]
où, au second membre, $k$ parcourt l'ensemble des $n-p$ indices distincts des
$i_h$, $i_r$ étant le plus grand indice $<k$.

On notera que la relation $d(dg)=0$ pour toute section
$g\in\Gamma(U,\mathcal{O}_X)$ s'exprime sous la forme
\[
  D_i(D_jg)=D_j(D_ig)\qquad\text{pour $i\neq j$};
\]
autrement dit, les dérivations $D_i$ définies dans (16.5.7) sont deux à deux
permutables.
\end{env}

\subsection{Les $\mathscr{P}_{X/S}^n(\mathscr{F})$.}
\label{IV.16.7}

\begin{env}[16.7.1]
\label{IV.16.7.1}
Soient $f:X\to S$ un morphisme de préschémas, $\mathscr{F}$ un
$\mathcal{O}_X$-Module. Désignons par $X_{\Delta_f}^{(n)}$ le \emph{$n$-ème
voisinage infinitésimal de $X$} pour le morphisme diagonal
\oldpage[IV-4]{37}
$\Delta_f:X\to X\times_S X$, par
$h_n:X_{\Delta_f}^{(n)}\to X\times_S X$ le morphisme canonique (16.1.2), et
considérons les deux morphismes composés
\[
  p_1^{(n)}:X_{\Delta_f}^{(n)}\xrightarrow{h_n}X\times_S X
  \xrightarrow{p_1}X,qquad
  p_2^{(n)}:X_{\Delta_f}^{(n)}\xrightarrow{h_n}X\times_S X
  \xrightarrow{p_2}X
\]
de sorte que, par définition, $p_1^{(n)}$ correspond à l'homomorphisme de
faisceaux d'anneaux $\mathcal{O}_X\to\mathscr{P}_{X/S}^n$ que nous avons
choisi pour définir la structure de $\mathcal{O}_X$-Algèbre sur
$\mathscr{P}_{X/S}^n$ (16.3.5), et $p_2^{(n)}$ à l'homomorphisme de faisceaux
d'anneaux $d_{X/S}^n:\mathcal{O}_X\to\mathscr{P}_{X/S}^n$ (16.3.6). Comme
$X_{\Delta_f}^{(n)}$ et $X$ ont même espace sous-jacent, on peut écrire
\[
  \mathscr{P}_{X/S}^n
  =(p_1^{(n)})_*((p_2^{(n)})^*(\mathcal{O}_X)).
  \tag{16.7.1.1}
  \label{IV.16.7.1.1}
\]
Plus généralement, nous poserons
\[
  \mathscr{P}_{X/S}^n(\mathscr{F})
  =(p_1^{(n)})_*((p_2^{(n)})^*(\mathscr{F}))
  \tag{16.7.1.2}
  \label{IV.16.7.1.2}
\]
de sorte que
$\mathscr{P}_{X/S}^n=\mathscr{P}_{X/S}^n(\mathcal{O}_X)$; par définition,
$\mathscr{P}_{X/S}^n(\mathscr{F})$ est donc un $\mathcal{O}_X$-Module.
\end{env}

\begin{env}[16.7.2]
\label{IV.16.7.2}
Si l'on revient aux définitions des images réciproques de Modules sur les
espaces annelés (0, 4.3.1) et que l'on tienne compte de ce que
$X_{\Delta_f}^{(n)}$ et $X$ ont même espace sous-jacent, on voit que l'on peut
aussi écrire la définition (16.7.1.2) sous la forme
\[
  \mathscr{P}_{X/S}^n(\mathscr{F})
  =\mathscr{P}_{X/S}^n\otimes_{\mathcal{O}_X}\mathscr{F}
  \tag{16.7.2.1}
  \label{IV.16.7.2.1}
\]
mais où il faut prendre garde que, dans l'interprétation du signe $\otimes$,
$\mathscr{P}_{X/S}^n$ est muni de sa structure de $\mathcal{O}_X$-Module
définie par l'homomorphisme de faisceaux d'anneaux
$d_{X/S}^n:\mathcal{O}_X\to\mathscr{P}_{X/S}^n$. Il résulte aussitôt de cette
formule (ou directement de (16.7.1.2)) que
$\mathscr{P}_{X/S}^n(\mathscr{F})$ est canoniquement muni d'une structure de
$\mathscr{P}_{X/S}^n$-Module.
\end{env}

\begin{proposition}[16.7.3]
\label{IV.16.7.3}
\begin{enumerate}
\item[(i)] Le foncteur
$\mathscr{F}\mapsto\mathscr{P}_{X/S}^n(\mathscr{F})$ de la catégorie des
$\mathcal{O}_X$-Modules dans celle des $\mathscr{P}_{X/S}^n$-Modules est
exact à droite, et commute aux limites inductives quelconques; il est exact
lorsque $\mathscr{P}_{X/S}^n$ est un $\mathcal{O}_X$-Module plat.
\item[(ii)] Si $\mathscr{F}$ est un $\mathcal{O}_X$-Module quasi-cohérent
(resp. de type fini, resp. de présentation finie),
$\mathscr{P}_{X/S}^n(\mathscr{F})$ est un $\mathscr{P}_{X/S}^n$-Module
quasi-cohérent (resp. de type fini, resp. de présentation finie).
\end{enumerate}
\end{proposition}

Les assertions de (i) résultent aussitôt de la formule (16.7.2.1) et de la
considération de la symétrie de $\mathscr{P}_{X/S}^n$ (16.3.4). Les assertions
de (ii) découlent de l'exactitude à droite du foncteur
$\mathscr{F}\mapsto\mathscr{P}_{X/S}^n(\mathscr{F})$.

\begin{env}[16.7.4]
\label{IV.16.7.4}
Les deux structures de $\mathcal{O}_X$-Module de $\mathscr{P}_{X/S}^n$
définissent sur $\mathscr{P}_{X/S}^n(\mathscr{F})$ deux structures de
$\mathcal{O}_X$-Module, d'ailleurs permutables, donc une structure de
$\mathcal{O}_X$-Bimodule. Il est commode de noter \emph{à gauche} celle de
ces structures provenant de l'homomorphisme structural
$\mathcal{O}_X\to\mathscr{P}_{X/S}^n$ (choisi dans (16.3.5)) et
\emph{à droite} celle provenant de l'homomorphisme
$d_{X/S}^n:\mathcal{O}_X\to\mathscr{P}_{X/S}^n$. Autrement dit, pour tout
ouvert $U$ de $X$, et tout triplet d'éléments
$a\in\Gamma(U,\mathcal{O}_X)$, $b\in\Gamma(U,\mathscr{P}_{X/S}^n)$,
$t\in\Gamma(U,\mathscr{F})$, on a par définition
\[
  a(b\otimes t)=(ab)\otimes t,qquad
  (b\otimes t)a=(b.d^na)\otimes t=b\otimes(at)
  =(d^na).(b\otimes t).
  \tag{16.7.4.1}
  \label{IV.16.7.4.1}
\]
La structure de $\mathcal{O}_X$-Module provenant de la définition (16.7.1.2)
est donc, avec ces conventions, la structure de $\mathcal{O}_X$-Module à
gauche.

\oldpage[IV-4]{38}
Si $\mathscr{F}$ est un $\mathcal{O}_X$-Module quasi-cohérent, il en est de
même de $\mathscr{P}_{X/S}^n(\mathscr{F})$ pour l'une ou l'autre de ses
structures de $\mathcal{O}_X$-Module. Si de plus $\mathscr{F}$ est de type
fini (resp. de présentation finie) et $f:X\to S$ localement de type fini
(resp. localement de présentation finie),
$\mathscr{P}_{X/S}^n(\mathscr{F})$ est (pour l'une ou l'autre de ses
structures de $\mathcal{O}_X$-Module) de type fini (resp. de présentation
finie), comme il résulte de (16.3.9) et de (16.4.22).
\end{env}

\begin{env}[16.7.5]
\label{IV.16.7.5}
La définition (16.7.2.1) entraîne l'existence d'un homomorphisme de faisceaux
de groupes commutatifs
\[
  d_{X/S,\mathscr{F}}^n:\mathscr{F}
  \longrightarrow\mathscr{P}_{X/S}^n(\mathscr{F})
  \qquad\text{(aussi noté $d_{X/S}^n$)}
  \tag{16.7.5.1}
  \label{IV.16.7.5.1}
\]
tel qu'avec les notations de (16.7.4), on ait
\[
  d_{X/S,\mathscr{F}}^n(t)=1\otimes t
  \tag{16.7.5.2}
  \label{IV.16.7.5.2}
\]
et par suite, en vertu de (16.7.4.1)
\[
  d_{X/S,\mathscr{F}}^n(at)
  =(1\otimes t)a=(d_{X/S,\mathscr{F}}^n(t)).a
  \tag{16.7.5.3}
  \label{IV.16.7.5.3}
\]
\[
  d_{X/S,\mathscr{F}}^n(at)
  =(d_{X/S}^n(a)).(1\otimes t)
  =(d_{X/S}^n(a))(d_{X/S,\mathscr{F}}^n(t)).
  \tag{16.7.5.4}
  \label{IV.16.7.5.4}
\]
Il est donc $\mathcal{O}_X$-linéaire pour la structure de
$\mathcal{O}_X$-Module à droite sur $\mathscr{P}_{X/S}^n(\mathscr{F})$, et
\emph{semi-linéaire} (relativement à l'automorphisme $\sigma$ (16.3.4)) pour
la structure de $\mathcal{O}_X$-Module à gauche.
\end{env}

\begin{proposition}[16.7.6]
\label{IV.16.7.6}
Le $\mathcal{O}_X$-Module à gauche $\mathscr{P}_{X/S}^n(\mathscr{F})$ est
engendré par l'image de $\mathscr{F}$ par l'homomorphisme canonique
$d_{X/S,\mathscr{F}}^n$.
\end{proposition}

Cela résulte aussitôt de (16.7.5.3) et du cas particulier correspondant à
$\mathscr{F}=\mathcal{O}_X$ (16.3.8).

\begin{env}[16.7.7]
\label{IV.16.7.7}
Les homomorphismes canoniques de faisceaux d'anneaux
\[
  \varphi_{nm}:\mathscr{P}_{X/S}^m\longrightarrow\mathscr{P}_{X/S}^n
\]
pour $n\leq m$ (16.1.2) définissent, en vertu de (16.7.2.1) des
homomorphismes canoniques
\[
  \mathscr{P}_{X/S}^m(\mathscr{F})
  \longrightarrow\mathscr{P}_{X/S}^n(\mathscr{F})qquad(n\leq m)
\]
qui sont des homomorphismes de $\mathcal{O}_X$-Bimodules en vertu de (16.1.6)
et (16.7.4.1); en outre on a des diagrammes commutatifs
\[
  \xymatrix{
    \mathscr{P}_{X/S}^m(\mathscr{F})\ar[rr]
      &&\mathscr{P}_{X/S}^n(\mathscr{F})\\
    &\mathscr{F}\ar[ul]^-{d_{X/S,\mathscr{F}}^m}
      \ar[ur]_-{d_{X/S,\mathscr{F}}^n}
  }
\]
On a donc ainsi un système projectif de $\mathcal{O}_X$-Bimodules
$(\mathscr{P}_{X/S}^n(\mathscr{F}))$, et l'on pose
\[
  \mathscr{P}_{X/S}^{\infty}(\mathscr{F})
  =\varprojlim\mathscr{P}_{X/S}^n(\mathscr{F}).
  \tag{16.7.7.1}
  \label{IV.16.7.7.1}
\]
En outre, ce qui précède montre que les homomorphismes (16.7.5.1) forment un
système projectif d'homomorphismes, et définissent donc un homomorphisme
canonique
\[
  d_{X/S,\mathscr{F}}^{\infty}:\mathscr{F}
  \longrightarrow\mathscr{P}_{X/S}^{\infty}(\mathscr{F}).
  \tag{16.7.7.2}
  \label{IV.16.7.7.2}
\]
\end{env}

\oldpage[IV-4]{39}

\begin{env}[16.7.8]
\label{IV.16.7.8}
Soient $\mathscr{F}$, $\mathscr{G}$ deux $\mathcal{O}_X$-Modules; il résulte
aussitôt de la définition (16.7.2.1) que l'on a un isomorphisme canonique de
$\mathscr{P}_{X/S}^n$-Modules
\[
  \mathscr{P}_{X/S}^n(\mathscr{F}\otimes_{\mathcal{O}_X}\mathscr{G})
  \xrightarrow{\sim}
  \mathscr{P}_{X/S}^n(\mathscr{F})
  \otimes_{\mathscr{P}_{X/S}^n}\mathscr{P}_{X/S}^n(\mathscr{G})
  \tag{16.7.8.1}
  \label{IV.16.7.8.1}
\]
(Bourbaki, \emph{Alg.}, chap.~II, 3\textsuperscript{e} éd., §~5,
n\textsuperscript{o}~1, prop.~3).

On en conclut en particulier (ou l'on voit directement sur la définition
(16.7.2.1)) que si $\mathscr{F}$ est muni d'une structure de
$\mathcal{O}_X$-Algèbre (non nécessairement associative),
$\mathscr{P}_{X/S}^n(\mathscr{F})$ est canoniquement muni d'une structure de
$\mathscr{P}_{X/S}^n$-Algèbre; cette dernière est associative (resp.
commutative, resp. unitaire, resp. une Algèbre de Lie) lorsqu'il en est ainsi
de $\mathscr{F}$. En outre les homomorphismes canoniques
$\mathscr{P}_{X/S}^m(\mathscr{F})\to\mathscr{P}_{X/S}^n(\mathscr{F})$ pour
$n\leq m$ (16.7.7) sont alors des di-homomorphismes d'Algèbres; de même
(16.7.5.1) est alors un homomorphisme de $\mathcal{O}_X$-Algèbres lorsque
$\mathscr{P}_{X/S}^n(\mathscr{F})$ est muni de sa structure de
$\mathcal{O}_X$-Algèbre provenant de sa structure de $\mathcal{O}_X$-Module
\emph{à droite}.

Avec les mêmes notations, on a également un homomorphisme canonique de
$\mathscr{P}_{X/S}^n$-Modules
\[
  \mathscr{P}_{X/S}^n
  (\mathscr{H}\!om_{\mathcal{O}_X}(\mathscr{F},\mathscr{G}))
  \longrightarrow
  \mathscr{H}\!om_{\mathscr{P}_{X/S}^n}
  (\mathscr{P}_{X/S}^n(\mathscr{F}),\mathscr{P}_{X/S}^n(\mathscr{G}))
  \tag{16.7.8.2}
  \label{IV.16.7.8.2}
\]
(Bourbaki, \emph{Alg.}, 3\textsuperscript{e} éd., §~5,
n\textsuperscript{o}~3), qui est bijectif lorsque $\mathscr{P}_{X/S}^n$ est
un $\mathcal{O}_X$-Module \emph{localement libre de type fini}
(\emph{loc. cit.}, prop.~7).
\end{env}

\begin{env}[16.7.9]
\label{IV.16.7.9}
Supposons qu'on soit dans la situation décrite dans (16.4.1); alors, de
l'homomorphisme canonique $P^n(u)$ (16.4.3.3), on déduit aussitôt un
homomorphisme canonique de $\mathcal{O}_{X'}$-Bimodules
\[
  u^*(\mathscr{P}_{X/S}^n(\mathscr{F}))
  \longrightarrow
  \mathscr{P}_{X'/S'}^n(u^*(\mathscr{F})).
  \tag{16.7.9.1}
  \label{IV.16.7.9.1}
\]
Nous laissons au lecteur le soin d'étendre à cet homomorphisme les propriétés
vues dans (16.4) pour le cas $\mathscr{F}=\mathcal{O}_X$.
\end{env}

\begin{remark}[16.7.10]
\label{IV.16.7.10}
La définition de $\mathscr{P}_{X/S}^n(\mathscr{F})$ sous la forme (16.7.1.2)
garde un sens lorsque $\mathscr{F}$ est un faisceau d'ensembles quelconque
(l'image réciproque d'un faisceau d'ensembles par $p_2^{(n)}$ étant définie
dans (0\textsubscript{I}, 3.7.1)); une variante de cette définition permet de
définir le «schéma des jets» (relativement à $S$) d'un $X$-préschéma
quelconque.
\end{remark}

\subsection{Opérateurs différentiels.}
\footnote{Pour un formalisme plus général, voir l'Exposé VII de [42]
(dû à P. Gabriel).}
\label{IV.16.8}

\begin{definition}[16.8.1]
\label{IV.16.8.1}
Soient $f=(\psi,\theta):X\to S$ un morphisme de préschémas,
$\mathscr{F}$, $\mathscr{G}$ deux $\mathcal{O}_X$-Modules, $n$ un entier
$\geq0$. On dit qu'un homomorphisme de faisceaux de groupes additifs
$D:\mathscr{F}\to\mathscr{G}$ est un \emph{opérateur différentiel d'ordre
$\leq n$} (relativement à $S$) s'il existe un homomorphisme de
$\mathcal{O}_X$-Modules
$u:\mathscr{P}_{X/S}^n(\mathscr{F})\to\mathscr{G}$ (où
$\mathscr{P}_{X/S}^n(\mathscr{F})$ est muni de sa structure de
$\mathcal{O}_X$-Module à gauche (16.7.4)) tel que l'on ait
$D=u\circ d_{X/S,\mathscr{F}}^n$.
\end{definition}

Il est clair, en vertu de l'existence des homomorphismes canoniques
\[
  \mathscr{P}_{X/S}^m(\mathscr{F})
  \longrightarrow\mathscr{P}_{X/S}^n(\mathscr{F})
\]
\oldpage[IV-4]{40}
pour $n\leq m$ (16.7.7) qu'un opérateur différentiel d'ordre $\leq n$ est
aussi un opérateur différentiel d'ordre $\leq m$ pour tout $m\geq n$. Si
$D:\mathscr{F}\to\mathscr{G}$ est un opérateur différentiel d'ordre
$\leq n$, alors, pour tout ouvert $U$ de $X$,
$D|U:\mathscr{F}|U\to\mathscr{G}|U$ est aussi un opérateur différentiel
d'ordre $\leq n$.

On dit qu'un homomorphisme $D:\mathscr{F}\to\mathscr{G}$ des faisceaux de
groupes additifs sous-jacents à $\mathscr{F}$ et à $\mathscr{G}$ est un
\emph{opérateur différentiel} (relativement à $S$) si, pour tout $x\in X$, il
existe un voisinage ouvert $U$ de $x$ et un entier $n\geq0$ tels que
$D|U:\mathscr{F}|U\to\mathscr{G}|U$ soit un opérateur différentiel d'ordre
$\leq n$. L'\emph{ordre} d'un opérateur différentiel
$D:\mathscr{F}\to\mathscr{G}$ est la borne inférieure des entiers $n$ tels
que $D$ soit un opérateur différentiel d'ordre $\leq n$ (et donc $+\infty$
s'il n'y a pas de tels entiers); cet ordre est toujours fini si $X$ est
\emph{quasi-compact}. Les opérateurs différentiels d'ordre $0$ ne sont autres
que les homomorphismes de $\mathcal{O}_X$-Modules
$\mathscr{F}\to\mathscr{G}$; on convient que tout opérateur différentiel
d'ordre $<0$ est \emph{nul}. Pour $n\geq0$, un opérateur différentiel n'est
pas en général un homomorphisme de $\mathcal{O}_X$-Modules mais est toujours
un homomorphisme de $\psi^*(\mathcal{O}_S)$-Modules.

Lorsque $\mathscr{F}=\mathcal{O}_X$, un opérateur différentiel d'ordre
$\leq1$ de $\mathcal{O}_X$ dans $\mathscr{G}$ s'écrit d'une seule manière
sous la forme $v+D$, où $v:\mathcal{O}_X\to\mathscr{G}$ est un
$\mathcal{O}_X$-homomorphisme, et $D$ une $S$-\emph{dérivation} (16.5.1) de
$\mathcal{O}_X$ dans $\mathscr{G}$: cela résulte de la structure de
$P_{B/A}^1$ (0, 20.4.8).

\begin{env}[16.8.2]
\label{IV.16.8.2}
Pour décrire de façon plus explicite un opérateur différentiel d'ordre
$\leq n$, $D:\mathscr{F}\to\mathscr{G}$, il suffit, pour tout ouvert affine
$U$ de $X$, dont l'image dans $S$ est contenue dans un ouvert affine $V$, de
caractériser l'homomorphisme
$D=D_U:\Gamma(U,\mathscr{F})\to\Gamma(U,\mathscr{G})$. Si l'on pose
$\Gamma(V,\mathcal{O}_S)=A$, $\Gamma(U,\mathcal{O}_X)=B$, de sorte que $B$
est une $A$-algèbre, on a
\[
  \Gamma(U,\mathscr{P}_{X/S}^n)
  =(B\otimes_A B)/\mathfrak{I}^{n+1},
\]
où l'on a posé pour abréger $\mathfrak{I}=\mathfrak{I}_{B/A}$. Posons en
outre $M=\Gamma(U,\mathscr{F})$, $N=\Gamma(U,\mathscr{G})$; alors la
définition de $D$ signifie que pour chaque couple $(U,V)$ vérifiant les
conditions précédentes, le $A$-homomorphisme $D:M\to N$ se factorise en
\[
  M\longrightarrow
  ((B\otimes_A B)/\mathfrak{I}^{n+1})\otimes_B M
  \xrightarrow{v}N
\]
où la première flèche est l'homomorphisme canonique $t\mapsto1\otimes t$, et
$v$ est un $B$-homomorphisme, la structure de $B$-module de
$((B\otimes_A B)/\mathfrak{I}^{n+1})\otimes_B M$ provenant du premier
facteur $B$ (alors qu'on rappelle que dans la formation du produit tensoriel
sur $B$, la structure de $B$-module de
$(B\otimes_A B)/\mathfrak{I}^{n+1}$ provient du second facteur $B$). Notons
maintenant que le $B$-module
$((B\otimes_A B)/\mathfrak{I}^{n+1})\otimes_B M$ est isomorphe à
$(B\otimes_A M)/\mathfrak{I}^{n+1}(B\otimes_A M)$, où $B\otimes_A M$ est
considéré comme $(B\otimes_A B)$-module et sa structure de $B$-module
provient de l'homomorphisme $b\mapsto b\otimes1$ de $B$ dans $B\otimes_A B$.
Soit alors $D'$ le $B$-homomorphisme de $B\otimes_A M$ dans $N$ tel que
$D'(b\otimes t)=bD(t)$; la condition de factorisation sur $D$ s'exprime encore
en disant que $D'$ doit \emph{être nul dans le $B$-module}
$\mathfrak{I}^{n+1}(B\otimes_A M)$.
\end{env}

\begin{env}[16.8.3]
\label{IV.16.8.3}
Il est clair que l'ensemble des opérateurs différentiels d'ordre $\leq n$ de
$\mathscr{F}$ dans $\mathscr{G}$ est un groupe additif, noté
$\operatorname{Diff}_{X/S}^n(\mathscr{F},\mathscr{G})$; lorsque
$\mathscr{F}=\mathscr{G}=\mathcal{O}_X$, on écrit aussi
$\operatorname{Diff}_{X/S}^n$ au lieu de
$\operatorname{Diff}_{X/S}^n(\mathcal{O}_X,\mathcal{O}_X)$.

On a vu (16.8.1) que l'on a pour deux ouverts $U\supset V$ de $X$, un
homomorphisme canonique de restriction
\[
  \operatorname{Diff}_{U/S}^n(\mathscr{F}|U,\mathscr{G}|U)
  \longrightarrow
  \operatorname{Diff}_{V/S}^n(\mathscr{F}|V,\mathscr{G}|V)
\]
\oldpage[IV-4]{41}
donc $U\mapsto\operatorname{Diff}_{U/S}^n(\mathscr{F}|U,\mathscr{G}|U)$ est
un préfaisceau de groupes additifs; en fait c'est même un \emph{faisceau}, car
pour $U$ ouvert variable dans $X$, les homomorphismes
$u\mapsto u\circ d_{U/S,\mathscr{F}|U}^n$ sont des \emph{isomorphismes} de
groupes additifs
\[
  \mathscr{H}\!om_{\mathcal{O}_U}
  (\mathscr{P}_{U/S}^n(\mathscr{F}|U),\mathscr{G}|U)
  \xrightarrow{\sim}
  \operatorname{Diff}_{U/S}^n(\mathscr{F}|U,\mathscr{G}|U),
  \tag{16.8.3.1}
  \label{IV.16.8.3.1}
\]
en vertu du fait que l'image de $\mathscr{F}$ par
$d_{X/S,\mathscr{F}}^n$ engendre $\mathscr{P}_{X/S}^n(\mathscr{F})$
(16.7.6). On note ce faisceau
$\mathscr{D}\!iff_{X/S}^n(\mathscr{F},\mathscr{G})$, et on a donc:
\end{env}

\begin{proposition}[16.8.4]
\label{IV.16.8.4}
Les isomorphismes (16.8.3.1) définissent un isomorphisme de faisceaux de
groupes additifs
\[
  \mathscr{H}\!om_{\mathcal{O}_X}
  (\mathscr{P}_{X/S}^n(\mathscr{F}),\mathscr{G})
  \xrightarrow{\sim}
  \mathscr{D}\!iff_{X/S}^n(\mathscr{F},\mathscr{G}).
  \tag{16.8.4.1}
  \label{IV.16.8.4.1}
\]
\end{proposition}

Lorsque $\mathscr{F}=\mathscr{G}=\mathcal{O}_X$, on écrit aussi
$\mathscr{D}\!iff_{X/S}^n$ au lieu de
$\mathscr{D}\!iff_{X/S}^n(\mathcal{O}_X,\mathcal{O}_X)$; il résulte de
(16.8.4) que $\mathscr{D}\!iff_{X/S}^n$ s'identifie canoniquement au
\emph{dual} du $\mathcal{O}_X$-Module $\mathscr{P}_{X/S}^n$; aussi écrit-on
$\langle t,D\rangle$ au lieu de $u(t)$ si $t$ est une section de
$\mathscr{P}_{X/S}^n$ au-dessus d'un ouvert et si $u$ est l'homomorphisme de
$\mathscr{P}_{X/S}^n$ dans $\mathcal{O}_X$ correspondant à $D$.

\begin{env}[16.8.5]
\label{IV.16.8.5}
Puisque $\mathscr{P}_{X/S}^n(\mathscr{F})$ est muni d'une structure de
$\mathcal{O}_X$-Bimodule (16.7.4), on en déduit canoniquement une structure de
$\mathcal{O}_X$-Bimodule sur
$\mathscr{H}\!om_{\mathcal{O}_X}
(\mathscr{P}_{X/S}^n(\mathscr{F}),\mathscr{G})$, donc aussi sur
$\mathscr{D}\!iff_{X/S}^n(\mathscr{F},\mathscr{G})$ en vertu de (16.8.4.1).
De façon précise, à la structure de $\mathcal{O}_X$-Module à gauche de
$\mathscr{P}_{X/S}^n(\mathscr{F})$ correspond, en vertu de la définition
(16.8.1), la structure de $\mathcal{O}_X$-Module à gauche sur
$\mathscr{D}\!iff_{X/S}^n(\mathscr{F},\mathscr{G})$ explicitée comme suit:
pour tout ouvert $U$ de $X$, toute section
$a\in\Gamma(U,\mathcal{O}_X)$ et tout opérateur différentiel
$D:\mathscr{F}|U\to\mathscr{G}|U$, $aD$ est l'opérateur différentiel qui, à
toute section $t\in\Gamma(U,\mathscr{F})$, fait correspondre la section
\[
  (aD)(t)=a(D(t))
  \tag{16.8.5.1}
  \label{IV.16.8.5.1}
\]
de $\Gamma(U,\mathscr{G})$. De même, à la structure de
$\mathcal{O}_X$-Module à droite de $\mathscr{P}_{X/S}^n(\mathscr{F})$
correspond la structure de $\mathcal{O}_X$-Module à droite sur
$\mathscr{D}\!iff_{X/S}^n(\mathscr{F},\mathscr{G})$ explicitée comme suit:
avec les mêmes notations que ci-dessus, $Da$ est l'opérateur différentiel qui,
à $t\in\Gamma(U,\mathscr{F})$, fait correspondre la section
\[
  (Da)(t)=D(at).
  \tag{16.8.5.2}
  \label{IV.16.8.5.2}
\]
\end{env}

\begin{proposition}[16.8.6]
\label{IV.16.8.6}
Si $f:X\to S$ est un morphisme localement de présentation finie,
$\mathscr{F}$ un $\mathcal{O}_X$-Module quasi-cohérent de présentation finie
et $\mathscr{G}$ un $\mathcal{O}_X$-Module quasi-cohérent, alors
$\mathscr{D}\!iff_{X/S}^n(\mathscr{F},\mathscr{G})$ est un
$\mathcal{O}_X$-Module quasi-cohérent pour l'une ou l'autre structure définies
dans (16.8.5).
\end{proposition}

La proposition résulte de ce que, sous les hypothèses faites,
$\mathscr{P}_{X/S}^n(\mathscr{F})$ est un $\mathcal{O}_X$-Module
quasi-cohérent de présentation finie (16.7.4), et de (I, 1.3.12).

\begin{env}[16.8.7]
\label{IV.16.8.7}
L'ensemble des opérateurs différentiels de $\mathscr{F}$ dans $\mathscr{G}$
(d'ordre non précisé (16.8.1)) se note
$\operatorname{Diff}_{X/S}(\mathscr{F},\mathscr{G})$; on voit encore comme
dans (16.8.3) que
$U\mapsto\operatorname{Diff}_{U/S}(\mathscr{F}|U,\mathscr{G}|U)$ est un
faisceau de groupes additifs, que nous désignerons par
$\mathscr{D}\!iff_{X/S}(\mathscr{F},\mathscr{G})$. Il est immédiat que
$\mathscr{D}\!iff_{X/S}(\mathscr{F},\mathscr{G})$ est réunion de la famille
filtrante croissante de ses sous-faisceaux
$\mathscr{D}\!iff_{X/S}^n(\mathscr{F},\mathscr{G})$; si $X$ est
quasi-compact, $\operatorname{Diff}_{X/S}(\mathscr{F},\mathscr{G})$
\oldpage[IV-4]{42}
est de même réunion de ses sous-groupes
$\operatorname{Diff}_{X/S}^n(\mathscr{F},\mathscr{G})$ (16.8.1). Les
structures de $\mathcal{O}_X$-Bimodule sur les
$\mathscr{D}\!iff_{X/S}^n(\mathscr{F},\mathscr{G})$ définissent donc une
structure de $\mathcal{O}_X$-Bimodule sur
$\mathscr{D}\!iff_{X/S}(\mathscr{F},\mathscr{G})$, explicitée encore par
(16.8.5.1) et (16.8.5.2).

Notons que, pour $n\leq m$, on a un diagramme commutatif
\[
  \xymatrix{
    \mathscr{H}\!om_{\mathcal{O}_X}
      (\mathscr{P}_{X/S}^n(\mathscr{F}),\mathscr{G})
      \ar[r]^-{\sim}\ar[d]
      &\mathscr{D}\!iff_{X/S}^n(\mathscr{F},\mathscr{G})\ar[d]\\
    \mathscr{H}\!om_{\mathcal{O}_X}
      (\mathscr{P}_{X/S}^m(\mathscr{F}),\mathscr{G})
      \ar[r]^-{\sim}
      &\mathscr{D}\!iff_{X/S}^m(\mathscr{F},\mathscr{G})
  }
  \tag{16.8.7.1}
  \label{IV.16.8.7.1}
\]
où les flèches horizontales sont les isomorphismes (16.8.4.2) et la flèche
verticale de gauche provient de l'homomorphisme canonique
$\mathscr{P}_{X/S}^m(\mathscr{F})\to
\mathscr{P}_{X/S}^n(\mathscr{F})$ (16.7.7). Pour tout ouvert $U$ de $X$,
munissons alors
\[
  \Gamma(U,\mathscr{P}_{X/S}^{\infty}(\mathscr{F}))
  =\varprojlim\Gamma(U,\mathscr{P}_{X/S}^n(\mathscr{F}))
\]
de la topologie limite projective des topologies discrètes sur les
$\Gamma(U,\mathscr{P}_{X/S}^n(\mathscr{F}))$, ce qui définit
$\Gamma(U,\mathscr{P}_{X/S}^{\infty}(\mathscr{F}))$ comme un
$\Gamma(U,\mathcal{O}_X)$-bimodule topologique, de sorte que
$\mathscr{P}_{X/S}^{\infty}(\mathscr{F})$ apparaît comme un faisceau à valeurs
dans la catégorie des groupes commutatifs topologiques
(0\textsubscript{I}, 3.2.6). Alors (G, II, 1.11) la limite du système inductif
de faisceaux de groupes commutatifs
\[
  (\mathscr{H}\!om_{\mathcal{O}_X}
  (\mathscr{P}_{X/S}^n(\mathscr{F}),\mathscr{G}))
\]
n'est autre que le faisceau des germes d'homomorphismes \emph{continus} de
$\mathscr{P}_{X/S}^{\infty}(\mathscr{F})$ dans $\mathscr{G}$ (ce dernier
étant muni de la topologie discrète): les homomorphismes continus de
$\Gamma(U,\mathscr{P}_{X/S}^{\infty}(\mathscr{F}))$ dans le groupe discret
$\Gamma(U,\mathscr{G})$ correspondent en effet biunivoquement aux systèmes
inductifs d'homomorphismes de groupes
$\Gamma(U,\mathscr{P}_{X/S}^n(\mathscr{F}))\to\Gamma(U,\mathscr{G})$. On peut
donc encore exprimer (16.8.4) en disant qu'on a un isomorphisme canonique
\[
  \mathscr{H}\!om.\mathrm{cont}_{\mathcal{O}_X}
  (\mathscr{P}_{X/S}^{\infty}(\mathscr{F}),\mathscr{G})
  \xrightarrow{\sim}
  \mathscr{D}\!iff_{X/S}(\mathscr{F},\mathscr{G})
\]
où le premier membre désigne le faisceau des germes d'homomorphismes continus
de $\mathscr{P}_{X/S}^{\infty}(\mathscr{F})$ dans $\mathscr{G}$.
\end{env}

\begin{proposition}[16.8.8]
\label{IV.16.8.8}
Soient $\mathscr{F}$, $\mathscr{G}$ deux $\mathcal{O}_X$-Modules,
$D:\mathscr{F}\to\mathscr{G}$ un homomorphisme de
$\psi^*(\mathcal{O}_S)$-Modules, $n$ un entier $\geq0$. Les conditions
suivantes sont équivalentes:
\begin{enumerate}
  \item[(a)] $D$ est un opérateur différentiel d'ordre $\leq n$.
  \item[(b)] Pour toute section $a$ de $\mathcal{O}_X$ au-dessus d'un ouvert
  $U$, l'homomorphisme $D_a:\mathscr{F}|U\to\mathscr{G}|U$ tel que, pour
  toute section $t$ de $\mathscr{F}$ au-dessus d'un ouvert $V\subset U$, on
  ait
  \[
    D_a(t)=D(at)-aD(t)
    \tag{16.8.8.1}
    \label{IV.16.8.8.1}
  \]
  est un opérateur différentiel d'ordre $\leq n-1$.
  \item[(c)] Pour tout ouvert $U$ de $X$, toute famille
  $(a_i)_{1\leq i\leq n+1}$ de $n+1$ sections de $\mathcal{O}_X$ au-dessus
  de $U$ et toute section $t$ de $\mathscr{F}$ au-dessus de $U$, on a
  l'identité
  \[
    \sum_{H\subset I_{n+1}}(-1)^{\operatorname{Card}(H)}
    \left(\prod_{i\in H}a_i\right)
    D\!\left(\left(\prod_{i\notin H}a_i\right)t\right)=0
    \tag{16.8.8.2}
    \label{IV.16.8.8.2}
  \]
  (où $I_{n+1}$ est l'intervalle $1\leq i\leq n+1$ de $\mathbb{N}$).
\end{enumerate}
\oldpage[IV-4]{43}
\end{proposition}

Prouvons d'abord l'équivalence de a) et c). Par définition, pour prouver que
$D$ est un opérateur différentiel d'ordre $\leq n$, il suffit de montrer qu'il
en est ainsi pour la restriction
$D|U:\mathscr{F}|U\to\mathscr{G}|U$ à tout ouvert affine $U$ de $X$, et
d'autre part la propriété c) est valable pour tout ouvert $U$ de $X$ si elle
l'est pour tout ouvert affine. On peut donc se borner au cas où
$S=\operatorname{Spec}(A)$ et $X=\operatorname{Spec}(B)$ sont affines. En
vertu de (16.8.2) (dont on garde les notations), la condition a) signifie
alors que le $A$-homomorphisme $D':B\otimes_A M\to N$ tel que
$D'(b\otimes t)=bD(t)$ s'annule dans
$\mathfrak{I}^{n+1}(B\otimes_A M)$, ce qui, en vertu de (0, 20.4.4), équivaut
à dire que $D'$ s'annule pour tous les éléments de la forme
\[
  \left(\prod_{i=1}^{n+1}(a_i\otimes1-1\otimes a_i)\right).(1\otimes t)
\]
où $a_i\in B$ et $t\in M$. Or, cet élément s'écrit
\[
  \sum_{H\subset I_{n+1}}
  \left(\prod_{i\in H}a_i\right)\otimes
  \left(\left(\prod_{i\notin H}a_i\right)t\right),
\]
et la valeur de $D'$ pour cet élément n'est autre que le premier membre de
(16.8.8.2), ce qui prouve l'équivalence de a) et c).

Prouvons maintenant l'équivalence de b) et c). Raisonnons par récurrence sur
$n$, l'assertion étant triviale pour $n=0$. Écrivant $a_{n+1}$ au lieu de $a$
dans la condition b), on voit, en vertu de l'hypothèse de récurrence, que la
condition b) signifie que pour toute famille $(a_i)_{1\leq i\leq n}$ de $n$
sections de $\mathcal{O}_X$ au-dessus de $U$ et toute section $t$ de
$\mathscr{F}$ au-dessus de $U$, on a
\[
  \sum_{H'\subset I_n}(-1)^{\operatorname{Card}(H')}
  \left(\prod_{i\in H'}a_i\right)
  D_{a_{n+1}}\!\left(\left(\prod_{i\notin H'}a_i\right)t\right)=0.
\]
Mais si l'on remplace dans cette relation $D_{a_{n+1}}$ par sa définition
(16.8.8.1), on constate aussitôt qu'on obtient, au signe près, le premier
membre de (16.8.8.2); d'où la conclusion.

\begin{proposition}[16.8.9]
\label{IV.16.8.9}
Si $D:\mathscr{F}\to\mathscr{G}$ est un opérateur différentiel d'ordre
$\leq n$, et $D':\mathscr{G}\to\mathscr{H}$ un opérateur différentiel d'ordre
$\leq n'$, alors $D'\circ D:\mathscr{F}\to\mathscr{H}$ est un opérateur
différentiel d'ordre $\leq n+n'$.
\end{proposition}

Par hypothèse, on peut écrire $D=u\circ d_{X/S,\mathscr{F}}^n$ et
$D'=v\circ d_{X/S,\mathscr{G}}^{n'}$, où
$u:\mathscr{P}_{X/S}^n\otimes_{\mathcal{O}_X}\mathscr{F}\to\mathscr{G}$ et
$v:\mathscr{P}_{X/S}^{n'}\otimes_{\mathcal{O}_X}\mathscr{G}\to\mathscr{H}$
sont des $\mathcal{O}_X$-homomorphismes. Tout revient à montrer que
l'homomorphisme composé de faisceaux de groupes additifs
\[
  \xymatrix{
    \mathscr{F}\ar[r]^-{d_{X/S,\mathscr{F}}^n}
      &\mathscr{P}_{X/S}^n\otimes_{\mathcal{O}_X}\mathscr{F}
        \ar[r]^-u
      &\mathscr{G}\ar[r]^-{d_{X/S,\mathscr{G}}^{n'}}
      &\mathscr{P}_{X/S}^{n'}\otimes_{\mathcal{O}_X}\mathscr{G}
  }
\]
se factorise en
\[
  \xymatrix{
    \mathscr{F}\ar[r]^-{d_{X/S,\mathscr{F}}^{n+n'}}
      &\mathscr{P}_{X/S}^{n+n'}\otimes_{\mathcal{O}_X}\mathscr{F}
        \ar[r]^-w
      &\mathscr{P}_{X/S}^{n'}\otimes_{\mathcal{O}_X}\mathscr{G}
  }
\]
où $w$ est un $\mathcal{O}_X$-homomorphisme. Il suffira de prouver le

\begin{lemma}[16.8.9.1]
\label{IV.16.8.9.1}
Il existe un $\mathcal{O}_X$-homomorphisme et un seul
\[
  \delta:\mathscr{P}_{X/S}^{n+n'}
  \longrightarrow
  \mathscr{P}_{X/S}^{n'}(\mathscr{P}_{X/S}^n)
  =\mathscr{P}_{X/S}^{n'}\otimes_{\mathcal{O}_X}\mathscr{P}_{X/S}^n
  \tag{16.8.9.2}
  \label{IV.16.8.9.2}
\]
\oldpage[IV-4]{44}
rendant commutatif le diagramme
\[
  \xymatrix{
    \mathcal{O}_X\ar[r]^-{d_{X/S}^{n+n'}}\ar[d]_-{d_{X/S}^n}
      &\mathscr{P}_{X/S}^{n+n'}\ar[d]^-\delta\\
    \mathscr{P}_{X/S}^n
      \ar[r]_-{d_{X/S,\mathscr{P}_{X/S}^n}^{n'}}
      &\mathscr{P}_{X/S}^{n'}(\mathscr{P}_{X/S}^n)
  }
  \tag{16.8.9.3}
  \label{IV.16.8.9.3}
\]
\end{lemma}

On aura alors en effet un diagramme commutatif déduit de (16.8.9.3) par
tensorisation avec $\mathscr{F}$
\[
  \xymatrix{
    \mathscr{F}\ar[r]^-{d_{X/S,\mathscr{F}}^{n+n'}}
      \ar[d]_-{d_{X/S,\mathscr{F}}^n}
      &\mathscr{P}_{X/S}^{n+n'}(\mathscr{F})
        \ar[d]^-{\delta\otimes1}\\
    \mathscr{P}_{X/S}^n(\mathscr{F})
      \ar[r]_-{d_{X/S,\mathscr{P}_{X/S}^n(\mathscr{F})}^{n'}}
      &\mathscr{P}_{X/S}^{n'}(\mathscr{P}_{X/S}^n(\mathscr{F}))
  }
\]
et d'autre part, on vérifie aussitôt sur la définition (16.7.5) que le
diagramme
\[
  \xymatrix{
    \mathscr{P}_{X/S}^n(\mathscr{F})\ar[r]^-u
      \ar[d]_-{d_{X/S,\mathscr{P}_{X/S}^n(\mathscr{F})}^{n'}}
      &\mathscr{G}\ar[d]^-{d_{X/S,\mathscr{G}}^{n'}}\\
    \mathscr{P}_{X/S}^{n'}(\mathscr{P}_{X/S}^n(\mathscr{F}))
      \ar[r]_-{1\otimes u}
      &\mathscr{P}_{X/S}^{n'}(\mathscr{G})
  }
\]
est commutatif. On répondra donc à la question en prenant pour $w$ le
$\mathcal{O}_X$-homomorphisme composé
\[
  \xymatrix{
    \mathscr{P}_{X/S}^{n+n'}(\mathscr{F})
      \ar[r]^-{\delta\otimes1}
      &\mathscr{P}_{X/S}^{n'}(\mathscr{P}_{X/S}^n(\mathscr{F}))
        \ar[r]^-{1\otimes u}
      &\mathscr{P}_{X/S}^{n'}(\mathscr{G}).
  }
\]

Reste à prouver le lemme (16.8.9.1). Compte tenu de (16.7.6), qui prouve
l'unicité de $\delta$, on est ramené au cas où $S=\operatorname{Spec}(A)$ et
$X=\operatorname{Spec}(B)$ sont affines; posant
$\mathfrak{I}=\mathfrak{I}_{B/A}$, il s'agit de définir un homomorphisme
canonique de $B$-modules
\[
  \varphi:(B\otimes_A B)/\mathfrak{I}^{n+n'+1}
  \longrightarrow
  ((B\otimes_A B)/\mathfrak{I}^{n'+1})
  \otimes_B((B\otimes_A B)/\mathfrak{I}^{n+1})
\]
les structures de $B$-module des deux membres provenant du premier facteur
$B$; rappelons que dans le produit tensoriel du second membre,
$(B\otimes_A B)/\mathfrak{I}^{n'+1}$ doit être considéré
\oldpage[IV-4]{45}
comme $B$-module à droite par son second facteur $B$ et
$(B\otimes_A B)/\mathfrak{I}^{n+1}$ comme $B$-module à gauche par son premier
facteur $B$ (16.7.2). Il revient au même de définir un homomorphisme de
$B$-modules
\[
  \varphi_0:B\otimes_A B\longrightarrow
  ((B\otimes_A B)/\mathfrak{I}^{n'+1})
  \otimes_B((B\otimes_A B)/\mathfrak{I}^{n+1})
\]
et de prouver qu'il s'annule dans $\mathfrak{I}^{n+n'+1}$. Or, on définit
aussitôt un tel homomorphisme par la condition que
\[
  \varphi_0(b\otimes b')
  =\pi_{n'}(b\otimes1)\otimes\pi_n(1\otimes b')
  \qquad\text{pour $b,b'$ dans $B$}
\]
avec les notations de (16.3.7). En outre, il est immédiat que $\varphi_0$ est
un homomorphisme d'\emph{anneaux}. Or, on peut écrire
\[
  \begin{aligned}
  \varphi_0(b\otimes1-1\otimes b)
  ={}&\pi_{n'}(b\otimes1-1\otimes b)\otimes\pi_n(1\otimes1)\\
  &+\pi_{n'}(1\otimes b)\otimes\pi_n(1\otimes1)
  -\pi_{n'}(1\otimes1)\otimes\pi_n(1\otimes b)
  \end{aligned}
\]
et on a
\[
  \begin{aligned}
  \pi_{n'}(1\otimes b)\otimes\pi_n(1\otimes1)
  &=\pi_{n'}(1\otimes1)b\otimes\pi_n(1\otimes1)\\
  &=\pi_{n'}(1\otimes1)\otimes b\pi_n(1\otimes1)\\
  &=\pi_{n'}(1\otimes1)\otimes\pi_n(b\otimes1)
  \end{aligned}
\]
d'où finalement
\[
  \begin{aligned}
  \varphi_0(b\otimes1-1\otimes b)
  ={}&\pi_{n'}(b\otimes1-1\otimes b)\otimes\pi_n(1\otimes1)\\
  &+\pi_{n'}(1\otimes1)\otimes\pi_n(b\otimes1-1\otimes b).
  \end{aligned}
  \tag{16.8.9.4}
  \label{IV.16.8.9.4}
\]
Un produit de $n+n'+1$ termes de la forme (16.8.9.4) est donc
nécessairement nul, car il en est ainsi d'un produit de $n+1$ termes de la
forme $\pi_n(b\otimes1-1\otimes b)$ et d'un produit de $n'+1$ termes de la
forme $\pi_{n'}(b\otimes1-1\otimes b)$. La conclusion résulte donc de
(0, 20.4.4).

\begin{corollary}[16.8.10]
\label{IV.16.8.10}
Le faisceau $\mathscr{D}\!iff_{X/S}(\mathcal{O}_X,\mathcal{O}_X)$ (aussi noté
$\mathscr{D}\!iff_{X/S}$) est canoniquement muni d'une structure de faisceau
d'anneaux, les $\mathscr{D}\!iff_{X/S}^n$ formant une filtration croissante
compatible avec cette structure.
\end{corollary}

En particulier, $\mathscr{D}\!iff_{X/S}^0$ est un faisceau de sous-anneaux de
$\mathscr{D}\!iff_{X/S}$, qui s'identifie canoniquement à $\mathcal{O}_X$
(16.8.1). Les formules (16.8.5.1) et (16.8.5.2) montrent que la structure de
$\mathcal{O}_X$-Bimodule de $\mathscr{D}\!iff_{X/S}$ provient de la
multiplication à gauche et à droite par les sections de $\mathcal{O}_X$
considéré comme faisceau de sous-anneaux de $\mathscr{D}\!iff_{X/S}$.

\begin{remarks}[16.8.11]
\label{IV.16.8.11}
\begin{enumerate}
  \item[(i)] Supposons que
  $\mathscr{F}=\bigoplus_{\lambda\in L}\mathscr{F}_\lambda$; alors il est
  clair (16.7.2.1) que
  $\mathscr{P}_{X/S}^n(\mathscr{F})=
  \bigoplus_{\lambda\in L}\mathscr{P}_{X/S}^n(\mathscr{F}_\lambda)$; comme le
  foncteur $\mathscr{F}\mapsto\Gamma(U,\mathscr{F})$ commute à la formation
  des sommes directes quelconques, $d_{X/S,\mathscr{F}}^n$ est
  l'homomorphisme dont la restriction à chaque $\mathscr{F}_\lambda$ est
  \[
    d_{X/S,\mathscr{F}_\lambda}^n:\mathscr{F}_\lambda
    \longrightarrow\mathscr{P}_{X/S}^n(\mathscr{F}_\lambda);
  \]
  on en conclut aussitôt que l'on a
  \[
    \operatorname{Diff}_{X/S}^n(\mathscr{F},\mathscr{G})
    =\prod_{\lambda\in L}
    \operatorname{Diff}_{X/S}^n(\mathscr{F}_\lambda,\mathscr{G}),
  \]
  et par suite aussi (0\textsubscript{I}, 3.2.6)
  \[
    \mathscr{D}\!iff_{X/S}^n(\mathscr{F},\mathscr{G})
    =\prod_{\lambda\in L}
    \mathscr{D}\!iff_{X/S}^n(\mathscr{F}_\lambda,\mathscr{G}).
  \]
  Par ailleurs, si $\mathscr{G}=\prod_{\mu\in M}\mathscr{G}_\mu$
  (0\textsubscript{I}, 3.2.6), on a
  \[
    \operatorname{Hom}_{\mathcal{O}_X}
    (\mathscr{P}_{X/S}^n(\mathscr{F}),\mathscr{G})
    =\prod_{\mu\in M}\operatorname{Hom}_{\mathcal{O}_X}
    (\mathscr{P}_{X/S}^n(\mathscr{F}),\mathscr{G}_\mu),
  \]
  \oldpage[IV-4]{46}
  tout homomorphisme $u$ de $\mathscr{P}_{X/S}^n(\mathscr{F})$ dans
  $\mathscr{G}$ correspondant biunivoquement à la famille de ses composés
  \[
    u_\mu:\mathscr{P}_{X/S}^n(\mathscr{F})
    \longrightarrow\mathscr{G}\longrightarrow\mathscr{G}_\mu.
  \]
  On a donc
  \[
    \operatorname{Diff}_{X/S}^n(\mathscr{F},\mathscr{G})
    =\prod_{\mu\in M}
    \operatorname{Diff}_{X/S}^n(\mathscr{F},\mathscr{G}_\mu),
  \]
  et par suite aussi
  \[
    \mathscr{D}\!iff_{X/S}^n(\mathscr{F},\mathscr{G})
    =\prod_{\mu\in M}
    \mathscr{D}\!iff_{X/S}^n(\mathscr{F},\mathscr{G}_\mu).
  \]
  \item[(ii)] Jusqu'ici, on n'a guère rencontré d'opérateurs différentiels
  $\mathscr{F}\to\mathscr{G}$ que lorsque $\mathscr{F}$ et $\mathscr{G}$ sont
  des $\mathcal{O}_X$-Modules localement libres de rang fini, auquel cas, leur
  structure se ramène localement, en vertu de (i), à celle du faisceau
  $\mathscr{D}\!iff_{X/S}$; cette dernière va être étudiée plus loin (16.11)
  dans un cas particulier.
\end{enumerate}
\end{remarks}

\subsection{Immersions régulières et quasi-régulières.}
\label{IV.16.9}

\begin{definition}[16.9.1]
\label{IV.16.9.1}
Soit $X$ un espace annelé. On dit qu'un Idéal $\mathscr{I}$ de
$\mathcal{O}_X$ est \emph{régulier} (resp. \emph{quasi-régulier}) si, pour
tout point $x\in\operatorname{Supp}(\mathcal{O}_X/\mathscr{I})$, il existe un
voisinage ouvert $U$ de $x$ dans $X$ et une suite régulière (0, 15.2.2) (resp.
quasi-régulière (0, 15.2.2)) d'éléments de $\Gamma(U,\mathcal{O}_X)$ qui
engendre $\mathscr{I}|U$.
\end{definition}

On dira qu'une suite régulière (resp. quasi-régulière) de sections de
$\mathcal{O}_X$ au-dessus de $U$ qui engendre $\mathscr{I}|U$ est un
\emph{système régulier} (resp. \emph{quasi-régulier}) \emph{de générateurs}
de $\mathscr{I}|U$.

\begin{definition}[16.9.2]
\label{IV.16.9.2}
Soit $j:Y\to X$ une immersion de préschémas et soit $U$ un ouvert de $X$ tel
que $j(Y)\subset U$ et que $j$ soit une immersion fermée de $Y$ dans $U$. On
dit que $j$ est \emph{régulière} (resp. \emph{quasi-régulière}) si le
sous-préschéma fermé $j(Y)$ de $U$ associé à $j$ est défini par un Idéal
régulier (resp. quasi-régulier) de $\mathcal{O}_U$ (condition indépendante de
l'ouvert $U$ choisi).
\end{definition}

On dit qu'un sous-préschéma $Y$ d'un préschéma $X$ est \emph{régulièrement
immergé} (resp. \emph{quasi-régulièrement immergé}) si l'injection canonique
$j:Y\to X$ est une immersion régulière (resp. quasi-régulière). Si $Y$ est un
sous-préschéma fermé de $X$ et $\mathscr{I}$ l'Idéal de $\mathcal{O}_X$ qui
définit $Y$, il revient au même de dire que $\mathscr{I}$ est \emph{régulier}
(resp. \emph{quasi-régulier}).

Par exemple, si $A$ est un anneau \emph{intègre}, $f$ un élément $\neq0$ de
$A$, le sous-préschéma fermé $V(f)$ de $\operatorname{Spec}(A)$ (isomorphe à
$\operatorname{Spec}(A/fA)$) est \emph{régulièrement immergé} dans
$\operatorname{Spec}(A)$.

Tout Idéal régulier est quasi-régulier (0, 15.2.2); toute immersion régulière
est quasi-régulière (cf. (16.9.11) pour une réciproque).

\begin{proposition}[16.9.3]
\label{IV.16.9.3}
Soient $X$ un espace annelé, $\mathscr{I}$ un Idéal de $\mathcal{O}_X$,
$(f_i)_{1\leq i\leq m}$ une suite finie de sections de $\mathcal{O}_X$
au-dessus de $X$ engendrant $\mathscr{I}$. Pour que $(f_i)$ soit une suite
quasi-régulière (0, 15.2.2), il faut et il suffit que les conditions suivantes
soient vérifiées:
\begin{enumerate}
  \item[(i)] Les images canoniques des $f_i$ dans
  $\mathscr{I}/\mathscr{I}^2$ forment une base de cet
  $\mathcal{O}_X/\mathscr{I}$-Module.
  \item[(ii)] L'homomorphisme surjectif canonique (16.1.2.2)
  \[
    \mathbf{S}_{\mathcal{O}_X/\mathscr{I}}^{\bullet}
    (\mathscr{I}/\mathscr{I}^2)
    \longrightarrow\mathscr{G}r_{\mathscr{I}}^{\bullet}(\mathcal{O}_X)
  \]
  est bijectif.
\end{enumerate}
En outre, s'il en est ainsi, toute suite $(f'_i)_{1\leq i\leq n}$ de $n$
sections de $\mathscr{I}$ au-dessus de $X$ qui engendre $\mathscr{I}$ est
quasi-régulière.
\end{proposition}

\oldpage[IV-4]{47}

Les deux conditions de l'énoncé ne font que traduire la définition donnée dans
(0, 15.2.2), compte tenu de la définition des homomorphismes canoniques
(0, 15.2.1.1). La dernière assertion résulte de ce que, si un module $M$ sur
un anneau commutatif $A$ admet une base de $n$ éléments, tout système de
générateurs de $M$ ayant $n$ éléments est une base de $M$ (Bourbaki,
\emph{Alg. comm.}, chap.~II, §~3, cor.~5 du th.~1).

\begin{corollary}[16.9.4]
\label{IV.16.9.4}
Soient $X$ un espace annelé en anneaux locaux, $\mathscr{I}$ un Idéal de
$\mathcal{O}_X$. Pour que $\mathscr{I}$ soit quasi-régulier, il faut et il
suffit qu'il vérifie les conditions suivantes:
\begin{enumerate}
  \item[(i)] $\mathscr{I}$ est de type fini.
  \item[(ii)] $\mathscr{I}/\mathscr{I}^2$ est un
  $(\mathcal{O}_X/\mathscr{I})$-Module localement libre.
  \item[(iii)] L'homomorphisme canonique
  \[
    \mathbf{S}_{\mathcal{O}_X/\mathscr{I}}^{\bullet}
    (\mathscr{I}/\mathscr{I}^2)
    \longrightarrow\mathscr{G}r_{\mathscr{I}}^{\bullet}(\mathcal{O}_X)
    \tag{16.9.4.1}
    \label{IV.16.9.4.1}
  \]
  est bijectif.
\end{enumerate}
\end{corollary}

La nécessité des conditions résulte aussitôt de (16.9.3). Pour voir que ces
conditions sont suffisantes, il faut montrer, en vertu de (16.9.3), que si, en
un point $x\in\operatorname{Supp}(\mathcal{O}_X/\mathscr{I})$, il existe un
voisinage ouvert $U$ de $x$ dans $X$ et $n$ sections $f_i$ ($1\leq i\leq n$)
de $\mathscr{I}$ au-dessus de $U$ dont les images canoniques dans
$\mathscr{I}/\mathscr{I}^2$ forment une base de
$(\mathscr{I}/\mathscr{I}^2)|U$ sur
$(\mathcal{O}_X/\mathscr{I})|U$, alors il y a un voisinage ouvert
$V\subset U$ de $x$ tel que les $f_i|V$ engendrent $\mathscr{I}|V$. Or, par
hypothèse, on a $\mathscr{I}_x\neq\mathcal{O}_x$, donc $\mathscr{I}_x$ est
contenu dans l'idéal maximal de $\mathcal{O}_x$; comme $\mathscr{I}_x$ est un
$\mathcal{O}_x$-module de type fini et que les classes des $(f_i)_x$ dans
$\mathscr{I}_x/\mathscr{I}_x^2$ engendrent cet
$(\mathcal{O}_x/\mathscr{I}_x)$-module, le lemme de Nakayama montre que les
$(f_i)_x$ engendrent $\mathscr{I}_x$. Comme $\mathscr{I}$ est de type fini,
on conclut par (0\textsubscript{I}, 5.2.2).

\begin{corollary}[16.9.5]
\label{IV.16.9.5}
Soient $X$ un espace annelé en anneaux locaux, $\mathscr{I}$ un Idéal
quasi-régulier de $\mathcal{O}_X$, $(f_i)_{1\leq i\leq n}$ une suite de
sections de $\mathscr{I}$ au-dessus de $X$,
$x$ un point de $\operatorname{Supp}(\mathcal{O}_X/\mathscr{I})$. Les
conditions suivantes sont équivalentes:
\begin{enumerate}
  \item[(a)] Il existe un voisinage ouvert $U$ de $x$ dans $X$ tel que les
  $f_i|U$ forment une suite quasi-régulière d'éléments de
  $\Gamma(U,\mathcal{O}_X)$ engendrant $\mathscr{I}|U$.
  \item[(b)] Les $(f_i)_x$ forment un système de générateurs de
  $\mathscr{I}_x$ dont le nombre d'éléments est le plus petit possible.
  \item[(b')] Les $(f_i)_x$ forment un système minimal de générateurs de
  $\mathscr{I}_x$.
  \item[(c)] Si $\bar f_i$ est l'image canonique de $f_i$ dans
  $\Gamma(X,\mathscr{I}/\mathscr{I}^2)$, les $(\bar f_i)_x$ forment une base
  du $(\mathcal{O}_x/\mathscr{I}_x)$-module
  $\mathscr{I}_x/\mathscr{I}_x^2$.
\end{enumerate}
\end{corollary}

Par hypothèse, $\mathcal{O}_x$ est un anneau local, $\mathscr{I}_x$ un idéal
de type fini de $\mathcal{O}_x$ contenu dans l'idéal maximal de
$\mathcal{O}_x$; l'équivalence de b), b') et c) résulte donc du lemme de
Nakayama (Bourbaki, \emph{Alg. comm.}, chap.~II, §~3,
n\textsuperscript{o}~2, prop.~5). Il est clair que a) entraîne c) en vertu de
(16.9.3); d'autre part, il résulte de (0\textsubscript{I}, 5.2.2) que si la
condition c) est vérifiée (donc aussi b)), il y a un voisinage $U$ de $x$ dans
$X$ tel que $(\mathscr{I}/\mathscr{I}^2)|U$ ait un rang constant égal à $n$,
et que les $f_i|U$ engendrent $\mathscr{I}|U$; il suffit donc d'appliquer dans
$U$ la dernière assertion de (16.9.3).

\begin{remarks}[16.9.6]
\label{IV.16.9.6}
\begin{enumerate}
  \item[(i)] Sous les hypothèses générales de (16.9.5), il ne suffit pas que
  les $(\bar f_i)_y$ forment une base du
  $(\mathcal{O}_y/\mathscr{I}_y)$-module
  $\mathscr{I}_y/\mathscr{I}_y^2$ pour tout $y\in X$ pour que la suite
  $(f_i)$ engendre $\mathscr{I}$. On en a un exemple en
\oldpage[IV-4]{48}
  prenant $X=\operatorname{Spec}(A)$, où $A$ est un anneau de Dedekind, et
  $\mathscr{I}=\widetilde{\mathfrak{I}}$, où $\mathfrak{I}$ est un idéal
  premier non principal de $A$; on a alors en effet
  $\mathscr{I}_y/\mathscr{I}_y^2=0$ en tout point $y$ distinct du point
  $x\in X$ correspondant à $\mathfrak{I}$, et
  $\mathscr{I}_x/\mathscr{I}_x^2$ est de rang $1$ sur le corps
  $\mathcal{O}_x/\mathscr{I}_x$; en outre $\mathscr{I}$ est évidemment un
  Idéal régulier.
  \item[(ii)] Dans (16.9.5), on ne peut remplacer « quasi-régulière » par
  « régulière », même lorsque $X$ est un préschéma (cf. (16.9.12)).
  Désignons en effet par $B$ l'anneau des germes de fonctions indéfiniment
  différentiables au point $0$ dans $\mathbb{R}$; il a un idéal maximal
  $\mathfrak{m}$ engendré par le germe $t$ de l'application identique de
  $\mathbb{R}$ au point $0$, et l'intersection $\mathfrak{n}$ des
  $\mathfrak{m}^k$ pour $k>0$ n'est pas réduite à $0$. Soit maintenant $A$
  l'anneau quotient $B[T]/\mathfrak{n}TB[T]$, et soient $f_1,f_2$ les images
  canoniques dans $A$ des éléments $t$ et $T$ de $B[T]$. La suite
  $(f_1,f_2)$ est régulière dans $A$: en effet, $f_1$ n'est pas diviseur de
  $0$ dans $A$, car la relation $tP[T]\in\mathfrak{n}TB[T]$, pour un
  polynôme $P\in B[T]$, entraîne que les produits de $t$ par les coefficients
  de $P$ appartiennent à l'idéal $\mathfrak{n}$, et il en résulte aussitôt que
  ces coefficients sont eux-mêmes dans $\mathfrak{n}$, donc
  $P[T]\in\mathfrak{n}TB[T]$. Comme $B/tB$ est isomorphe à $\mathbb{R}$,
  $A/f_1A$ est isomorphe à l'anneau de polynômes $\mathbb{R}[T]$, donc
  intègre, et l'image de $f_2$ dans $A/f_1A$, étant égale à $T$, n'est donc
  pas diviseur de $0$, ce qui prouve notre assertion. Cependant $f_2$ est
  diviseur de $0$ dans $A$, car pour tout élément $x\in\mathfrak{n}$ non nul,
  l'image de $x$ dans $A$ est $\ne0$, mais l'image de $xT$ est nulle. On en
  conclut que la suite $(f_2,f_1)$ n'est pas régulière dans $A$; d'autre part
  l'idéal $\mathfrak{I}=f_1A+f_2A$ est distinct de $A$, donc les conditions
  b), b') et c) de (16.9.5) n'entraînent pas la condition a) lorsqu'on y
  remplace « quasi-régulière » par « régulière ».
\end{enumerate}
\end{remarks}

\begin{env}[16.9.7]
\label{IV.16.9.7}
Si $X=\operatorname{Spec}(A)$ est un schéma affine, nous dirons qu'un idéal
$\mathfrak{I}$ de $A$ est régulier (resp. quasi-régulier) si l'Idéal
$\mathscr{I}=\widetilde{\mathfrak{I}}$ de $\mathcal{O}_X$ est régulier
(resp. quasi-régulier); on notera que cette notion est locale et n'implique
nullement l'existence d'un système de générateurs de $\mathfrak{I}$ formant
dans $A$ une suite régulière (resp. quasi-régulière) comme le montre l'exemple
(16.9.6, (i)); il en est toutefois ainsi lorsque $A$ est local (16.9.5).

La proposition (16.9.4) se traduit en termes d'immersions quasi-régulières de
la façon suivante:
\end{env}

\begin{proposition}[16.9.8]
\label{IV.16.9.8}
Soit $j:Y\to X$ un morphisme de préschémas; pour que $j$ soit une immersion
quasi-régulière, il faut et il suffit que $j$ satisfasse aux conditions
suivantes:
\begin{enumerate}
  \item[(i)] $j$ est une immersion localement de présentation finie.
  \item[(ii)] Le faisceau conormal
  $\operatorname{Gr}^1(j)=\mathcal{N}_{Y/X}$ (16.1.2) est un
  $\mathcal{O}_Y$-Module localement libre.
  \item[(iii)] L'homomorphisme canonique
  \[
    \operatorname{S}_{\mathcal{O}_Y}^{\bullet}
      (\operatorname{Gr}^1(j))\longrightarrow\operatorname{Gr}^{\bullet}(j)
  \]
  (16.1.2.2) est bijectif.
\end{enumerate}
\end{proposition}

La question étant locale sur $Y$, on peut se borner au cas où $j$ est
l'injection canonique d'un sous-préschéma fermé $Y$ de $X$, et alors la
traduction de (16.9.4) en (16.9.8) résulte de l'explicitation de
$\operatorname{Gr}^1(j)$ et $\operatorname{Gr}^{\bullet}(j)$ en termes de
l'Idéal $\mathscr{I}$ de $\mathcal{O}_X$ définissant le sous-préschéma $Y$
(16.1.3, (ii)).

\begin{corollary}[16.9.9]
\label{IV.16.9.9}
Soient $Y$ un préschéma, $X$ un $Y$-préschéma, $j:Y\to X$ une $Y$-section
de $X$, de sorte que le $n$-ème invariant normal $\mathscr{A}^{(n)}$ de $j$
(16.1.2) est une $\mathcal{O}_Y$-Algèbre augmentée (16.1.7); posons
$\mathscr{A}^{(\infty)}=\varprojlim\mathscr{A}^{(n)}$. Pour que $j$ soit une
immersion quasi-régulière, il faut et il suffit que $j$ soit localement de
présentation finie, et que tout $y\in Y$ admette un voisinage ouvert affine
$U$ d'anneau $C$ tel que $\mathscr{A}^{(\infty)}|U$ soit isomorphe comme
$\mathcal{O}_U$-Algèbre topologique augmentée à
$\mathcal{O}_U[[T_1,\ldots,T_n]]$.
\end{corollary}

On peut se borner au cas où $j$ est une immersion fermée en se restreignant
à un voisinage assez petit de $y$ (voir le raisonnement de (16.4.11)), et
alors $\mathcal{O}_Y$ s'identifie à une Algèbre quotient
$\mathcal{O}_X/\mathscr{I}$ et l'homomorphisme surjectif canonique
$\mathcal{O}_X\to\mathcal{O}_Y$ admet un
\oldpage[IV-4]{49}
inverse à droite (16.1.7). On peut donc supposer
$X=\operatorname{Spec}(B)$ et $Y=\operatorname{Spec}(A)$ affines, $B$ étant
une $A$-algèbre augmentée, et l'idéal d'augmentation $\mathfrak{I}$ étant de
type fini. Comme $\mathscr{A}^{(n)}$ s'identifie alors à
$(B/\mathfrak{I}^{n+1})^{\sim}$, le corollaire résulte de l'équivalence de b)
et c) dans (0, 19.5.4) puisque $B/\mathfrak{I}=A$.

On notera que, dans le cas affine envisagé, le fait que $j$ soit une immersion
quasi-régulière équivaut encore, en vertu de (0, 19.5.4), à dire que $B$ est
une $A$-algèbre formellement lisse pour la topologie
$\mathfrak{I}$-pré-adique.

On notera aussi que la condition que $j$ soit une immersion localement de
présentation finie est toujours vérifiée lorsque le morphisme $X\to Y$ est
localement de type fini (1.4.3, (v)).

\begin{proposition}[16.9.10]
\label{IV.16.9.10}
Soient $X$ un préschéma localement noethérien, $Y$ un sous-préschéma de $X$,
$j:Y\to X$ l'injection canonique, $y$ un point de $Y$.
\begin{enumerate}
  \item[(i)] Pour qu'il existe un voisinage ouvert $U$ de $y$ dans $X$ tel que
  la restriction $Y\cap U\to U$ de $j$ soit une immersion régulière, il faut
  et il suffit que le noyau $\mathscr{I}_y$ de l'homomorphisme surjectif
  $\mathcal{O}_{X,y}\to\mathcal{O}_{Y,y}$ soit engendré par une suite régulière
  d'éléments de $\mathcal{O}_{X,y}$.
  \item[(ii)] Pour que l'immersion $j$ soit régulière, il faut et il suffit
  qu'elle soit quasi-régulière.
\end{enumerate}
\end{proposition}

\begin{enumerate}
  \item[(i)] On peut se borner au cas où $Y$ est un sous-préschéma fermé de
  $X$ défini par un Idéal \emph{cohérent} $\mathscr{I}$ de $\mathcal{O}_X$.
  La condition est évidemment nécessaire. Inversement, si $\mathscr{I}_y$
  est engendré par une suite régulière $(s_i)_y$, où les $s_i$ sont des
  sections de $\mathscr{I}$ au-dessus d'un voisinage ouvert $U$ de $y$ dans
  $X$, on peut supposer que les $s_i$ engendrent $\mathscr{I}|U$
  (0\textsubscript{I}, 5.2.2) et forment une suite régulière (0, 15.2.4),
  d'où l'assertion.
  \item[(ii)] Le fait qu'une immersion quasi-régulière soit régulière résulte
  de (i) et de l'identité des suites quasi-régulières et des suites régulières
  dans $\mathcal{O}_{X,y}$, formées d'éléments de l'idéal maximal
  (0, 15.1.11).
\end{enumerate}

Lorsque (sans hypothèse noethérienne sur $X$) le noyau $\mathscr{I}_y$ de
$\mathcal{O}_{X,y}\to\mathcal{O}_{Y,y}$ est engendré par une suite régulière
d'éléments de $\mathcal{O}_{X,y}$, on dit que l'immersion $j$ est
\emph{régulière au point $y$}.

\begin{corollary}[16.9.11]
\label{IV.16.9.11}
Soit $X$ un préschéma localement noethérien; alors tout Idéal quasi-régulier de
$\mathcal{O}_X$ est régulier.
\end{corollary}

\begin{remarks}[16.9.12]
\label{IV.16.9.12}
\begin{enumerate}
  \item[(i)] On notera qu'une immersion régulière n'est pas en général un
  morphisme plat, ni \emph{a fortiori} un morphisme \emph{régulier} au sens de
  (6.8.1).
  \item[(ii)] Soit $A$ un anneau local noethérien; il résulte aussitôt de
  (16.9.4) et de (0, 17.1.1) que pour que $A$ soit \emph{régulier}, il faut et
  il suffit que son idéal maximal $\mathfrak{m}$ soit \emph{quasi-régulier}
  (ou \emph{régulier}, ce qui revient au même puisque $A$ est noethérien).
  Pour qu'un schéma affine noethérien $X$ soit \emph{régulier}, il faut et il
  suffit que pour tout point fermé $x\in X$, l'injection canonique
  $\operatorname{Spec}(k(x))\to X$ soit une immersion \emph{régulière}.
\end{enumerate}
\end{remarks}

\begin{proposition}[16.9.13]
\label{IV.16.9.13}
Soient $X$ un préschéma localement noethérien, $Y$ un sous-préschéma de $X$,
$Y'$ un sous-préschéma de $Y$, tel que l'injection canonique $j:Y'\to Y$ soit
régulière. Alors la suite de $\mathcal{O}_{Y'}$-Modules
\[
  \xymatrix{
    0\ar[r] & j^*(\mathcal{N}_{Y/X})\ar[r]
      & \mathcal{N}_{Y'/X}\ar[r] & \mathcal{N}_{Y'/Y}\ar[r] & 0
  }
  \tag{16.9.13.1}
  \label{IV.16.9.13.1}
\]
\oldpage[IV-4]{50}
est exacte; en outre, pour tout $x\in X$, il existe un voisinage ouvert $U$
de $x$ tel que les restrictions à $U$ des homomorphismes de (16.9.13.1)
forment une suite exacte et scindée.
\end{proposition}

Prouvons d'abord le lemme suivant:

\begin{lemma}[16.9.13.2]
\label{IV.16.9.13.2}
Soient $A$ un anneau, $\mathfrak{I}$ un idéal de $A$,
$A'=A/\mathfrak{I}$, $(f_i)_{1\leq i\leq r}$ une suite d'éléments de $A$ qui
est $A'$-régulière, $\mathfrak{K}=\sum_i f_iA$,
$\mathfrak{L}=\mathfrak{I}+\mathfrak{K}$,
$\mathfrak{K}'=\sum_i f_iA'$, de sorte que $C=A/\mathfrak{L}$ est isomorphe à
$A'/\mathfrak{K}'$. Alors pour tout entier $n>0$ et tout entier $N\geq n$,
on a la relation
\[
  \mathfrak{I}\cap\mathfrak{K}^n
  =\mathfrak{I}\mathfrak{K}^n+(\mathfrak{I}\cap\mathfrak{K}^N).
  \tag{16.9.13.3}
  \label{IV.16.9.13.3}
\]
\end{lemma}

Il suffit évidemment de prouver que tout élément du premier membre est
contenu dans le second, et, par récurrence sur $n$, on se ramène au cas où
$N=n+1$. Un élément du premier membre de (16.9.13.3), étant dans
$\mathfrak{K}^n$, s'écrit $P(f_1,\ldots,f_r)$, où
$P\in A[T_1,\ldots,T_r]$ est homogène et de degré $n$. Si $f_i'$ est l'image
canonique de $f_i$ dans $A'$, l'hypothèse
$P(f_1,\ldots,f_r)\in\mathfrak{I}$ signifie que
$P(f_1',\ldots,f_r')=0$. Mais
$P(f_1',\ldots,f_r')\in\mathfrak{K}'^n$, donc l'image canonique de
$P(f_1',\ldots,f_r')$ dans
$\mathfrak{K}'^n/\mathfrak{K}'^{n+1}$ est nulle. Or l'hypothèse que la suite
$(f_i)$ est $A'$-régulière implique que l'homomorphisme canonique
\[
  \mathbf{S}_C^n(\mathfrak{K}'/\mathfrak{K}'^2)
  \longrightarrow\mathfrak{K}'^n/\mathfrak{K}'^{n+1}
\]
est bijectif (0, 15.1.9); on conclut que les coefficients de $P$ appartiennent
à $\mathfrak{L}=\mathfrak{I}+\mathfrak{K}$. Il en résulte aussitôt que l'on a
$P(f_1,\ldots,f_r)\in\mathfrak{I}\mathfrak{K}^n+\mathfrak{K}^{n+1}$, et comme
$P(f_1,\ldots,f_r)\in\mathfrak{I}$, on a finalement
\[
  P(f_1,\ldots,f_r)
  \in\mathfrak{I}\mathfrak{K}^n
    +(\mathfrak{I}\cap\mathfrak{K}^{n+1}),
\]
ce qui prouve le lemme.

En prenant les quotients des deux membres de (16.9.13.3) par
$\mathfrak{I}\mathfrak{K}^n$, on voit que les relations (16.9.13.3) pour
$N\geq n$ entraînent
\[
  (\mathfrak{I}\cap\mathfrak{K}^n)/\mathfrak{I}\mathfrak{K}^n
  \subset\bigcap_{N\geq n}\mathfrak{K}^N
    \cdot(A/(\mathfrak{I}\mathfrak{K}^n)).
  \tag{16.9.13.4}
  \label{IV.16.9.13.4}
\]
On en déduit le

\begin{corollary}[16.9.13.5]
\label{IV.16.9.13.5}
Supposons vérifiées les hypothèses de (16.9.13.2) et supposons en outre que
l'anneau $A$ soit noethérien et que $\mathfrak{K}$ soit contenu dans le
radical de $A$. Alors on a pour tout entier $n>0$,
\[
  \mathfrak{I}\cap\mathfrak{K}^n=\mathfrak{I}\mathfrak{K}^n.
  \tag{16.9.13.6}
  \label{IV.16.9.13.6}
\]
\end{corollary}

En effet, le second membre de (16.9.13.4) est alors nul, puisque
$A/\mathfrak{I}\mathfrak{K}^n$ est un $A$-module de type fini (Bourbaki,
\emph{Alg. comm.}, chap.~III, §~3, n\textsuperscript{o}~3, prop.~6).

Prenons en particulier $n=2$ dans (16.9.13.6), et remarquons que l'on a
$\mathfrak{L}^2=\mathfrak{I}^2+\mathfrak{I}\mathfrak{K}+\mathfrak{K}^2
=\mathfrak{I}\mathfrak{L}+\mathfrak{K}^2$; puisque
$\mathfrak{I}\mathfrak{L}\subset\mathfrak{L}^2$, on en déduit que
\[
  \mathfrak{I}\cap\mathfrak{L}^2
  =\mathfrak{I}\mathfrak{L}+(\mathfrak{I}\cap\mathfrak{K}^2)
  =\mathfrak{I}\mathfrak{L}+\mathfrak{I}\mathfrak{K}^2
  =\mathfrak{I}\mathfrak{L},
\]
autrement dit
\[
  \mathfrak{I}\cap\mathfrak{L}^2=\mathfrak{I}\mathfrak{L},
  \tag{16.9.13.7}
  \label{IV.16.9.13.7}
\]
ce qu'on peut encore exprimer en disant que l'homomorphisme canonique
\[
  \mathfrak{I}/\mathfrak{I}\mathfrak{L}
  \longrightarrow(\mathfrak{I}+\mathfrak{L}^2)/\mathfrak{L}^2
\]
est bijectif.

Ces lemmes étant démontrés, prouvons la première assertion de (16.9.13): il
suffit
\oldpage[IV-4]{51}
évidemment de prouver que la suite des fibres des faisceaux figurant dans
(16.9.13.1), en un point $x\in Y'$, est exacte. Or, si l'on pose
$A=\mathcal{O}_{X,x}$, on peut écrire
$\mathcal{O}_{Y,x}=A'=A/\mathfrak{I}$ où $\mathfrak{I}$ est un idéal contenu
dans l'idéal maximal de $A$, puis
$\mathcal{O}_{Y',x}=A'/\mathfrak{K}'$, où $\mathfrak{K}'$ est engendré par une
suite $A'$-régulière d'éléments de $A'$, eux-mêmes images des éléments d'une
suite $A'$-régulière d'éléments de $A$ appartenant à l'idéal maximal de $A$.
Si $\mathfrak{K}$ est l'idéal engendré par ces derniers et
$\mathfrak{L}=\mathfrak{I}+\mathfrak{K}$, on a
$\mathcal{O}_{Y',x}=A/\mathfrak{L}$, et comme on est dans la situation de
(16.9.13.5), l'homomorphisme canonique
\[
  \mathfrak{I}/\mathfrak{I}\mathfrak{L}
  \longrightarrow(\mathfrak{I}+\mathfrak{L}^2)/\mathfrak{L}^2
\]
est bijectif. Mais cela montre que la suite
\[
  0\longrightarrow\mathfrak{I}/\mathfrak{I}\mathfrak{L}
  \longrightarrow\mathfrak{L}/\mathfrak{L}^2
  \longrightarrow(\mathfrak{L}/\mathfrak{I})/
    (\mathfrak{L}/\mathfrak{I})^2\longrightarrow0
\]
est exacte (voir la démonstration de (16.2.7)), et les modules figurant dans
cette suite sont précisément les fibres en $x$ des faisceaux de (16.9.13.1).
La seconde assertion résulte de ce que $\mathcal{N}_{Y'/Y}$ est un
$\mathcal{O}_{Y'}$-Module localement \emph{libre} (16.9.8) et de Bourbaki,
\emph{Alg.}, chap.~II, 3\textsuperscript{e} éd., §~1, n\textsuperscript{o}~11,
prop.~21.

\subsection{Morphismes différentiellement lisses}
\label{IV.16.10}

\begin{definition}[16.10.1]
\label{IV.16.10.1}
On dit qu'un morphisme de préschémas $f:X\to S$ est
\emph{différentiellement lisse} (ou que $X$ est \emph{différentiellement
lisse sur $S$}) s'il vérifie les conditions suivantes:
\begin{enumerate}
  \item[(i)] $\Omega_{X/S}^1$ est un $\mathcal{O}_X$-Module localement
  projectif, c'est-à-dire que tout point de $X$ admet un voisinage ouvert
  affine $U$ tel que $\Gamma(U,\Omega_{X/S}^1)$ soit un
  $\Gamma(U,\mathcal{O}_X)$-module projectif (non nécessairement de type
  fini).
  \item[(ii)] L'homomorphisme canonique (16.3.1.1)
  \[
    \mathbf{S}_{\mathcal{O}_X}^{\bullet}(\Omega_{X/S}^1)
    \longrightarrow\mathscr{G}r_{\bullet}(\mathscr{P}_{X/S})
  \]
  est bijectif.
\end{enumerate}
En particulier, si $\Omega_{X/S}^1$ est localement libre de rang fini, les
$\mathscr{P}_{X/S}^n$ sont des $\mathcal{O}_X$-Modules localement libres de
rang fini (étant des extensions de tels Modules).
\end{definition}

On dit que $f$ est \emph{différentiellement lisse en un point $x\in X$} (ou
que $X$ est différentiellement lisse sur $S$ au point $x$) s'il existe un
voisinage ouvert $U$ de $x$ dans $X$ tel que $f|U$ soit différentiellement
lisse.

Nous verrons plus loin (17.12.4) qu'un morphisme lisse est différentiellement
lisse, ce qui justifie la terminologie; mais la réciproque est inexacte; en
effet, un \emph{monomorphisme} $f:X\to S$ est différentiellement lisse,
puisque $\Omega_{X/S}^1=0$ en vertu de (I, 5.3.8), et par suite
l'homomorphisme surjectif (16.3.1.1) est évidemment bijectif; or un
monomorphisme n'est même pas nécessairement plat, ni \emph{a fortiori} lisse.
Bornons-nous ici à noter la proposition suivante:

\begin{proposition}[16.10.2]
\label{IV.16.10.2}
Soient $A$ un anneau, $B$ une $A$-algèbre formellement lisse pour les
topologies discrètes (0, 19.3.1). Alors $\operatorname{Spec}(B)$ est
différentiellement lisse sur $\operatorname{Spec}(A)$.
\end{proposition}

En effet, $B\otimes_A B$ est alors (pour les topologies discrètes) une
$B$-algèbre formellement lisse (pour l'un ou l'autre des homomorphismes
canoniques $b\mapsto b\otimes1$, $b\mapsto1\otimes b$ de $B$
\oldpage[IV-4]{52}
dans $B\otimes_A B$) (0, 19.3.5, (iii)); donc $B\otimes_A B$ est aussi une
$A$-algèbre formellement lisse pour les topologies discrètes (0, 19.3.5,
(ii)). Posant $\mathfrak{I}=\mathfrak{I}_{B/A}$, il en résulte que
$B\otimes_A B$ est aussi une $A$-algèbre formellement lisse pour la topologie
$\mathfrak{I}$-préadique (0, 19.3.8); comme par hypothèse
$B=(B\otimes_A B)/\mathfrak{I}$ est une $A$-algèbre formellement lisse pour
les topologies discrètes, la proposition résulte de l'équivalence de a) et b)
dans (0, 19.5.4).

\begin{proposition}[16.10.3]
\label{IV.16.10.3}
Pour qu'un morphisme $f:X\to S$ soit différentiellement lisse, il faut et il
suffit que pour tout $x\in X$, il existe un voisinage ouvert affine $U$ de
$x$, d'anneau $A$, tel que $\Gamma(U,\mathscr{P}_{X/S}^{\infty})$ soit une
$A$-algèbre topologique augmentée isomorphe à l'algèbre complétée $\widehat B$,
où $B=\mathbf{S}_A(V)$, $V$ étant un $A$-module projectif et $B$ étant muni de
la topologie $B^+$-préadique (où $B^+$ est l'idéal d'augmentation). Si
$\Omega_{X/S}^1$ est localement libre de rang fini, on peut remplacer
$\widehat B$ par l'algèbre de séries formelles $A[[T_1,\ldots,T_n]]$.
\end{proposition}

La notion de morphisme différentiellement lisse étant évidemment locale sur
$X$, on peut se borner au cas où $S=\operatorname{Spec}(B)$,
$X=\operatorname{Spec}(C)$. Considérons $C\otimes_B C$ comme une $C$-algèbre
(pour le premier facteur); posons $\mathfrak{I}=\mathfrak{I}_{C/B}$ et
munissons $C\otimes_B C$ de la topologie $\mathfrak{I}$-préadique; on peut
appliquer à la $C$-algèbre topologique $C\otimes_B C$ et à l'idéal
$\mathfrak{I}$ de $C\otimes_B C$ l'équivalence de b) et c) dans (0, 19.5.4),
puisque $(C\otimes_B C)/\mathfrak{I}=C$ est évidemment une $C$-algèbre
formellement lisse pour les topologies discrètes. La topologie sur
$\Gamma(U,\mathscr{P}_{X/S}^{\infty})$ est évidemment la topologie de limite
projective de cet anneau (16.1.11).

On notera que l'entier $n$ de l'énoncé de (16.10.3) est le \emph{rang} de
$\Omega_{X/S}^1$ au point $x$. Nous verrons plus loin (17.13.5) que lorsque
$f$ est différentiellement lisse et localement de type fini, $n$ est aussi
égal à la dimension de la fibre $f^{-1}(f(x))$ au point $x$.

\begin{proposition}[16.10.4]
\label{IV.16.10.4}
Soient $f:X\to S$, $g:S'\to S$ deux morphismes, et posons
$X'=X\times_S S'$, $f'=f_{(S')}:X'\to S'$.
\begin{enumerate}
  \item[(i)] Si $f$ est différentiellement lisse, il en est de même de $f'$.
  \item[(ii)] Inversement, si $g$ est fidèlement plat et quasi-compact, et si
  $f'$ est différentiellement lisse et $\Omega_{X'/S'}^1$ un
  $\mathcal{O}_{X'}$-Module de type fini, $f$ est différentiellement lisse et
  $\Omega_{X/S}^1$ est un $\mathcal{O}_X$-Module de type fini.
\end{enumerate}
\end{proposition}

En effet, si $f$ est différentiellement lisse, les
$\mathscr{G}r_n(\mathscr{P}_{X/S})$ sont des $\mathcal{O}_X$-Modules
\emph{plats}; par suite (16.4.6), l'homomorphisme
\[
  \mathscr{G}r_n(\mathscr{P}_{X/S})\otimes_{\mathcal{O}_X}\mathcal{O}_{X'}
  \longrightarrow\mathscr{G}r_n(\mathscr{P}_{X'/S'})
\]
est bijectif pour tout $n$, et en vertu de la commutativité du diagramme
(16.2.1.3), il résulte de la définition (16.10.1) que $f'$ est
différentiellement lisse. D'autre part, si $g$ est fidèlement plat et
quasi-compact, il résulte encore de (16.4.6) que
\[
  \mathscr{G}r_n(\mathscr{P}_{X/S})\otimes_{\mathcal{O}_X}\mathcal{O}_{X'}
  \longrightarrow\mathscr{G}r_n(\mathscr{P}_{X'/S'})
\]
est bijectif pour tout $n$. Supposons alors $f'$ différentiellement lisse et
$\Omega_{X'/S'}^1$ de rang fini. Comme la projection canonique $X'\to X$ est
un morphisme fidèlement plat et quasi-compact, il résulte d'abord de (2.5.2)
que $\Omega_{X/S}^1$ est un $\mathcal{O}_X$-Module localement libre de rang
fini, puis de (2.2.7) que l'homomorphisme canonique (16.3.1.1) est bijectif,
donc $f$ est différentiellement lisse.

\begin{proposition}[16.10.5]
\label{IV.16.10.5}
Pour qu'un morphisme localement de type fini $f:X\to S$ soit
différentiellement lisse, il faut et il suffit que l'immersion diagonale
$\Delta_f:X\to X\times_S X$ soit quasi-régulière.
\end{proposition}

\oldpage[IV-4]{53}

La question étant locale, on peut se borner au cas où $S$ et $X$ sont
affines, et par suite le sous-préschéma diagonal de $X\times_S X$ est fermé.
L'hypothèse que $f$ est localement de type fini entraîne que $\Delta_f$ est
localement de présentation finie (I, 4.3.1), donc le sous-préschéma diagonal
de $X\times_S X$ est défini par un Idéal de type fini $\mathscr{I}$, et
$\Omega_{X/S}^1=\mathscr{I}/\mathscr{I}^2$ est un $\mathcal{O}_X$-Module de
type fini. La proposition résulte alors aussitôt de la comparaison des
conditions dans (16.10.1) et (16.9.4).

\begin{remark}[16.10.6]
\label{IV.16.10.6}
Soit $f:X\to S$ un morphisme tel que le $\mathcal{O}_X$-Module
$\Omega_{X/S}^1$ soit localement libre de rang fini. Il résulte alors de
(0, 20.4.7) que tout $x\in X$ possède un voisinage ouvert $U$ tel qu'il existe
une famille finie $(z_\lambda)_{\lambda\in L}$ de sections de
$\mathcal{O}_X$ au-dessus de $U$ pour lesquelles $(dz_\lambda)_{\lambda\in L}$
forme une \emph{base} du $\Gamma(U,\mathcal{O}_X)$-module
$\Gamma(U,\Omega_{X/S}^1)$.
\end{remark}

\subsection{Opérateurs différentiels sur un $S$-préschéma différentiellement
lisse}
\label{IV.16.11}

\begin{env}[16.11.1]
\label{IV.16.11.1}
Soient $f:X\to S$ un morphisme, $U$ un ouvert de $X$,
$(z_\lambda)_{\lambda\in L}$ une famille de sections de $\mathcal{O}_X$
au-dessus de $U$ telle que les $dz_\lambda$ forment un système de générateurs
de $\Omega_{X/S}^1|U=\Omega_{U/S}^1$. Soit $m$ un entier ou le symbole
$\infty$, et posons, pour tout $\lambda$
\[
  \zeta_\lambda=\delta z_\lambda=d^m z_\lambda-z_\lambda
  \in\Gamma(U,\mathscr{P}_{X/S}^m).
  \tag{16.11.1.1}
  \label{IV.16.11.1.1}
\]

Nous utiliserons d'autre part les notations usuelles de l'analyse; pour tout
$\mathbf{p}=(p_\lambda)\in\mathbb{N}^{(L)}$ (avec $p_\lambda=0$ sauf pour un
nombre fini d'indices), nous poserons
\[
  |\mathbf{p}|=\sum_\lambda p_\lambda,
  \qquad \mathbf{p}!=\prod_\lambda(p_\lambda!)
  \tag{16.11.1.2}
  \label{IV.16.11.1.2}
\]
\[
  \binom{\mathbf{p}}{\mathbf{q}}
  =\mathbf{p}!/(\mathbf{q}!(\mathbf{p}-\mathbf{q})!)
  \quad\text{pour $\mathbf{p},\mathbf{q}$ dans $\mathbb{N}^{(L)}$,
  $\mathbf{q}\leq\mathbf{p}$}
  \tag{16.11.1.3}
  \label{IV.16.11.1.3}
\]
et on convient que $\binom{\mathbf{p}}{\mathbf{q}}=0$ si
$\mathbf{q}\nleq\mathbf{p}$,
\[
  \mathbf{z}^{\mathbf{p}}=\prod_\lambda(z_\lambda)^{p_\lambda},
  \qquad
  \boldsymbol{\zeta}^{\mathbf{p}}=\prod_\lambda(\zeta_\lambda)^{p_\lambda}.
  \tag{16.11.1.4}
  \label{IV.16.11.1.4}
\]

On aura donc, avec ces notations
\[
  d^m(\mathbf{z}^{\mathbf{p}})
  =(d^m(\mathbf{z}))^{\mathbf{p}}
  =(\boldsymbol{\zeta}+\mathbf{z})^{\mathbf{p}}
  =\sum_{\mathbf{q}\leq\mathbf{p}}
    \binom{\mathbf{p}}{\mathbf{q}}
    \mathbf{z}^{\mathbf{p}-\mathbf{q}}\boldsymbol{\zeta}^{\mathbf{q}}
  \tag{16.11.1.5}
  \label{IV.16.11.1.5}
\]
\[
  \boldsymbol{\zeta}^{\mathbf{p}}
  =(d^m\mathbf{z}-\mathbf{z})^{\mathbf{p}}
  =\sum_{\mathbf{q}\leq\mathbf{p}}
    (-1)^{|\mathbf{p}-\mathbf{q}|}
    \binom{\mathbf{p}}{\mathbf{q}}
    \mathbf{z}^{\mathbf{p}-\mathbf{q}}d^m(\mathbf{z}^{\mathbf{q}}).
  \tag{16.11.1.6}
  \label{IV.16.11.1.6}
\]

Comme les $dz_\lambda$ engendrent $\Omega_{X/S}^1$ et sont les images des
$\delta z_\lambda$, et que l'homomorphisme canonique (16.3.1.1) est
surjectif, on en conclut que pour $m$ fini, les $\delta z_\lambda$ engendrent
la $\mathcal{O}_U$-Algèbre $\mathscr{P}_{U/S}^m$ (Bourbaki,
\emph{Alg. comm.}, chap.~III, §~2, n\textsuperscript{o}~8, cor.~2 du th.~1).
Donc les $\boldsymbol{\zeta}^{\mathbf{p}}$ (pour $|\mathbf{p}|\leq m$)
engendrent le $\mathcal{O}_U$-Module $\mathscr{P}_{U/S}^m$. Un opérateur
différentiel $D\in\operatorname{Diff}_{U/S}^m$ est par suite entièrement
déterminé par les valeurs des
$\langle\boldsymbol{\zeta}^{\mathbf{p}},D\rangle$ pour
$|\mathbf{p}|\leq m$, ou, ce qui revient au même par (16.11.1.5) et
(16.11.1.6), par les valeurs
\oldpage[IV-4]{54}
des $\langle d^m(\mathbf{z}^{\mathbf{p}}),D\rangle=D(\mathbf{z}^{\mathbf{p}})$
pour $|\mathbf{p}|\leq m$; de façon précise, il résulte de (16.11.1.5) que
l'on a
\[
  D(\mathbf{z}^{\mathbf{p}})
  =\langle d^m(\mathbf{z}^{\mathbf{p}}),D\rangle
  =\sum_{\mathbf{q}\leq\mathbf{p}}
    \binom{\mathbf{p}}{\mathbf{q}}
    \langle\boldsymbol{\zeta}^{\mathbf{q}},D\rangle
    \mathbf{z}^{\mathbf{p}-\mathbf{q}}.
  \tag{16.11.1.7}
  \label{IV.16.11.1.7}
\]
\end{env}

\begin{theorem}[16.11.2]
\label{IV.16.11.2}
Soient $f:X\to S$ un morphisme, $U$ un ouvert de $X$,
$(z_\lambda)_{\lambda\in L}$ une famille de sections de $\mathcal{O}_X$
au-dessus de $U$ telle que la famille $(dz_\lambda)_{\lambda\in L}$ engendre
$\Omega_{X/S}^1|U=\Omega_{U/S}^1$. Les conditions suivantes sont
équivalentes:
\begin{enumerate}
  \item[a)] $f|U$ est différentiellement lisse et $(dz_\lambda)$ est une base
  du $\mathcal{O}_U$-Module $\Omega_{U/S}^1$.
  \item[b)] Il existe une famille
  $(D_{\mathbf{p}})_{\mathbf{p}\in\mathbb{N}^{(L)}}$ d'opérateurs
  différentiels de $\mathcal{O}_U$ dans lui-même vérifiant les conditions
  \[
    D_{\mathbf{p}}(\mathbf{z}^{\mathbf{q}})
    =\binom{\mathbf{q}}{\mathbf{p}}
      \mathbf{z}^{\mathbf{q}-\mathbf{p}}
    \quad(\mathbf{p},\mathbf{q}\text{ dans }\mathbb{N}^{(L)}).
    \tag{16.11.2.1}
    \label{IV.16.11.2.1}
  \]
\end{enumerate}
De plus, lorsque ces conditions sont vérifiées, la famille
$(D_{\mathbf{p}})$ est déterminée de façon unique par les conditions
(16.11.2.1) et vérifie les relations
\[
  D_{\mathbf{p}}\circ D_{\mathbf{q}}
  =D_{\mathbf{q}}\circ D_{\mathbf{p}}
  =\frac{(\mathbf{p}+\mathbf{q})!}{\mathbf{p}!\mathbf{q}!}
    D_{\mathbf{p}+\mathbf{q}}
  \quad(\mathbf{p},\mathbf{q}\text{ dans }\mathbb{N}^{(L)}).
  \tag{16.11.2.2}
  \label{IV.16.11.2.2}
\]
Enfin, si $L$ est fini, pour tout entier $m$, les $D_{\mathbf{p}}$ tels que
$|\mathbf{p}|\leq m$ forment une base du $\mathcal{O}_U$-Module
$\operatorname{Diff}_{U/S}^m$, autrement dit, tout opérateur différentiel
d'ordre $\leq m$ sur $U$ s'écrit d'une seule façon sous la forme
\[
  D=\sum_{|\mathbf{p}|\leq m}a_{\mathbf{p}}D_{\mathbf{p}}
\]
où les $a_{\mathbf{p}}$ sont des sections de $\mathcal{O}_X$ au-dessus de
$U$.
\end{theorem}

Notons d'abord qu'en vertu de (16.11.1.6) et (16.11.1.5), on vérifie aussitôt
que les conditions (16.11.2.1) sont équivalentes à
\[
  \langle\boldsymbol{\zeta}^{\mathbf{p}},D_{\mathbf{q}}\rangle
  =\delta_{\mathbf{p}\mathbf{q}}
  \quad\text{(indice de Kronecker)}.
  \tag{16.11.2.3}
  \label{IV.16.11.2.3}
\]
L'existence de la famille $(D_{\mathbf{p}})$ vérifiant ces relations entraîne
donc d'abord (en prenant $|\mathbf{p}|=1$) que les $dz_\lambda$ sont
linéairement indépendants, donc forment une \emph{base} du
$\mathcal{O}_U$-Module $\Omega_{U/S}^1$. Puis, pour tout entier $m\geq1$, on
déduit de même de (16.11.2.3) que les $\boldsymbol{\zeta}^{\mathbf{p}}$ tels
que $|\mathbf{p}|\leq m$ sont linéairement indépendants; par suite
l'homomorphisme canonique (16.3.1.1) est injectif, donc bijectif, et cela
prouve que b) entraîne a). La réciproque découle aussitôt de la définition
(16.10.1), le fait que les $\boldsymbol{\zeta}^{\mathbf{p}}$ forment une base
de $\mathscr{P}_{U/S}^m$ pour $|\mathbf{p}|\leq m$ entraînant l'existence et
l'unicité d'une famille d'homomorphismes
$u_{\mathbf{q},m}:\mathscr{P}_{U/S}^m\to\mathcal{O}_U$
($|\mathbf{q}|\leq m$) tels que
$\langle\boldsymbol{\zeta}^{\mathbf{p}},u_{\mathbf{q},m}\rangle
=\delta_{\mathbf{p}\mathbf{q}}$ pour $|\mathbf{p}|\leq m$,
$|\mathbf{q}|\leq m$. Pour une valeur de $\mathbf{q}$ donnée, les opérateurs
différentiels correspondant aux $u_{\mathbf{q},m}$ pour
$m\geq|\mathbf{q}|$ sont identifiés à un même opérateur $D_{\mathbf{q}}$.
Cela prouve donc que a) entraîne b), et en outre que la famille
$(D_{\mathbf{p}})$ est déterminée de façon unique et que, si $L$ est fini,
pour $|\mathbf{p}|\leq m$, les $D_{\mathbf{p}}$ forment une base du dual
$\operatorname{Diff}_{U/S}^m$ de $\mathscr{P}_{U/S}^m$. Enfin, les relations
(16.11.2.2) découlent aussitôt de l'expression des valeurs des trois
opérateurs considérés pour les $\mathbf{z}^{\mathbf{r}}$, et du fait que les
$\boldsymbol{\zeta}^{\mathbf{r}}$ pour $|\mathbf{r}|\leq m$ engendrent
$\mathscr{P}_{U/S}^m$.

\oldpage[IV-4]{55}

\begin{remarks}[16.11.3]
\label{IV.16.11.3}
\begin{enumerate}
  \item[(i)] Le fait que les $D_{\mathbf{p}}$ sont deux à deux permutables en
  vertu de (16.11.2.2) n'implique naturellement pas que la
  $\mathcal{O}_U$-Algèbre $\operatorname{Diff}_{U/S}$ soit commutative, les
  $D_{\mathbf{p}}$ ne permutant aux produits par les sections de
  $\mathcal{O}_U$ que si $n=0$.
  \item[(ii)] Les indices $\mathbf{p}$ tels que $|\mathbf{p}|=1$ sont les
  $\boldsymbol{\varepsilon}_\lambda
  =(\varepsilon_{\lambda\mu})_{\mu\in L}$ avec
  $\varepsilon_{\lambda\mu}=0$ si $\mu\neq\lambda$ et
  $\varepsilon_{\lambda\lambda}=1$; lorsque $L$ est fini, les opérateurs
  $D_{\boldsymbol{\varepsilon}_\lambda}$ ne sont autres que les
  $S$-dérivations $D_i$ introduites dans (16.5.7). On notera qu'en général
  (et contrairement à ce qui se passe en Analyse classique), il n'est pas
  vrai qu'un opérateur différentiel d'ordre quelconque puisse s'écrire comme
  combinaison linéaire de puissances des $D_i$ (cf. (16.12)).
  \item[(iii)] Pour tout entier $r\geq1$, on peut définir la notion de
  morphisme \emph{différentiellement lisse jusqu'à l'ordre $r$} en remplaçant
  dans (16.10.1) la condition (ii) par la condition que les homomorphismes
  \[
    \mathbf{S}_{\mathcal{O}_X}^m(\Omega_{X/S}^1)
    \longrightarrow\mathscr{G}r_m(\mathscr{P}_{X/S})
  \]
  soient bijectifs \emph{pour tout $m\leq r$}. Le raisonnement de (16.11.2)
  prouve alors que si, dans la condition a), on remplace «différentiellement
  lisse» par «différentiellement lisse jusqu'à l'ordre $r$», cette condition
  est équivalente à la condition b) dans laquelle on se borne aux
  $\mathbf{p}\in\mathbb{N}^{(L)}$, $\mathbf{q}\in\mathbb{N}^{(L)}$ tels que
  $|\mathbf{p}|\leq r$, $|\mathbf{q}|\leq r$.
\end{enumerate}
\end{remarks}

\subsection{Cas de la caractéristique nulle: critère jacobien pour les
morphismes différentiellement lisses}
\label{IV.16.12}

\begin{env}[16.12.1]
\label{IV.16.12.1}
On dit qu'un préschéma $X$ est \emph{de caractéristique $p$} ($p$ égal à $0$
ou à un nombre premier) si, pour tout ouvert affine $U$ de $X$, l'anneau
$\Gamma(U,\mathcal{O}_X)$ est de caractéristique $p$ (0, 21.1.1). Il résulte
de (0, 21.1.3) que pour que $X$ soit de caractéristique $0$, il faut et il
suffit que pour tout point \emph{fermé} $x$ de $X$, le corps résiduel $k(x)$
soit de caractéristique $0$, ou encore que $X$ puisse être muni d'une
structure de $\mathbb{Q}$-préschéma (nécessairement unique).
\end{env}

\begin{theorem}[16.12.2]
\label{IV.16.12.2}
Soient $X$ un préschéma de caractéristique $0$, $f:X\to S$ un morphisme. Si
$\Omega_{X/S}^1$ est un $\mathcal{O}_X$-Module localement libre (non
nécessairement de type fini), $f$ est différentiellement lisse.
\end{theorem}

La question étant locale sur $X$, on peut supposer qu'il existe une famille
$(z_\lambda)$ de sections de $\mathcal{O}_X$ au-dessus de $X$ telle que
$(dz_\lambda)$ soit une base du $\mathcal{O}_X$-Module $\Omega_{X/S}^1$.
Appliquant le critère (16.11.2), il suffit de remarquer que les opérateurs
\[
  D_{\mathbf{p}}=(\mathbf{p}!)^{-1}
  \prod_\lambda D_\lambda^{p_\lambda}
\]
(où les $D_\lambda$ sont les formes coordonnées correspondant à la base
$(dz_\lambda)$) vérifient les relations (16.11.2.1), ce qui est une
conséquence du fait que les $D_\lambda$ sont des dérivations.

\begin{env}[16.12.3]
\label{IV.16.12.3}
Le théorème précédent n'est plus exact lorsqu'on abandonne l'hypothèse que
$X$ est de caractéristique $0$. Par exemple, si
$S=\operatorname{Spec}(k)$, où $k$ est un corps de caractéristique $p>0$,
$X=\operatorname{Spec}(K)$ où $K=k(\alpha)$ avec
$\alpha\notin k$, $\alpha^p\in k$, on vérifie aussitôt
\oldpage[IV-4]{56}
que $\Omega_{X/S}^1$ est de rang $1$, et que le morphisme $X\to S$ est
différentiellement lisse jusqu'à l'ordre $p-1$ (16.11.3, (iii)), mais non
jusqu'à l'ordre $p$. Cependant, la démonstration de (16.12.2) prouve que si
$\Omega_{X/S}^1$ est localement libre, et si $n!1_{\mathcal{O}_X}$ est
inversible dans $\Gamma(X,\mathcal{O}_X)$, alors $X$ est différentiellement
lisse sur $S$ jusqu'à l'ordre $n$.
\end{env}
